\documentclass[11pt]{article}

    \usepackage[breakable]{tcolorbox}
    \usepackage{parskip} % Stop auto-indenting (to mimic markdown behaviour)
    

    % Basic figure setup, for now with no caption control since it's done
    % automatically by Pandoc (which extracts ![](path) syntax from Markdown).
    \usepackage{graphicx}
    % Keep aspect ratio if custom image width or height is specified
    \setkeys{Gin}{keepaspectratio}
    % Maintain compatibility with old templates. Remove in nbconvert 6.0
    \let\Oldincludegraphics\includegraphics
    % Ensure that by default, figures have no caption (until we provide a
    % proper Figure object with a Caption API and a way to capture that
    % in the conversion process - todo).
    \usepackage{caption}
    \DeclareCaptionFormat{nocaption}{}
    \captionsetup{format=nocaption,aboveskip=0pt,belowskip=0pt}

    \usepackage{float}
    \floatplacement{figure}{H} % forces figures to be placed at the correct location
    \usepackage{xcolor} % Allow colors to be defined
    \usepackage{enumerate} % Needed for markdown enumerations to work
    \usepackage{geometry} % Used to adjust the document margins
    \usepackage{amsmath} % Equations
    \usepackage{amssymb} % Equations
    \usepackage{textcomp} % defines textquotesingle
    % Hack from http://tex.stackexchange.com/a/47451/13684:
    \AtBeginDocument{%
        \def\PYZsq{\textquotesingle}% Upright quotes in Pygmentized code
    }
    \usepackage{upquote} % Upright quotes for verbatim code
    \usepackage{eurosym} % defines \euro

    \usepackage{iftex}
    \ifPDFTeX
        \usepackage[T1]{fontenc}
        \IfFileExists{alphabeta.sty}{
              \usepackage{alphabeta}
          }{
              \usepackage[mathletters]{ucs}
              \usepackage[utf8x]{inputenc}
          }
    \else
        \usepackage{fontspec}
        \usepackage{unicode-math}
    \fi

    \usepackage{fancyvrb} % verbatim replacement that allows latex
    \usepackage{grffile} % extends the file name processing of package graphics
                         % to support a larger range
    \makeatletter % fix for old versions of grffile with XeLaTeX
    \@ifpackagelater{grffile}{2019/11/01}
    {
      % Do nothing on new versions
    }
    {
      \def\Gread@@xetex#1{%
        \IfFileExists{"\Gin@base".bb}%
        {\Gread@eps{\Gin@base.bb}}%
        {\Gread@@xetex@aux#1}%
      }
    }
    \makeatother
    \usepackage[Export]{adjustbox} % Used to constrain images to a maximum size
    \adjustboxset{max size={0.9\linewidth}{0.9\paperheight}}

    % The hyperref package gives us a pdf with properly built
    % internal navigation ('pdf bookmarks' for the table of contents,
    % internal cross-reference links, web links for URLs, etc.)
    \usepackage{hyperref}
    % The default LaTeX title has an obnoxious amount of whitespace. By default,
    % titling removes some of it. It also provides customization options.
    \usepackage{titling}
    \usepackage{longtable} % longtable support required by pandoc >1.10
    \usepackage{booktabs}  % table support for pandoc > 1.12.2
    \usepackage{array}     % table support for pandoc >= 2.11.3
    \usepackage{calc}      % table minipage width calculation for pandoc >= 2.11.1
    \usepackage[inline]{enumitem} % IRkernel/repr support (it uses the enumerate* environment)
    \usepackage[normalem]{ulem} % ulem is needed to support strikethroughs (\sout)
                                % normalem makes italics be italics, not underlines
    \usepackage{soul}      % strikethrough (\st) support for pandoc >= 3.0.0
    \usepackage{mathrsfs}
    

    
    % Colors for the hyperref package
    \definecolor{urlcolor}{rgb}{0,.145,.698}
    \definecolor{linkcolor}{rgb}{.71,0.21,0.01}
    \definecolor{citecolor}{rgb}{.12,.54,.11}

    % ANSI colors
    \definecolor{ansi-black}{HTML}{3E424D}
    \definecolor{ansi-black-intense}{HTML}{282C36}
    \definecolor{ansi-red}{HTML}{E75C58}
    \definecolor{ansi-red-intense}{HTML}{B22B31}
    \definecolor{ansi-green}{HTML}{00A250}
    \definecolor{ansi-green-intense}{HTML}{007427}
    \definecolor{ansi-yellow}{HTML}{DDB62B}
    \definecolor{ansi-yellow-intense}{HTML}{B27D12}
    \definecolor{ansi-blue}{HTML}{208FFB}
    \definecolor{ansi-blue-intense}{HTML}{0065CA}
    \definecolor{ansi-magenta}{HTML}{D160C4}
    \definecolor{ansi-magenta-intense}{HTML}{A03196}
    \definecolor{ansi-cyan}{HTML}{60C6C8}
    \definecolor{ansi-cyan-intense}{HTML}{258F8F}
    \definecolor{ansi-white}{HTML}{C5C1B4}
    \definecolor{ansi-white-intense}{HTML}{A1A6B2}
    \definecolor{ansi-default-inverse-fg}{HTML}{FFFFFF}
    \definecolor{ansi-default-inverse-bg}{HTML}{000000}

    % common color for the border for error outputs.
    \definecolor{outerrorbackground}{HTML}{FFDFDF}

    % commands and environments needed by pandoc snippets
    % extracted from the output of `pandoc -s`
    \providecommand{\tightlist}{%
      \setlength{\itemsep}{0pt}\setlength{\parskip}{0pt}}
    \DefineVerbatimEnvironment{Highlighting}{Verbatim}{commandchars=\\\{\}}
    % Add ',fontsize=\small' for more characters per line
    \newenvironment{Shaded}{}{}
    \newcommand{\KeywordTok}[1]{\textcolor[rgb]{0.00,0.44,0.13}{\textbf{{#1}}}}
    \newcommand{\DataTypeTok}[1]{\textcolor[rgb]{0.56,0.13,0.00}{{#1}}}
    \newcommand{\DecValTok}[1]{\textcolor[rgb]{0.25,0.63,0.44}{{#1}}}
    \newcommand{\BaseNTok}[1]{\textcolor[rgb]{0.25,0.63,0.44}{{#1}}}
    \newcommand{\FloatTok}[1]{\textcolor[rgb]{0.25,0.63,0.44}{{#1}}}
    \newcommand{\CharTok}[1]{\textcolor[rgb]{0.25,0.44,0.63}{{#1}}}
    \newcommand{\StringTok}[1]{\textcolor[rgb]{0.25,0.44,0.63}{{#1}}}
    \newcommand{\CommentTok}[1]{\textcolor[rgb]{0.38,0.63,0.69}{\textit{{#1}}}}
    \newcommand{\OtherTok}[1]{\textcolor[rgb]{0.00,0.44,0.13}{{#1}}}
    \newcommand{\AlertTok}[1]{\textcolor[rgb]{1.00,0.00,0.00}{\textbf{{#1}}}}
    \newcommand{\FunctionTok}[1]{\textcolor[rgb]{0.02,0.16,0.49}{{#1}}}
    \newcommand{\RegionMarkerTok}[1]{{#1}}
    \newcommand{\ErrorTok}[1]{\textcolor[rgb]{1.00,0.00,0.00}{\textbf{{#1}}}}
    \newcommand{\NormalTok}[1]{{#1}}

    % Additional commands for more recent versions of Pandoc
    \newcommand{\ConstantTok}[1]{\textcolor[rgb]{0.53,0.00,0.00}{{#1}}}
    \newcommand{\SpecialCharTok}[1]{\textcolor[rgb]{0.25,0.44,0.63}{{#1}}}
    \newcommand{\VerbatimStringTok}[1]{\textcolor[rgb]{0.25,0.44,0.63}{{#1}}}
    \newcommand{\SpecialStringTok}[1]{\textcolor[rgb]{0.73,0.40,0.53}{{#1}}}
    \newcommand{\ImportTok}[1]{{#1}}
    \newcommand{\DocumentationTok}[1]{\textcolor[rgb]{0.73,0.13,0.13}{\textit{{#1}}}}
    \newcommand{\AnnotationTok}[1]{\textcolor[rgb]{0.38,0.63,0.69}{\textbf{\textit{{#1}}}}}
    \newcommand{\CommentVarTok}[1]{\textcolor[rgb]{0.38,0.63,0.69}{\textbf{\textit{{#1}}}}}
    \newcommand{\VariableTok}[1]{\textcolor[rgb]{0.10,0.09,0.49}{{#1}}}
    \newcommand{\ControlFlowTok}[1]{\textcolor[rgb]{0.00,0.44,0.13}{\textbf{{#1}}}}
    \newcommand{\OperatorTok}[1]{\textcolor[rgb]{0.40,0.40,0.40}{{#1}}}
    \newcommand{\BuiltInTok}[1]{{#1}}
    \newcommand{\ExtensionTok}[1]{{#1}}
    \newcommand{\PreprocessorTok}[1]{\textcolor[rgb]{0.74,0.48,0.00}{{#1}}}
    \newcommand{\AttributeTok}[1]{\textcolor[rgb]{0.49,0.56,0.16}{{#1}}}
    \newcommand{\InformationTok}[1]{\textcolor[rgb]{0.38,0.63,0.69}{\textbf{\textit{{#1}}}}}
    \newcommand{\WarningTok}[1]{\textcolor[rgb]{0.38,0.63,0.69}{\textbf{\textit{{#1}}}}}
    \makeatletter
    \newsavebox\pandoc@box
    \newcommand*\pandocbounded[1]{%
      \sbox\pandoc@box{#1}%
      % scaling factors for width and height
      \Gscale@div\@tempa\textheight{\dimexpr\ht\pandoc@box+\dp\pandoc@box\relax}%
      \Gscale@div\@tempb\linewidth{\wd\pandoc@box}%
      % select the smaller of both
      \ifdim\@tempb\p@<\@tempa\p@
        \let\@tempa\@tempb
      \fi
      % scaling accordingly (\@tempa < 1)
      \ifdim\@tempa\p@<\p@
        \scalebox{\@tempa}{\usebox\pandoc@box}%
      % scaling not needed, use as it is
      \else
        \usebox{\pandoc@box}%
      \fi
    }
    \makeatother

    % Define a nice break command that doesn't care if a line doesn't already
    % exist.
    \def\br{\hspace*{\fill} \\* }
    % Math Jax compatibility definitions
    \def\gt{>}
    \def\lt{<}
    \let\Oldtex\TeX
    \let\Oldlatex\LaTeX
    \renewcommand{\TeX}{\textrm{\Oldtex}}
    \renewcommand{\LaTeX}{\textrm{\Oldlatex}}
    % Document parameters
    % Document title
    \title{Ch02-statlearn-lab}
    
    
    
    
    
    
    
% Pygments definitions
\makeatletter
\def\PY@reset{\let\PY@it=\relax \let\PY@bf=\relax%
    \let\PY@ul=\relax \let\PY@tc=\relax%
    \let\PY@bc=\relax \let\PY@ff=\relax}
\def\PY@tok#1{\csname PY@tok@#1\endcsname}
\def\PY@toks#1+{\ifx\relax#1\empty\else%
    \PY@tok{#1}\expandafter\PY@toks\fi}
\def\PY@do#1{\PY@bc{\PY@tc{\PY@ul{%
    \PY@it{\PY@bf{\PY@ff{#1}}}}}}}
\def\PY#1#2{\PY@reset\PY@toks#1+\relax+\PY@do{#2}}

\@namedef{PY@tok@w}{\def\PY@tc##1{\textcolor[rgb]{0.73,0.73,0.73}{##1}}}
\@namedef{PY@tok@c}{\let\PY@it=\textit\def\PY@tc##1{\textcolor[rgb]{0.24,0.48,0.48}{##1}}}
\@namedef{PY@tok@cp}{\def\PY@tc##1{\textcolor[rgb]{0.61,0.40,0.00}{##1}}}
\@namedef{PY@tok@k}{\let\PY@bf=\textbf\def\PY@tc##1{\textcolor[rgb]{0.00,0.50,0.00}{##1}}}
\@namedef{PY@tok@kp}{\def\PY@tc##1{\textcolor[rgb]{0.00,0.50,0.00}{##1}}}
\@namedef{PY@tok@kt}{\def\PY@tc##1{\textcolor[rgb]{0.69,0.00,0.25}{##1}}}
\@namedef{PY@tok@o}{\def\PY@tc##1{\textcolor[rgb]{0.40,0.40,0.40}{##1}}}
\@namedef{PY@tok@ow}{\let\PY@bf=\textbf\def\PY@tc##1{\textcolor[rgb]{0.67,0.13,1.00}{##1}}}
\@namedef{PY@tok@nb}{\def\PY@tc##1{\textcolor[rgb]{0.00,0.50,0.00}{##1}}}
\@namedef{PY@tok@nf}{\def\PY@tc##1{\textcolor[rgb]{0.00,0.00,1.00}{##1}}}
\@namedef{PY@tok@nc}{\let\PY@bf=\textbf\def\PY@tc##1{\textcolor[rgb]{0.00,0.00,1.00}{##1}}}
\@namedef{PY@tok@nn}{\let\PY@bf=\textbf\def\PY@tc##1{\textcolor[rgb]{0.00,0.00,1.00}{##1}}}
\@namedef{PY@tok@ne}{\let\PY@bf=\textbf\def\PY@tc##1{\textcolor[rgb]{0.80,0.25,0.22}{##1}}}
\@namedef{PY@tok@nv}{\def\PY@tc##1{\textcolor[rgb]{0.10,0.09,0.49}{##1}}}
\@namedef{PY@tok@no}{\def\PY@tc##1{\textcolor[rgb]{0.53,0.00,0.00}{##1}}}
\@namedef{PY@tok@nl}{\def\PY@tc##1{\textcolor[rgb]{0.46,0.46,0.00}{##1}}}
\@namedef{PY@tok@ni}{\let\PY@bf=\textbf\def\PY@tc##1{\textcolor[rgb]{0.44,0.44,0.44}{##1}}}
\@namedef{PY@tok@na}{\def\PY@tc##1{\textcolor[rgb]{0.41,0.47,0.13}{##1}}}
\@namedef{PY@tok@nt}{\let\PY@bf=\textbf\def\PY@tc##1{\textcolor[rgb]{0.00,0.50,0.00}{##1}}}
\@namedef{PY@tok@nd}{\def\PY@tc##1{\textcolor[rgb]{0.67,0.13,1.00}{##1}}}
\@namedef{PY@tok@s}{\def\PY@tc##1{\textcolor[rgb]{0.73,0.13,0.13}{##1}}}
\@namedef{PY@tok@sd}{\let\PY@it=\textit\def\PY@tc##1{\textcolor[rgb]{0.73,0.13,0.13}{##1}}}
\@namedef{PY@tok@si}{\let\PY@bf=\textbf\def\PY@tc##1{\textcolor[rgb]{0.64,0.35,0.47}{##1}}}
\@namedef{PY@tok@se}{\let\PY@bf=\textbf\def\PY@tc##1{\textcolor[rgb]{0.67,0.36,0.12}{##1}}}
\@namedef{PY@tok@sr}{\def\PY@tc##1{\textcolor[rgb]{0.64,0.35,0.47}{##1}}}
\@namedef{PY@tok@ss}{\def\PY@tc##1{\textcolor[rgb]{0.10,0.09,0.49}{##1}}}
\@namedef{PY@tok@sx}{\def\PY@tc##1{\textcolor[rgb]{0.00,0.50,0.00}{##1}}}
\@namedef{PY@tok@m}{\def\PY@tc##1{\textcolor[rgb]{0.40,0.40,0.40}{##1}}}
\@namedef{PY@tok@gh}{\let\PY@bf=\textbf\def\PY@tc##1{\textcolor[rgb]{0.00,0.00,0.50}{##1}}}
\@namedef{PY@tok@gu}{\let\PY@bf=\textbf\def\PY@tc##1{\textcolor[rgb]{0.50,0.00,0.50}{##1}}}
\@namedef{PY@tok@gd}{\def\PY@tc##1{\textcolor[rgb]{0.63,0.00,0.00}{##1}}}
\@namedef{PY@tok@gi}{\def\PY@tc##1{\textcolor[rgb]{0.00,0.52,0.00}{##1}}}
\@namedef{PY@tok@gr}{\def\PY@tc##1{\textcolor[rgb]{0.89,0.00,0.00}{##1}}}
\@namedef{PY@tok@ge}{\let\PY@it=\textit}
\@namedef{PY@tok@gs}{\let\PY@bf=\textbf}
\@namedef{PY@tok@ges}{\let\PY@bf=\textbf\let\PY@it=\textit}
\@namedef{PY@tok@gp}{\let\PY@bf=\textbf\def\PY@tc##1{\textcolor[rgb]{0.00,0.00,0.50}{##1}}}
\@namedef{PY@tok@go}{\def\PY@tc##1{\textcolor[rgb]{0.44,0.44,0.44}{##1}}}
\@namedef{PY@tok@gt}{\def\PY@tc##1{\textcolor[rgb]{0.00,0.27,0.87}{##1}}}
\@namedef{PY@tok@err}{\def\PY@bc##1{{\setlength{\fboxsep}{\string -\fboxrule}\fcolorbox[rgb]{1.00,0.00,0.00}{1,1,1}{\strut ##1}}}}
\@namedef{PY@tok@kc}{\let\PY@bf=\textbf\def\PY@tc##1{\textcolor[rgb]{0.00,0.50,0.00}{##1}}}
\@namedef{PY@tok@kd}{\let\PY@bf=\textbf\def\PY@tc##1{\textcolor[rgb]{0.00,0.50,0.00}{##1}}}
\@namedef{PY@tok@kn}{\let\PY@bf=\textbf\def\PY@tc##1{\textcolor[rgb]{0.00,0.50,0.00}{##1}}}
\@namedef{PY@tok@kr}{\let\PY@bf=\textbf\def\PY@tc##1{\textcolor[rgb]{0.00,0.50,0.00}{##1}}}
\@namedef{PY@tok@bp}{\def\PY@tc##1{\textcolor[rgb]{0.00,0.50,0.00}{##1}}}
\@namedef{PY@tok@fm}{\def\PY@tc##1{\textcolor[rgb]{0.00,0.00,1.00}{##1}}}
\@namedef{PY@tok@vc}{\def\PY@tc##1{\textcolor[rgb]{0.10,0.09,0.49}{##1}}}
\@namedef{PY@tok@vg}{\def\PY@tc##1{\textcolor[rgb]{0.10,0.09,0.49}{##1}}}
\@namedef{PY@tok@vi}{\def\PY@tc##1{\textcolor[rgb]{0.10,0.09,0.49}{##1}}}
\@namedef{PY@tok@vm}{\def\PY@tc##1{\textcolor[rgb]{0.10,0.09,0.49}{##1}}}
\@namedef{PY@tok@sa}{\def\PY@tc##1{\textcolor[rgb]{0.73,0.13,0.13}{##1}}}
\@namedef{PY@tok@sb}{\def\PY@tc##1{\textcolor[rgb]{0.73,0.13,0.13}{##1}}}
\@namedef{PY@tok@sc}{\def\PY@tc##1{\textcolor[rgb]{0.73,0.13,0.13}{##1}}}
\@namedef{PY@tok@dl}{\def\PY@tc##1{\textcolor[rgb]{0.73,0.13,0.13}{##1}}}
\@namedef{PY@tok@s2}{\def\PY@tc##1{\textcolor[rgb]{0.73,0.13,0.13}{##1}}}
\@namedef{PY@tok@sh}{\def\PY@tc##1{\textcolor[rgb]{0.73,0.13,0.13}{##1}}}
\@namedef{PY@tok@s1}{\def\PY@tc##1{\textcolor[rgb]{0.73,0.13,0.13}{##1}}}
\@namedef{PY@tok@mb}{\def\PY@tc##1{\textcolor[rgb]{0.40,0.40,0.40}{##1}}}
\@namedef{PY@tok@mf}{\def\PY@tc##1{\textcolor[rgb]{0.40,0.40,0.40}{##1}}}
\@namedef{PY@tok@mh}{\def\PY@tc##1{\textcolor[rgb]{0.40,0.40,0.40}{##1}}}
\@namedef{PY@tok@mi}{\def\PY@tc##1{\textcolor[rgb]{0.40,0.40,0.40}{##1}}}
\@namedef{PY@tok@il}{\def\PY@tc##1{\textcolor[rgb]{0.40,0.40,0.40}{##1}}}
\@namedef{PY@tok@mo}{\def\PY@tc##1{\textcolor[rgb]{0.40,0.40,0.40}{##1}}}
\@namedef{PY@tok@ch}{\let\PY@it=\textit\def\PY@tc##1{\textcolor[rgb]{0.24,0.48,0.48}{##1}}}
\@namedef{PY@tok@cm}{\let\PY@it=\textit\def\PY@tc##1{\textcolor[rgb]{0.24,0.48,0.48}{##1}}}
\@namedef{PY@tok@cpf}{\let\PY@it=\textit\def\PY@tc##1{\textcolor[rgb]{0.24,0.48,0.48}{##1}}}
\@namedef{PY@tok@c1}{\let\PY@it=\textit\def\PY@tc##1{\textcolor[rgb]{0.24,0.48,0.48}{##1}}}
\@namedef{PY@tok@cs}{\let\PY@it=\textit\def\PY@tc##1{\textcolor[rgb]{0.24,0.48,0.48}{##1}}}

\def\PYZbs{\char`\\}
\def\PYZus{\char`\_}
\def\PYZob{\char`\{}
\def\PYZcb{\char`\}}
\def\PYZca{\char`\^}
\def\PYZam{\char`\&}
\def\PYZlt{\char`\<}
\def\PYZgt{\char`\>}
\def\PYZsh{\char`\#}
\def\PYZpc{\char`\%}
\def\PYZdl{\char`\$}
\def\PYZhy{\char`\-}
\def\PYZsq{\char`\'}
\def\PYZdq{\char`\"}
\def\PYZti{\char`\~}
% for compatibility with earlier versions
\def\PYZat{@}
\def\PYZlb{[}
\def\PYZrb{]}
\makeatother


    % For linebreaks inside Verbatim environment from package fancyvrb.
    \makeatletter
        \newbox\Wrappedcontinuationbox
        \newbox\Wrappedvisiblespacebox
        \newcommand*\Wrappedvisiblespace {\textcolor{red}{\textvisiblespace}}
        \newcommand*\Wrappedcontinuationsymbol {\textcolor{red}{\llap{\tiny$\m@th\hookrightarrow$}}}
        \newcommand*\Wrappedcontinuationindent {3ex }
        \newcommand*\Wrappedafterbreak {\kern\Wrappedcontinuationindent\copy\Wrappedcontinuationbox}
        % Take advantage of the already applied Pygments mark-up to insert
        % potential linebreaks for TeX processing.
        %        {, <, #, %, $, ' and ": go to next line.
        %        _, }, ^, &, >, - and ~: stay at end of broken line.
        % Use of \textquotesingle for straight quote.
        \newcommand*\Wrappedbreaksatspecials {%
            \def\PYGZus{\discretionary{\char`\_}{\Wrappedafterbreak}{\char`\_}}%
            \def\PYGZob{\discretionary{}{\Wrappedafterbreak\char`\{}{\char`\{}}%
            \def\PYGZcb{\discretionary{\char`\}}{\Wrappedafterbreak}{\char`\}}}%
            \def\PYGZca{\discretionary{\char`\^}{\Wrappedafterbreak}{\char`\^}}%
            \def\PYGZam{\discretionary{\char`\&}{\Wrappedafterbreak}{\char`\&}}%
            \def\PYGZlt{\discretionary{}{\Wrappedafterbreak\char`\<}{\char`\<}}%
            \def\PYGZgt{\discretionary{\char`\>}{\Wrappedafterbreak}{\char`\>}}%
            \def\PYGZsh{\discretionary{}{\Wrappedafterbreak\char`\#}{\char`\#}}%
            \def\PYGZpc{\discretionary{}{\Wrappedafterbreak\char`\%}{\char`\%}}%
            \def\PYGZdl{\discretionary{}{\Wrappedafterbreak\char`\$}{\char`\$}}%
            \def\PYGZhy{\discretionary{\char`\-}{\Wrappedafterbreak}{\char`\-}}%
            \def\PYGZsq{\discretionary{}{\Wrappedafterbreak\textquotesingle}{\textquotesingle}}%
            \def\PYGZdq{\discretionary{}{\Wrappedafterbreak\char`\"}{\char`\"}}%
            \def\PYGZti{\discretionary{\char`\~}{\Wrappedafterbreak}{\char`\~}}%
        }
        % Some characters . , ; ? ! / are not pygmentized.
        % This macro makes them "active" and they will insert potential linebreaks
        \newcommand*\Wrappedbreaksatpunct {%
            \lccode`\~`\.\lowercase{\def~}{\discretionary{\hbox{\char`\.}}{\Wrappedafterbreak}{\hbox{\char`\.}}}%
            \lccode`\~`\,\lowercase{\def~}{\discretionary{\hbox{\char`\,}}{\Wrappedafterbreak}{\hbox{\char`\,}}}%
            \lccode`\~`\;\lowercase{\def~}{\discretionary{\hbox{\char`\;}}{\Wrappedafterbreak}{\hbox{\char`\;}}}%
            \lccode`\~`\:\lowercase{\def~}{\discretionary{\hbox{\char`\:}}{\Wrappedafterbreak}{\hbox{\char`\:}}}%
            \lccode`\~`\?\lowercase{\def~}{\discretionary{\hbox{\char`\?}}{\Wrappedafterbreak}{\hbox{\char`\?}}}%
            \lccode`\~`\!\lowercase{\def~}{\discretionary{\hbox{\char`\!}}{\Wrappedafterbreak}{\hbox{\char`\!}}}%
            \lccode`\~`\/\lowercase{\def~}{\discretionary{\hbox{\char`\/}}{\Wrappedafterbreak}{\hbox{\char`\/}}}%
            \catcode`\.\active
            \catcode`\,\active
            \catcode`\;\active
            \catcode`\:\active
            \catcode`\?\active
            \catcode`\!\active
            \catcode`\/\active
            \lccode`\~`\~
        }
    \makeatother

    \let\OriginalVerbatim=\Verbatim
    \makeatletter
    \renewcommand{\Verbatim}[1][1]{%
        %\parskip\z@skip
        \sbox\Wrappedcontinuationbox {\Wrappedcontinuationsymbol}%
        \sbox\Wrappedvisiblespacebox {\FV@SetupFont\Wrappedvisiblespace}%
        \def\FancyVerbFormatLine ##1{\hsize\linewidth
            \vtop{\raggedright\hyphenpenalty\z@\exhyphenpenalty\z@
                \doublehyphendemerits\z@\finalhyphendemerits\z@
                \strut ##1\strut}%
        }%
        % If the linebreak is at a space, the latter will be displayed as visible
        % space at end of first line, and a continuation symbol starts next line.
        % Stretch/shrink are however usually zero for typewriter font.
        \def\FV@Space {%
            \nobreak\hskip\z@ plus\fontdimen3\font minus\fontdimen4\font
            \discretionary{\copy\Wrappedvisiblespacebox}{\Wrappedafterbreak}
            {\kern\fontdimen2\font}%
        }%

        % Allow breaks at special characters using \PYG... macros.
        \Wrappedbreaksatspecials
        % Breaks at punctuation characters . , ; ? ! and / need catcode=\active
        \OriginalVerbatim[#1,codes*=\Wrappedbreaksatpunct]%
    }
    \makeatother

    % Exact colors from NB
    \definecolor{incolor}{HTML}{303F9F}
    \definecolor{outcolor}{HTML}{D84315}
    \definecolor{cellborder}{HTML}{CFCFCF}
    \definecolor{cellbackground}{HTML}{F7F7F7}

    % prompt
    \makeatletter
    \newcommand{\boxspacing}{\kern\kvtcb@left@rule\kern\kvtcb@boxsep}
    \makeatother
    \newcommand{\prompt}[4]{
        {\ttfamily\llap{{\color{#2}[#3]:\hspace{3pt}#4}}\vspace{-\baselineskip}}
    }
    

    
    % Prevent overflowing lines due to hard-to-break entities
    \sloppy
    % Setup hyperref package
    \hypersetup{
      breaklinks=true,  % so long urls are correctly broken across lines
      colorlinks=true,
      urlcolor=urlcolor,
      linkcolor=linkcolor,
      citecolor=citecolor,
      }
    % Slightly bigger margins than the latex defaults
    
    \geometry{verbose,tmargin=1in,bmargin=1in,lmargin=1in,rmargin=1in}
    
    

\begin{document}
    
    \maketitle
    
    

    
    \hypertarget{introduction-to-python}{%
\section{Introduction to Python}\label{introduction-to-python}}

\href{https://mybinder.org/v2/gh/intro-stat-learning/ISLP_labs/v2.2?labpath=Ch02-statlearn-lab.ipynb}{\includesvg{https://mybinder.org/badge_logo.svg}}

    \hypertarget{getting-started}{%
\subsection{Getting Started}\label{getting-started}}

    To run the labs in this book, you will need two things:

\begin{itemize}
\tightlist
\item
  An installation of \texttt{Python3}, which is the specific version of
  \texttt{Python} used in the labs.
\item
  Access to \texttt{Jupyter}, a very popular \texttt{Python} interface
  that runs code through a file called a \emph{notebook}.
\end{itemize}

    You can download and install \texttt{Python3} by following the
instructions available at \href{http://anaconda.com}{anaconda.com}.

    There are a number of ways to get access to \texttt{Jupyter}. Here are
just a few:

\begin{itemize}
\tightlist
\item
  Using Google's \texttt{Colaboratory} service:
  \href{https://colab.research.google.com/}{colab.research.google.com/}.
\item
  Using \texttt{JupyterHub}, available at
  \href{https://jupyter.org/hub}{jupyter.org/hub}.
\item
  Using your own \texttt{jupyter} installation. Installation
  instructions are available at
  \href{https://jupyter.org/install}{jupyter.org/install}.
\end{itemize}

Please see the \texttt{Python} resources page on the book website
\href{https://www.statlearning.com}{statlearning.com} for up-to-date
information about getting \texttt{Python} and \texttt{Jupyter} working
on your computer.

You will need to install the \texttt{ISLP} package, which provides
access to the datasets and custom-built functions that we provide.
Inside a macOS or Linux terminal type \texttt{pip\ install\ ISLP}; this
also installs most other packages needed in the labs. The
\texttt{Python} resources page has a link to the \texttt{ISLP}
documentation website.

To run this lab, download the file \texttt{Ch2-statlearn-lab.ipynb} from
the \texttt{Python} resources page. Now run the following code at the
command line: \texttt{jupyter\ lab\ Ch2-statlearn-lab.ipynb}.

If you're using Windows, you can use the \texttt{start\ menu} to access
\texttt{anaconda}, and follow the links. For example, to install
\texttt{ISLP} and run this lab, you can run the same code above in an
\texttt{anaconda} shell.

    \hypertarget{basic-commands}{%
\subsection{Basic Commands}\label{basic-commands}}

    In this lab, we will introduce some simple \texttt{Python} commands. For
more resources about \texttt{Python} in general, readers may want to
consult the tutorial at
\href{https://docs.python.org/3/tutorial/}{docs.python.org/3/tutorial/}.

    Like most programming languages, \texttt{Python} uses \emph{functions}
to perform operations. To run a function called \texttt{fun}, we type
\texttt{fun(input1,input2)}, where the inputs (or \emph{arguments})
\texttt{input1} and \texttt{input2} tell \texttt{Python} how to run the
function. A function can have any number of inputs. For example, the
\texttt{print()} function outputs a text representation of all of its
arguments to the console.

    \begin{tcolorbox}[breakable, size=fbox, boxrule=1pt, pad at break*=1mm,colback=cellbackground, colframe=cellborder]
\prompt{In}{incolor}{1}{\boxspacing}
\begin{Verbatim}[commandchars=\\\{\}]
\PY{n+nb}{print}\PY{p}{(}\PY{l+s+s1}{\PYZsq{}}\PY{l+s+s1}{fit a model with}\PY{l+s+s1}{\PYZsq{}}\PY{p}{,} \PY{l+m+mi}{11}\PY{p}{,} \PY{l+s+s1}{\PYZsq{}}\PY{l+s+s1}{variables}\PY{l+s+s1}{\PYZsq{}}\PY{p}{)}
\end{Verbatim}
\end{tcolorbox}

    \begin{Verbatim}[commandchars=\\\{\}]
fit a model with 11 variables
    \end{Verbatim}

    The following command will provide information about the
\texttt{print()} function.

    \begin{tcolorbox}[breakable, size=fbox, boxrule=1pt, pad at break*=1mm,colback=cellbackground, colframe=cellborder]
\prompt{In}{incolor}{2}{\boxspacing}
\begin{Verbatim}[commandchars=\\\{\}]
print\PY{o}{?}
\end{Verbatim}
\end{tcolorbox}

    
    \begin{Verbatim}[commandchars=\\\{\}]
\textcolor{ansi-red}{Signature:} print(*args, sep=\textcolor{ansi-yellow}{' '}, end=\textcolor{ansi-yellow}{'\textbackslash{}n'}, file=\def\tcRGB{\textcolor[RGB]}\expandafter\tcRGB\expandafter{\detokenize{0,135,0}}{\textbf{None}}, flush=\def\tcRGB{\textcolor[RGB]}\expandafter\tcRGB\expandafter{\detokenize{0,135,0}}{\textbf{False}})
\textcolor{ansi-red}{Docstring:}
Prints the values to a stream, or to sys.stdout by default.

sep
  string inserted between values, default a space.
end
  string appended after the last value, default a newline.
file
  a file-like object (stream); defaults to the current sys.stdout.
flush
  whether to forcibly flush the stream.
\textcolor{ansi-red}{Type:}      builtin\_function\_or\_method
    \end{Verbatim}

    
    Adding two integers in \texttt{Python} is pretty intuitive.

    \begin{tcolorbox}[breakable, size=fbox, boxrule=1pt, pad at break*=1mm,colback=cellbackground, colframe=cellborder]
\prompt{In}{incolor}{3}{\boxspacing}
\begin{Verbatim}[commandchars=\\\{\}]
\PY{l+m+mi}{3} \PY{o}{+} \PY{l+m+mi}{5}
\end{Verbatim}
\end{tcolorbox}

            \begin{tcolorbox}[breakable, size=fbox, boxrule=.5pt, pad at break*=1mm, opacityfill=0]
\prompt{Out}{outcolor}{3}{\boxspacing}
\begin{Verbatim}[commandchars=\\\{\}]
8
\end{Verbatim}
\end{tcolorbox}
        
    In \texttt{Python}, textual data is handled using \emph{strings}. For
instance, \texttt{"hello"} and
\texttt{\textquotesingle{}hello\textquotesingle{}} are strings. We can
concatenate them using the addition \texttt{+} symbol.

    \begin{tcolorbox}[breakable, size=fbox, boxrule=1pt, pad at break*=1mm,colback=cellbackground, colframe=cellborder]
\prompt{In}{incolor}{4}{\boxspacing}
\begin{Verbatim}[commandchars=\\\{\}]
\PY{l+s+s2}{\PYZdq{}}\PY{l+s+s2}{hello}\PY{l+s+s2}{\PYZdq{}} \PY{o}{+} \PY{l+s+s2}{\PYZdq{}}\PY{l+s+s2}{ }\PY{l+s+s2}{\PYZdq{}} \PY{o}{+} \PY{l+s+s2}{\PYZdq{}}\PY{l+s+s2}{world}\PY{l+s+s2}{\PYZdq{}}
\end{Verbatim}
\end{tcolorbox}

            \begin{tcolorbox}[breakable, size=fbox, boxrule=.5pt, pad at break*=1mm, opacityfill=0]
\prompt{Out}{outcolor}{4}{\boxspacing}
\begin{Verbatim}[commandchars=\\\{\}]
'hello world'
\end{Verbatim}
\end{tcolorbox}
        
    A string is actually a type of \emph{sequence}: this is a generic term
for an ordered list. The three most important types of sequences are
lists, tuples, and strings.\\
We introduce lists now.

    The following command instructs \texttt{Python} to join together the
numbers 3, 4, and 5, and to save them as a \emph{list} named \texttt{x}.
When we type \texttt{x}, it gives us back the list.

    \begin{tcolorbox}[breakable, size=fbox, boxrule=1pt, pad at break*=1mm,colback=cellbackground, colframe=cellborder]
\prompt{In}{incolor}{5}{\boxspacing}
\begin{Verbatim}[commandchars=\\\{\}]
\PY{n}{x} \PY{o}{=} \PY{p}{[}\PY{l+m+mi}{3}\PY{p}{,} \PY{l+m+mi}{4}\PY{p}{,} \PY{l+m+mi}{5}\PY{p}{]}
\PY{n}{x}
\end{Verbatim}
\end{tcolorbox}

            \begin{tcolorbox}[breakable, size=fbox, boxrule=.5pt, pad at break*=1mm, opacityfill=0]
\prompt{Out}{outcolor}{5}{\boxspacing}
\begin{Verbatim}[commandchars=\\\{\}]
[3, 4, 5]
\end{Verbatim}
\end{tcolorbox}
        
    Note that we used the brackets \texttt{{[}{]}} to construct this list.

We will often want to add two sets of numbers together. It is reasonable
to try the following code, though it will not produce the desired
results.

    \begin{tcolorbox}[breakable, size=fbox, boxrule=1pt, pad at break*=1mm,colback=cellbackground, colframe=cellborder]
\prompt{In}{incolor}{6}{\boxspacing}
\begin{Verbatim}[commandchars=\\\{\}]
\PY{n}{y} \PY{o}{=} \PY{p}{[}\PY{l+m+mi}{4}\PY{p}{,} \PY{l+m+mi}{9}\PY{p}{,} \PY{l+m+mi}{7}\PY{p}{]}
\PY{n}{x} \PY{o}{+} \PY{n}{y}
\end{Verbatim}
\end{tcolorbox}

            \begin{tcolorbox}[breakable, size=fbox, boxrule=.5pt, pad at break*=1mm, opacityfill=0]
\prompt{Out}{outcolor}{6}{\boxspacing}
\begin{Verbatim}[commandchars=\\\{\}]
[3, 4, 5, 4, 9, 7]
\end{Verbatim}
\end{tcolorbox}
        
    The result may appear slightly counterintuitive: why did \texttt{Python}
not add the entries of the lists element-by-element? In \texttt{Python},
lists hold \emph{arbitrary} objects, and are added using
\emph{concatenation}. In fact, concatenation is the behavior that we saw
earlier when we entered \texttt{"hello"\ +\ "\ "\ +\ "world"}.

    This example reflects the fact that \texttt{Python} is a general-purpose
programming language. Much of \texttt{Python}'s data-specific
functionality comes from other packages, notably \texttt{numpy} and
\texttt{pandas}. In the next section, we will introduce the
\texttt{numpy} package. See
\href{https://docs.scipy.org/doc/numpy/user/quickstart.html}{docs.scipy.org/doc/numpy/user/quickstart.html}
for more information about \texttt{numpy}.

    \hypertarget{introduction-to-numerical-python}{%
\subsection{Introduction to Numerical
Python}\label{introduction-to-numerical-python}}

As mentioned earlier, this book makes use of functionality that is
contained in the \texttt{numpy} \emph{library}, or \emph{package}. A
package is a collection of modules that are not necessarily included in
the base \texttt{Python} distribution. The name \texttt{numpy} is an
abbreviation for \emph{numerical Python}.

    To access \texttt{numpy}, we must first \texttt{import} it.

    \begin{tcolorbox}[breakable, size=fbox, boxrule=1pt, pad at break*=1mm,colback=cellbackground, colframe=cellborder]
\prompt{In}{incolor}{7}{\boxspacing}
\begin{Verbatim}[commandchars=\\\{\}]
\PY{k+kn}{import}\PY{+w}{ }\PY{n+nn}{numpy}\PY{+w}{ }\PY{k}{as}\PY{+w}{ }\PY{n+nn}{np} 
\end{Verbatim}
\end{tcolorbox}

    In the previous line, we named the \texttt{numpy} \emph{module}
\texttt{np}; an abbreviation for easier referencing.

    In \texttt{numpy}, an \emph{array} is a generic term for a
multidimensional set of numbers. We use the \texttt{np.array()} function
to define \texttt{x} and \texttt{y}, which are one-dimensional arrays,
i.e.~vectors.

    \begin{tcolorbox}[breakable, size=fbox, boxrule=1pt, pad at break*=1mm,colback=cellbackground, colframe=cellborder]
\prompt{In}{incolor}{8}{\boxspacing}
\begin{Verbatim}[commandchars=\\\{\}]
\PY{n}{x} \PY{o}{=} \PY{n}{np}\PY{o}{.}\PY{n}{array}\PY{p}{(}\PY{p}{[}\PY{l+m+mi}{3}\PY{p}{,} \PY{l+m+mi}{4}\PY{p}{,} \PY{l+m+mi}{5}\PY{p}{]}\PY{p}{)}
\PY{n}{y} \PY{o}{=} \PY{n}{np}\PY{o}{.}\PY{n}{array}\PY{p}{(}\PY{p}{[}\PY{l+m+mi}{4}\PY{p}{,} \PY{l+m+mi}{9}\PY{p}{,} \PY{l+m+mi}{7}\PY{p}{]}\PY{p}{)}
\end{Verbatim}
\end{tcolorbox}

    Note that if you forgot to run the \texttt{import\ numpy\ as\ np}
command earlier, then you will encounter an error in calling the
\texttt{np.array()} function in the previous line. The syntax
\texttt{np.array()} indicates that the function being called is part of
the \texttt{numpy} package, which we have abbreviated as \texttt{np}.

    Since \texttt{x} and \texttt{y} have been defined using
\texttt{np.array()}, we get a sensible result when we add them together.
Compare this to our results in the previous section, when we tried to
add two lists without using \texttt{numpy}.

    \begin{tcolorbox}[breakable, size=fbox, boxrule=1pt, pad at break*=1mm,colback=cellbackground, colframe=cellborder]
\prompt{In}{incolor}{9}{\boxspacing}
\begin{Verbatim}[commandchars=\\\{\}]
\PY{n}{x} \PY{o}{+} \PY{n}{y}
\end{Verbatim}
\end{tcolorbox}

            \begin{tcolorbox}[breakable, size=fbox, boxrule=.5pt, pad at break*=1mm, opacityfill=0]
\prompt{Out}{outcolor}{9}{\boxspacing}
\begin{Verbatim}[commandchars=\\\{\}]
array([ 7, 13, 12])
\end{Verbatim}
\end{tcolorbox}
        
    

    In \texttt{numpy}, matrices are typically represented as two-dimensional
arrays, and vectors as one-dimensional arrays. \{While it is also
possible to create matrices using \texttt{np.matrix()}, we will use
\texttt{np.array()} throughout the labs in this book.\} We can create a
two-dimensional array as follows.

    \begin{tcolorbox}[breakable, size=fbox, boxrule=1pt, pad at break*=1mm,colback=cellbackground, colframe=cellborder]
\prompt{In}{incolor}{10}{\boxspacing}
\begin{Verbatim}[commandchars=\\\{\}]
\PY{n}{x} \PY{o}{=} \PY{n}{np}\PY{o}{.}\PY{n}{array}\PY{p}{(}\PY{p}{[}\PY{p}{[}\PY{l+m+mi}{1}\PY{p}{,} \PY{l+m+mi}{2}\PY{p}{]}\PY{p}{,} \PY{p}{[}\PY{l+m+mi}{3}\PY{p}{,} \PY{l+m+mi}{4}\PY{p}{]}\PY{p}{]}\PY{p}{)}
\PY{n}{x}
\end{Verbatim}
\end{tcolorbox}

            \begin{tcolorbox}[breakable, size=fbox, boxrule=.5pt, pad at break*=1mm, opacityfill=0]
\prompt{Out}{outcolor}{10}{\boxspacing}
\begin{Verbatim}[commandchars=\\\{\}]
array([[1, 2],
       [3, 4]])
\end{Verbatim}
\end{tcolorbox}
        
    

    The object \texttt{x} has several \emph{attributes}, or associated
objects. To access an attribute of \texttt{x}, we type
\texttt{x.attribute}, where we replace \texttt{attribute} with the name
of the attribute. For instance, we can access the \texttt{ndim}
attribute of \texttt{x} as follows.

    \begin{tcolorbox}[breakable, size=fbox, boxrule=1pt, pad at break*=1mm,colback=cellbackground, colframe=cellborder]
\prompt{In}{incolor}{11}{\boxspacing}
\begin{Verbatim}[commandchars=\\\{\}]
\PY{n}{x}\PY{o}{.}\PY{n}{ndim}
\end{Verbatim}
\end{tcolorbox}

            \begin{tcolorbox}[breakable, size=fbox, boxrule=.5pt, pad at break*=1mm, opacityfill=0]
\prompt{Out}{outcolor}{11}{\boxspacing}
\begin{Verbatim}[commandchars=\\\{\}]
2
\end{Verbatim}
\end{tcolorbox}
        
    The output indicates that \texttt{x} is a two-dimensional array.\\
Similarly, \texttt{x.dtype} is the \emph{data type} attribute of the
object \texttt{x}. This indicates that \texttt{x} is comprised of 64-bit
integers:

    \begin{tcolorbox}[breakable, size=fbox, boxrule=1pt, pad at break*=1mm,colback=cellbackground, colframe=cellborder]
\prompt{In}{incolor}{12}{\boxspacing}
\begin{Verbatim}[commandchars=\\\{\}]
\PY{n}{x}\PY{o}{.}\PY{n}{dtype}
\end{Verbatim}
\end{tcolorbox}

            \begin{tcolorbox}[breakable, size=fbox, boxrule=.5pt, pad at break*=1mm, opacityfill=0]
\prompt{Out}{outcolor}{12}{\boxspacing}
\begin{Verbatim}[commandchars=\\\{\}]
dtype('int64')
\end{Verbatim}
\end{tcolorbox}
        
    Why is \texttt{x} comprised of integers? This is because we created
\texttt{x} by passing in exclusively integers to the \texttt{np.array()}
function. If we had passed in any decimals, then we would have obtained
an array of \emph{floating point numbers} (i.e.~real-valued numbers).

    \begin{tcolorbox}[breakable, size=fbox, boxrule=1pt, pad at break*=1mm,colback=cellbackground, colframe=cellborder]
\prompt{In}{incolor}{13}{\boxspacing}
\begin{Verbatim}[commandchars=\\\{\}]
\PY{n}{np}\PY{o}{.}\PY{n}{array}\PY{p}{(}\PY{p}{[}\PY{p}{[}\PY{l+m+mi}{1}\PY{p}{,} \PY{l+m+mi}{2}\PY{p}{]}\PY{p}{,} \PY{p}{[}\PY{l+m+mf}{3.0}\PY{p}{,} \PY{l+m+mi}{4}\PY{p}{]}\PY{p}{]}\PY{p}{)}\PY{o}{.}\PY{n}{dtype}
\end{Verbatim}
\end{tcolorbox}

            \begin{tcolorbox}[breakable, size=fbox, boxrule=.5pt, pad at break*=1mm, opacityfill=0]
\prompt{Out}{outcolor}{13}{\boxspacing}
\begin{Verbatim}[commandchars=\\\{\}]
dtype('float64')
\end{Verbatim}
\end{tcolorbox}
        
    Typing \texttt{fun?} will cause \texttt{Python} to display documentation
associated with the function \texttt{fun}, if it exists. We can try this
for \texttt{np.array()}.

    \begin{tcolorbox}[breakable, size=fbox, boxrule=1pt, pad at break*=1mm,colback=cellbackground, colframe=cellborder]
\prompt{In}{incolor}{14}{\boxspacing}
\begin{Verbatim}[commandchars=\\\{\}]
np.array\PY{o}{?}
\end{Verbatim}
\end{tcolorbox}

    
    \begin{Verbatim}[commandchars=\\\{\}]
\textcolor{ansi-red}{Docstring:}
array(object, dtype=None, *, copy=True, order='K', subok=False, ndmin=0,
      like=None)

Create an array.

Parameters
----------
object : array\_like
    An array, any object exposing the array interface, an object whose
    ``\_\_array\_\_`` method returns an array, or any (nested) sequence.
    If object is a scalar, a 0-dimensional array containing object is
    returned.
dtype : data-type, optional
    The desired data-type for the array. If not given, NumPy will try to use
    a default ``dtype`` that can represent the values (by applying promotion
    rules when necessary.)
copy : bool, optional
    If true (default), then the object is copied.  Otherwise, a copy will
    only be made if ``\_\_array\_\_`` returns a copy, if obj is a nested
    sequence, or if a copy is needed to satisfy any of the other
    requirements (``dtype``, ``order``, etc.).
order : \{'K', 'A', 'C', 'F'\}, optional
    Specify the memory layout of the array. If object is not an array, the
    newly created array will be in C order (row major) unless 'F' is
    specified, in which case it will be in Fortran order (column major).
    If object is an array the following holds.

    ===== ========= ===================================================
    order  no copy                     copy=True
    ===== ========= ===================================================
    'K'   unchanged F \& C order preserved, otherwise most similar order
    'A'   unchanged F order if input is F and not C, otherwise C order
    'C'   C order   C order
    'F'   F order   F order
    ===== ========= ===================================================

    When ``copy=False`` and a copy is made for other reasons, the result is
    the same as if ``copy=True``, with some exceptions for 'A', see the
    Notes section. The default order is 'K'.
subok : bool, optional
    If True, then sub-classes will be passed-through, otherwise
    the returned array will be forced to be a base-class array (default).
ndmin : int, optional
    Specifies the minimum number of dimensions that the resulting
    array should have.  Ones will be prepended to the shape as
    needed to meet this requirement.
like : array\_like, optional
    Reference object to allow the creation of arrays which are not
    NumPy arrays. If an array-like passed in as ``like`` supports
    the ``\_\_array\_function\_\_`` protocol, the result will be defined
    by it. In this case, it ensures the creation of an array object
    compatible with that passed in via this argument.

    .. versionadded:: 1.20.0

Returns
-------
out : ndarray
    An array object satisfying the specified requirements.

See Also
--------
empty\_like : Return an empty array with shape and type of input.
ones\_like : Return an array of ones with shape and type of input.
zeros\_like : Return an array of zeros with shape and type of input.
full\_like : Return a new array with shape of input filled with value.
empty : Return a new uninitialized array.
ones : Return a new array setting values to one.
zeros : Return a new array setting values to zero.
full : Return a new array of given shape filled with value.


Notes
-----
When order is 'A' and ``object`` is an array in neither 'C' nor 'F' order,
and a copy is forced by a change in dtype, then the order of the result is
not necessarily 'C' as expected. This is likely a bug.

Examples
--------
>>> np.array([1, 2, 3])
array([1, 2, 3])

Upcasting:

>>> np.array([1, 2, 3.0])
array([ 1.,  2.,  3.])

More than one dimension:

>>> np.array([[1, 2], [3, 4]])
array([[1, 2],
       [3, 4]])

Minimum dimensions 2:

>>> np.array([1, 2, 3], ndmin=2)
array([[1, 2, 3]])

Type provided:

>>> np.array([1, 2, 3], dtype=complex)
array([ 1.+0.j,  2.+0.j,  3.+0.j])

Data-type consisting of more than one element:

>>> x = np.array([(1,2),(3,4)],dtype=[('a','<i4'),('b','<i4')])
>>> x['a']
array([1, 3])

Creating an array from sub-classes:

>>> np.array(np.mat('1 2; 3 4'))
array([[1, 2],
       [3, 4]])

>>> np.array(np.mat('1 2; 3 4'), subok=True)
matrix([[1, 2],
        [3, 4]])
\textcolor{ansi-red}{Type:}      builtin\_function\_or\_method
    \end{Verbatim}

    
    This documentation indicates that we could create a floating point array
by passing a \texttt{dtype} argument into \texttt{np.array()}.

    \begin{tcolorbox}[breakable, size=fbox, boxrule=1pt, pad at break*=1mm,colback=cellbackground, colframe=cellborder]
\prompt{In}{incolor}{15}{\boxspacing}
\begin{Verbatim}[commandchars=\\\{\}]
\PY{n}{np}\PY{o}{.}\PY{n}{array}\PY{p}{(}\PY{p}{[}\PY{p}{[}\PY{l+m+mi}{1}\PY{p}{,} \PY{l+m+mi}{2}\PY{p}{]}\PY{p}{,} \PY{p}{[}\PY{l+m+mi}{3}\PY{p}{,} \PY{l+m+mi}{4}\PY{p}{]}\PY{p}{]}\PY{p}{,} \PY{n+nb}{float}\PY{p}{)}\PY{o}{.}\PY{n}{dtype}
\end{Verbatim}
\end{tcolorbox}

            \begin{tcolorbox}[breakable, size=fbox, boxrule=.5pt, pad at break*=1mm, opacityfill=0]
\prompt{Out}{outcolor}{15}{\boxspacing}
\begin{Verbatim}[commandchars=\\\{\}]
dtype('float64')
\end{Verbatim}
\end{tcolorbox}
        
    The array \texttt{x} is two-dimensional. We can find out the number of
rows and columns by looking at its \texttt{shape} attribute.

    \begin{tcolorbox}[breakable, size=fbox, boxrule=1pt, pad at break*=1mm,colback=cellbackground, colframe=cellborder]
\prompt{In}{incolor}{16}{\boxspacing}
\begin{Verbatim}[commandchars=\\\{\}]
\PY{n}{x}\PY{o}{.}\PY{n}{shape}
\end{Verbatim}
\end{tcolorbox}

            \begin{tcolorbox}[breakable, size=fbox, boxrule=.5pt, pad at break*=1mm, opacityfill=0]
\prompt{Out}{outcolor}{16}{\boxspacing}
\begin{Verbatim}[commandchars=\\\{\}]
(2, 2)
\end{Verbatim}
\end{tcolorbox}
        
    A \emph{method} is a function that is associated with an object. For
instance, given an array \texttt{x}, the expression \texttt{x.sum()}
sums all of its elements, using the \texttt{sum()} method for arrays.
The call \texttt{x.sum()} automatically provides \texttt{x} as the first
argument to its \texttt{sum()} method.

    \begin{tcolorbox}[breakable, size=fbox, boxrule=1pt, pad at break*=1mm,colback=cellbackground, colframe=cellborder]
\prompt{In}{incolor}{17}{\boxspacing}
\begin{Verbatim}[commandchars=\\\{\}]
\PY{n}{x} \PY{o}{=} \PY{n}{np}\PY{o}{.}\PY{n}{array}\PY{p}{(}\PY{p}{[}\PY{l+m+mi}{1}\PY{p}{,} \PY{l+m+mi}{2}\PY{p}{,} \PY{l+m+mi}{3}\PY{p}{,} \PY{l+m+mi}{4}\PY{p}{]}\PY{p}{)}
\PY{n}{x}\PY{o}{.}\PY{n}{sum}\PY{p}{(}\PY{p}{)}
\end{Verbatim}
\end{tcolorbox}

            \begin{tcolorbox}[breakable, size=fbox, boxrule=.5pt, pad at break*=1mm, opacityfill=0]
\prompt{Out}{outcolor}{17}{\boxspacing}
\begin{Verbatim}[commandchars=\\\{\}]
10
\end{Verbatim}
\end{tcolorbox}
        
    We could also sum the elements of \texttt{x} by passing in \texttt{x} as
an argument to the \texttt{np.sum()} function.

    \begin{tcolorbox}[breakable, size=fbox, boxrule=1pt, pad at break*=1mm,colback=cellbackground, colframe=cellborder]
\prompt{In}{incolor}{18}{\boxspacing}
\begin{Verbatim}[commandchars=\\\{\}]
\PY{n}{x} \PY{o}{=} \PY{n}{np}\PY{o}{.}\PY{n}{array}\PY{p}{(}\PY{p}{[}\PY{l+m+mi}{1}\PY{p}{,} \PY{l+m+mi}{2}\PY{p}{,} \PY{l+m+mi}{3}\PY{p}{,} \PY{l+m+mi}{4}\PY{p}{]}\PY{p}{)}
\PY{n}{np}\PY{o}{.}\PY{n}{sum}\PY{p}{(}\PY{n}{x}\PY{p}{)}
\end{Verbatim}
\end{tcolorbox}

            \begin{tcolorbox}[breakable, size=fbox, boxrule=.5pt, pad at break*=1mm, opacityfill=0]
\prompt{Out}{outcolor}{18}{\boxspacing}
\begin{Verbatim}[commandchars=\\\{\}]
10
\end{Verbatim}
\end{tcolorbox}
        
    As another example, the \texttt{reshape()} method returns a new array
with the same elements as \texttt{x}, but a different shape. We do this
by passing in a \texttt{tuple} in our call to \texttt{reshape()}, in
this case \texttt{(2,\ 3)}. This tuple specifies that we would like to
create a two-dimensional array with \(2\) rows and \(3\) columns. \{Like
lists, tuples represent a sequence of objects. Why do we need more than
one way to create a sequence? There are a few differences between tuples
and lists, but perhaps the most important is that elements of a tuple
cannot be modified, whereas elements of a list can be.\}

In what follows, the \texttt{\textbackslash{}n} character creates a
\emph{new line}.

    \begin{tcolorbox}[breakable, size=fbox, boxrule=1pt, pad at break*=1mm,colback=cellbackground, colframe=cellborder]
\prompt{In}{incolor}{19}{\boxspacing}
\begin{Verbatim}[commandchars=\\\{\}]
\PY{n}{x} \PY{o}{=} \PY{n}{np}\PY{o}{.}\PY{n}{array}\PY{p}{(}\PY{p}{[}\PY{l+m+mi}{1}\PY{p}{,} \PY{l+m+mi}{2}\PY{p}{,} \PY{l+m+mi}{3}\PY{p}{,} \PY{l+m+mi}{4}\PY{p}{,} \PY{l+m+mi}{5}\PY{p}{,} \PY{l+m+mi}{6}\PY{p}{]}\PY{p}{)}
\PY{n+nb}{print}\PY{p}{(}\PY{l+s+s1}{\PYZsq{}}\PY{l+s+s1}{beginning x:}\PY{l+s+se}{\PYZbs{}n}\PY{l+s+s1}{\PYZsq{}}\PY{p}{,} \PY{n}{x}\PY{p}{)}
\PY{n}{x\PYZus{}reshape} \PY{o}{=} \PY{n}{x}\PY{o}{.}\PY{n}{reshape}\PY{p}{(}\PY{p}{(}\PY{l+m+mi}{2}\PY{p}{,} \PY{l+m+mi}{3}\PY{p}{)}\PY{p}{)}
\PY{n+nb}{print}\PY{p}{(}\PY{l+s+s1}{\PYZsq{}}\PY{l+s+s1}{reshaped x:}\PY{l+s+se}{\PYZbs{}n}\PY{l+s+s1}{\PYZsq{}}\PY{p}{,} \PY{n}{x\PYZus{}reshape}\PY{p}{)}
\end{Verbatim}
\end{tcolorbox}

    \begin{Verbatim}[commandchars=\\\{\}]
beginning x:
 [1 2 3 4 5 6]
reshaped x:
 [[1 2 3]
 [4 5 6]]
    \end{Verbatim}

    The previous output reveals that \texttt{numpy} arrays are specified as
a sequence of \emph{rows}. This is called \emph{row-major ordering}, as
opposed to \emph{column-major ordering}.

    \texttt{Python} (and hence \texttt{numpy}) uses 0-based indexing. This
means that to access the top left element of \texttt{x\_reshape}, we
type in \texttt{x\_reshape{[}0,0{]}}.

    \begin{tcolorbox}[breakable, size=fbox, boxrule=1pt, pad at break*=1mm,colback=cellbackground, colframe=cellborder]
\prompt{In}{incolor}{20}{\boxspacing}
\begin{Verbatim}[commandchars=\\\{\}]
\PY{n}{x\PYZus{}reshape}\PY{p}{[}\PY{l+m+mi}{0}\PY{p}{,} \PY{l+m+mi}{0}\PY{p}{]} 
\end{Verbatim}
\end{tcolorbox}

            \begin{tcolorbox}[breakable, size=fbox, boxrule=.5pt, pad at break*=1mm, opacityfill=0]
\prompt{Out}{outcolor}{20}{\boxspacing}
\begin{Verbatim}[commandchars=\\\{\}]
1
\end{Verbatim}
\end{tcolorbox}
        
    Similarly, \texttt{x\_reshape{[}1,2{]}} yields the element in the second
row and the third column of \texttt{x\_reshape}.

    \begin{tcolorbox}[breakable, size=fbox, boxrule=1pt, pad at break*=1mm,colback=cellbackground, colframe=cellborder]
\prompt{In}{incolor}{21}{\boxspacing}
\begin{Verbatim}[commandchars=\\\{\}]
\PY{n}{x\PYZus{}reshape}\PY{p}{[}\PY{l+m+mi}{1}\PY{p}{,} \PY{l+m+mi}{2}\PY{p}{]} 
\end{Verbatim}
\end{tcolorbox}

            \begin{tcolorbox}[breakable, size=fbox, boxrule=.5pt, pad at break*=1mm, opacityfill=0]
\prompt{Out}{outcolor}{21}{\boxspacing}
\begin{Verbatim}[commandchars=\\\{\}]
6
\end{Verbatim}
\end{tcolorbox}
        
    Similarly, \texttt{x{[}2{]}} yields the third entry of \texttt{x}.

Now, let's modify the top left element of \texttt{x\_reshape}. To our
surprise, we discover that the first element of \texttt{x} has been
modified as well!

    \begin{tcolorbox}[breakable, size=fbox, boxrule=1pt, pad at break*=1mm,colback=cellbackground, colframe=cellborder]
\prompt{In}{incolor}{22}{\boxspacing}
\begin{Verbatim}[commandchars=\\\{\}]
\PY{n+nb}{print}\PY{p}{(}\PY{l+s+s1}{\PYZsq{}}\PY{l+s+s1}{x before we modify x\PYZus{}reshape:}\PY{l+s+se}{\PYZbs{}n}\PY{l+s+s1}{\PYZsq{}}\PY{p}{,} \PY{n}{x}\PY{p}{)}
\PY{n+nb}{print}\PY{p}{(}\PY{l+s+s1}{\PYZsq{}}\PY{l+s+s1}{x\PYZus{}reshape before we modify x\PYZus{}reshape:}\PY{l+s+se}{\PYZbs{}n}\PY{l+s+s1}{\PYZsq{}}\PY{p}{,} \PY{n}{x\PYZus{}reshape}\PY{p}{)}
\PY{n}{x\PYZus{}reshape}\PY{p}{[}\PY{l+m+mi}{0}\PY{p}{,} \PY{l+m+mi}{0}\PY{p}{]} \PY{o}{=} \PY{l+m+mi}{5}
\PY{n+nb}{print}\PY{p}{(}\PY{l+s+s1}{\PYZsq{}}\PY{l+s+s1}{x\PYZus{}reshape after we modify its top left element:}\PY{l+s+se}{\PYZbs{}n}\PY{l+s+s1}{\PYZsq{}}\PY{p}{,} \PY{n}{x\PYZus{}reshape}\PY{p}{)}
\PY{n+nb}{print}\PY{p}{(}\PY{l+s+s1}{\PYZsq{}}\PY{l+s+s1}{x after we modify top left element of x\PYZus{}reshape:}\PY{l+s+se}{\PYZbs{}n}\PY{l+s+s1}{\PYZsq{}}\PY{p}{,} \PY{n}{x}\PY{p}{)}
\end{Verbatim}
\end{tcolorbox}

    \begin{Verbatim}[commandchars=\\\{\}]
x before we modify x\_reshape:
 [1 2 3 4 5 6]
x\_reshape before we modify x\_reshape:
 [[1 2 3]
 [4 5 6]]
x\_reshape after we modify its top left element:
 [[5 2 3]
 [4 5 6]]
x after we modify top left element of x\_reshape:
 [5 2 3 4 5 6]
    \end{Verbatim}

    Modifying \texttt{x\_reshape} also modified \texttt{x} because the two
objects occupy the same space in memory.

    We just saw that we can modify an element of an array. Can we also
modify a tuple? It turns out that we cannot --- and trying to do so
introduces an \emph{exception}, or error.

    \begin{tcolorbox}[breakable, size=fbox, boxrule=1pt, pad at break*=1mm,colback=cellbackground, colframe=cellborder]
\prompt{In}{incolor}{23}{\boxspacing}
\begin{Verbatim}[commandchars=\\\{\}]
\PY{n}{my\PYZus{}tuple} \PY{o}{=} \PY{p}{(}\PY{l+m+mi}{3}\PY{p}{,} \PY{l+m+mi}{4}\PY{p}{,} \PY{l+m+mi}{5}\PY{p}{)}
\PY{n}{my\PYZus{}tuple}\PY{p}{[}\PY{l+m+mi}{0}\PY{p}{]} \PY{o}{=} \PY{l+m+mi}{2}
\end{Verbatim}
\end{tcolorbox}

    \begin{Verbatim}[commandchars=\\\{\}, frame=single, framerule=2mm, rulecolor=\color{outerrorbackground}]
\textcolor{ansi-red}{---------------------------------------------------------------------------}
\textcolor{ansi-red}{TypeError}                                 Traceback (most recent call last)
\textcolor{ansi-cyan}{Cell}\textcolor{ansi-cyan}{ }\textcolor{ansi-green}{In[23]}\textcolor{ansi-green}{, line 2}
\textcolor{ansi-green}{      1} my\_tuple = (\textcolor{ansi-green}{3}, \textcolor{ansi-green}{4}, \textcolor{ansi-green}{5})
\textcolor{ansi-green}{----> }\textcolor{ansi-green}{2} \setlength{\fboxsep}{0pt}\colorbox{ansi-yellow}{my\_tuple\strut}\setlength{\fboxsep}{0pt}\colorbox{ansi-yellow}{[\strut}\textcolor{ansi-green}{\setlength{\fboxsep}{0pt}\colorbox{ansi-yellow}{0\strut}}\setlength{\fboxsep}{0pt}\colorbox{ansi-yellow}{]\strut} = \textcolor{ansi-green}{2}

\textcolor{ansi-red}{TypeError}: 'tuple' object does not support item assignment
    \end{Verbatim}

    We now briefly mention some attributes of arrays that will come in
handy. An array's \texttt{shape} attribute contains its dimension; this
is always a tuple. The \texttt{ndim} attribute yields the number of
dimensions, and \texttt{T} provides its transpose.

    \begin{tcolorbox}[breakable, size=fbox, boxrule=1pt, pad at break*=1mm,colback=cellbackground, colframe=cellborder]
\prompt{In}{incolor}{24}{\boxspacing}
\begin{Verbatim}[commandchars=\\\{\}]
\PY{n}{x\PYZus{}reshape}\PY{o}{.}\PY{n}{shape}\PY{p}{,} \PY{n}{x\PYZus{}reshape}\PY{o}{.}\PY{n}{ndim}\PY{p}{,} \PY{n}{x\PYZus{}reshape}\PY{o}{.}\PY{n}{T}
\end{Verbatim}
\end{tcolorbox}

            \begin{tcolorbox}[breakable, size=fbox, boxrule=.5pt, pad at break*=1mm, opacityfill=0]
\prompt{Out}{outcolor}{24}{\boxspacing}
\begin{Verbatim}[commandchars=\\\{\}]
((2, 3),
 2,
 array([[5, 4],
        [2, 5],
        [3, 6]]))
\end{Verbatim}
\end{tcolorbox}
        
    Notice that the three individual outputs \texttt{(2,3)}, \texttt{2}, and
\texttt{array({[}{[}5,\ 4{]},{[}2,\ 5{]},\ {[}3,6{]}{]})} are themselves
output as a tuple.

We will often want to apply functions to arrays. For instance, we can
compute the square root of the entries using the \texttt{np.sqrt()}
function:

    \begin{tcolorbox}[breakable, size=fbox, boxrule=1pt, pad at break*=1mm,colback=cellbackground, colframe=cellborder]
\prompt{In}{incolor}{25}{\boxspacing}
\begin{Verbatim}[commandchars=\\\{\}]
\PY{n}{np}\PY{o}{.}\PY{n}{sqrt}\PY{p}{(}\PY{n}{x}\PY{p}{)}
\end{Verbatim}
\end{tcolorbox}

            \begin{tcolorbox}[breakable, size=fbox, boxrule=.5pt, pad at break*=1mm, opacityfill=0]
\prompt{Out}{outcolor}{25}{\boxspacing}
\begin{Verbatim}[commandchars=\\\{\}]
array([2.23606798, 1.41421356, 1.73205081, 2.        , 2.23606798,
       2.44948974])
\end{Verbatim}
\end{tcolorbox}
        
    We can also square the elements:

    \begin{tcolorbox}[breakable, size=fbox, boxrule=1pt, pad at break*=1mm,colback=cellbackground, colframe=cellborder]
\prompt{In}{incolor}{26}{\boxspacing}
\begin{Verbatim}[commandchars=\\\{\}]
\PY{n}{x}\PY{o}{*}\PY{o}{*}\PY{l+m+mi}{2}
\end{Verbatim}
\end{tcolorbox}

            \begin{tcolorbox}[breakable, size=fbox, boxrule=.5pt, pad at break*=1mm, opacityfill=0]
\prompt{Out}{outcolor}{26}{\boxspacing}
\begin{Verbatim}[commandchars=\\\{\}]
array([25,  4,  9, 16, 25, 36])
\end{Verbatim}
\end{tcolorbox}
        
    We can compute the square roots using the same notation, raising to the
power of \(1/2\) instead of 2.

    \begin{tcolorbox}[breakable, size=fbox, boxrule=1pt, pad at break*=1mm,colback=cellbackground, colframe=cellborder]
\prompt{In}{incolor}{27}{\boxspacing}
\begin{Verbatim}[commandchars=\\\{\}]
\PY{n}{x}\PY{o}{*}\PY{o}{*}\PY{l+m+mf}{0.5}
\end{Verbatim}
\end{tcolorbox}

            \begin{tcolorbox}[breakable, size=fbox, boxrule=.5pt, pad at break*=1mm, opacityfill=0]
\prompt{Out}{outcolor}{27}{\boxspacing}
\begin{Verbatim}[commandchars=\\\{\}]
array([2.23606798, 1.41421356, 1.73205081, 2.        , 2.23606798,
       2.44948974])
\end{Verbatim}
\end{tcolorbox}
        
    Throughout this book, we will often want to generate random data. The
\texttt{np.random.normal()} function generates a vector of random normal
variables. We can learn more about this function by looking at the help
page, via a call to \texttt{np.random.normal?}. The first line of the
help page reads \texttt{normal(loc=0.0,\ scale=1.0,\ size=None)}. This
\emph{signature} line tells us that the function's arguments are
\texttt{loc}, \texttt{scale}, and \texttt{size}. These are
\emph{keyword} arguments, which means that when they are passed into the
function, they can be referred to by name (in any order).
\{\texttt{Python} also uses \emph{positional} arguments. Positional
arguments do not need to use a keyword. To see an example, type in
\texttt{np.sum?}. We see that \texttt{a} is a positional argument,
i.e.~this function assumes that the first unnamed argument that it
receives is the array to be summed. By contrast, \texttt{axis} and
\texttt{dtype} are keyword arguments: the position in which these
arguments are entered into \texttt{np.sum()} does not matter.\} By
default, this function will generate random normal variable(s) with mean
(\texttt{loc}) \(0\) and standard deviation (\texttt{scale}) \(1\);
furthermore, a single random variable will be generated unless the
argument to \texttt{size} is changed.

We now generate 50 independent random variables from a \(N(0,1)\)
distribution.

    \begin{tcolorbox}[breakable, size=fbox, boxrule=1pt, pad at break*=1mm,colback=cellbackground, colframe=cellborder]
\prompt{In}{incolor}{28}{\boxspacing}
\begin{Verbatim}[commandchars=\\\{\}]
\PY{n}{x} \PY{o}{=} \PY{n}{np}\PY{o}{.}\PY{n}{random}\PY{o}{.}\PY{n}{normal}\PY{p}{(}\PY{n}{size}\PY{o}{=}\PY{l+m+mi}{50}\PY{p}{)}
\PY{n}{x}
\end{Verbatim}
\end{tcolorbox}

            \begin{tcolorbox}[breakable, size=fbox, boxrule=.5pt, pad at break*=1mm, opacityfill=0]
\prompt{Out}{outcolor}{28}{\boxspacing}
\begin{Verbatim}[commandchars=\\\{\}]
array([-0.06755286, -0.59515975, -0.59124445,  0.93345285,  0.50593652,
       -1.76907045, -0.25741405,  0.87796801,  0.29867187,  1.19609283,
        0.16468045,  0.11386838, -1.82591518, -0.53335271, -0.85678655,
       -0.19322137,  0.65343707,  0.47160945,  0.33741741,  0.58690539,
       -1.00195451,  1.32539572, -0.84977831,  0.6590251 ,  0.4127158 ,
       -0.01590489, -0.33792676, -1.0412565 ,  0.57664837, -0.89242264,
       -0.03146975, -2.94619409,  0.84020772, -0.51395234, -1.48913189,
       -1.67531694,  0.34131073,  0.01197368, -0.13921446,  0.37209259,
       -1.1667784 , -0.14614259, -1.11819663,  1.43098144,  0.22207642,
        0.64787668, -0.06510998, -1.15919834, -0.84805457, -0.73155173])
\end{Verbatim}
\end{tcolorbox}
        
    We create an array \texttt{y} by adding an independent \(N(50,1)\)
random variable to each element of \texttt{x}.

    \begin{tcolorbox}[breakable, size=fbox, boxrule=1pt, pad at break*=1mm,colback=cellbackground, colframe=cellborder]
\prompt{In}{incolor}{29}{\boxspacing}
\begin{Verbatim}[commandchars=\\\{\}]
\PY{n}{y} \PY{o}{=} \PY{n}{x} \PY{o}{+} \PY{n}{np}\PY{o}{.}\PY{n}{random}\PY{o}{.}\PY{n}{normal}\PY{p}{(}\PY{n}{loc}\PY{o}{=}\PY{l+m+mi}{50}\PY{p}{,} \PY{n}{scale}\PY{o}{=}\PY{l+m+mi}{1}\PY{p}{,} \PY{n}{size}\PY{o}{=}\PY{l+m+mi}{50}\PY{p}{)}
\end{Verbatim}
\end{tcolorbox}

    The \texttt{np.corrcoef()} function computes the correlation matrix
between \texttt{x} and \texttt{y}. The off-diagonal elements give the
correlation between \texttt{x} and \texttt{y}.

    \begin{tcolorbox}[breakable, size=fbox, boxrule=1pt, pad at break*=1mm,colback=cellbackground, colframe=cellborder]
\prompt{In}{incolor}{30}{\boxspacing}
\begin{Verbatim}[commandchars=\\\{\}]
\PY{n}{np}\PY{o}{.}\PY{n}{corrcoef}\PY{p}{(}\PY{n}{x}\PY{p}{,} \PY{n}{y}\PY{p}{)}
\end{Verbatim}
\end{tcolorbox}

            \begin{tcolorbox}[breakable, size=fbox, boxrule=.5pt, pad at break*=1mm, opacityfill=0]
\prompt{Out}{outcolor}{30}{\boxspacing}
\begin{Verbatim}[commandchars=\\\{\}]
array([[1.        , 0.56659826],
       [0.56659826, 1.        ]])
\end{Verbatim}
\end{tcolorbox}
        
    If you're following along in your own \texttt{Jupyter} notebook, then
you probably noticed that you got a different set of results when you
ran the past few commands. In particular, each time we call
\texttt{np.random.normal()}, we will get a different answer, as shown in
the following example.

    \begin{tcolorbox}[breakable, size=fbox, boxrule=1pt, pad at break*=1mm,colback=cellbackground, colframe=cellborder]
\prompt{In}{incolor}{31}{\boxspacing}
\begin{Verbatim}[commandchars=\\\{\}]
\PY{n+nb}{print}\PY{p}{(}\PY{n}{np}\PY{o}{.}\PY{n}{random}\PY{o}{.}\PY{n}{normal}\PY{p}{(}\PY{n}{scale}\PY{o}{=}\PY{l+m+mi}{5}\PY{p}{,} \PY{n}{size}\PY{o}{=}\PY{l+m+mi}{2}\PY{p}{)}\PY{p}{)}
\PY{n+nb}{print}\PY{p}{(}\PY{n}{np}\PY{o}{.}\PY{n}{random}\PY{o}{.}\PY{n}{normal}\PY{p}{(}\PY{n}{scale}\PY{o}{=}\PY{l+m+mi}{5}\PY{p}{,} \PY{n}{size}\PY{o}{=}\PY{l+m+mi}{2}\PY{p}{)}\PY{p}{)} 
\end{Verbatim}
\end{tcolorbox}

    \begin{Verbatim}[commandchars=\\\{\}]
[-1.07468791 -1.07888433]
[ 3.84695783 -4.83482283]
    \end{Verbatim}

    

    In order to ensure that our code provides exactly the same results each
time it is run, we can set a \emph{random seed} using the
\texttt{np.random.default\_rng()} function. This function takes an
arbitrary, user-specified integer argument. If we set a random seed
before generating random data, then re-running our code will yield the
same results. The object \texttt{rng} has essentially all the random
number generating methods found in \texttt{np.random}. Hence, to
generate normal data we use \texttt{rng.normal()}.

    \begin{tcolorbox}[breakable, size=fbox, boxrule=1pt, pad at break*=1mm,colback=cellbackground, colframe=cellborder]
\prompt{In}{incolor}{32}{\boxspacing}
\begin{Verbatim}[commandchars=\\\{\}]
\PY{n}{rng} \PY{o}{=} \PY{n}{np}\PY{o}{.}\PY{n}{random}\PY{o}{.}\PY{n}{default\PYZus{}rng}\PY{p}{(}\PY{l+m+mi}{1303}\PY{p}{)}
\PY{n+nb}{print}\PY{p}{(}\PY{n}{rng}\PY{o}{.}\PY{n}{normal}\PY{p}{(}\PY{n}{scale}\PY{o}{=}\PY{l+m+mi}{5}\PY{p}{,} \PY{n}{size}\PY{o}{=}\PY{l+m+mi}{2}\PY{p}{)}\PY{p}{)}
\PY{n}{rng2} \PY{o}{=} \PY{n}{np}\PY{o}{.}\PY{n}{random}\PY{o}{.}\PY{n}{default\PYZus{}rng}\PY{p}{(}\PY{l+m+mi}{1303}\PY{p}{)}
\PY{n+nb}{print}\PY{p}{(}\PY{n}{rng2}\PY{o}{.}\PY{n}{normal}\PY{p}{(}\PY{n}{scale}\PY{o}{=}\PY{l+m+mi}{5}\PY{p}{,} \PY{n}{size}\PY{o}{=}\PY{l+m+mi}{2}\PY{p}{)}\PY{p}{)} 
\end{Verbatim}
\end{tcolorbox}

    \begin{Verbatim}[commandchars=\\\{\}]
[ 4.09482632 -1.07485605]
[ 4.09482632 -1.07485605]
    \end{Verbatim}

    Throughout the labs in this book, we use
\texttt{np.random.default\_rng()} whenever we perform calculations
involving random quantities within \texttt{numpy}. In principle, this
should enable the reader to exactly reproduce the stated results.
However, as new versions of \texttt{numpy} become available, it is
possible that some small discrepancies may occur between the output in
the labs and the output from \texttt{numpy}.

The \texttt{np.mean()}, \texttt{np.var()}, and \texttt{np.std()}
functions can be used to compute the mean, variance, and standard
deviation of arrays. These functions are also available as methods on
the arrays.

    \begin{tcolorbox}[breakable, size=fbox, boxrule=1pt, pad at break*=1mm,colback=cellbackground, colframe=cellborder]
\prompt{In}{incolor}{33}{\boxspacing}
\begin{Verbatim}[commandchars=\\\{\}]
\PY{n}{rng} \PY{o}{=} \PY{n}{np}\PY{o}{.}\PY{n}{random}\PY{o}{.}\PY{n}{default\PYZus{}rng}\PY{p}{(}\PY{l+m+mi}{3}\PY{p}{)}
\PY{n}{y} \PY{o}{=} \PY{n}{rng}\PY{o}{.}\PY{n}{standard\PYZus{}normal}\PY{p}{(}\PY{l+m+mi}{10}\PY{p}{)}
\PY{n}{np}\PY{o}{.}\PY{n}{mean}\PY{p}{(}\PY{n}{y}\PY{p}{)}\PY{p}{,} \PY{n}{y}\PY{o}{.}\PY{n}{mean}\PY{p}{(}\PY{p}{)}
\end{Verbatim}
\end{tcolorbox}

            \begin{tcolorbox}[breakable, size=fbox, boxrule=.5pt, pad at break*=1mm, opacityfill=0]
\prompt{Out}{outcolor}{33}{\boxspacing}
\begin{Verbatim}[commandchars=\\\{\}]
(-0.1126795190952861, -0.1126795190952861)
\end{Verbatim}
\end{tcolorbox}
        
    

    \begin{tcolorbox}[breakable, size=fbox, boxrule=1pt, pad at break*=1mm,colback=cellbackground, colframe=cellborder]
\prompt{In}{incolor}{34}{\boxspacing}
\begin{Verbatim}[commandchars=\\\{\}]
\PY{n}{np}\PY{o}{.}\PY{n}{var}\PY{p}{(}\PY{n}{y}\PY{p}{)}\PY{p}{,} \PY{n}{y}\PY{o}{.}\PY{n}{var}\PY{p}{(}\PY{p}{)}\PY{p}{,} \PY{n}{np}\PY{o}{.}\PY{n}{mean}\PY{p}{(}\PY{p}{(}\PY{n}{y} \PY{o}{\PYZhy{}} \PY{n}{y}\PY{o}{.}\PY{n}{mean}\PY{p}{(}\PY{p}{)}\PY{p}{)}\PY{o}{*}\PY{o}{*}\PY{l+m+mi}{2}\PY{p}{)}
\end{Verbatim}
\end{tcolorbox}

            \begin{tcolorbox}[breakable, size=fbox, boxrule=.5pt, pad at break*=1mm, opacityfill=0]
\prompt{Out}{outcolor}{34}{\boxspacing}
\begin{Verbatim}[commandchars=\\\{\}]
(2.7243406406465125, 2.7243406406465125, 2.7243406406465125)
\end{Verbatim}
\end{tcolorbox}
        
    Notice that by default \texttt{np.var()} divides by the sample size
\(n\) rather than \(n-1\); see the \texttt{ddof} argument in
\texttt{np.var?}.

    \begin{tcolorbox}[breakable, size=fbox, boxrule=1pt, pad at break*=1mm,colback=cellbackground, colframe=cellborder]
\prompt{In}{incolor}{35}{\boxspacing}
\begin{Verbatim}[commandchars=\\\{\}]
\PY{n}{np}\PY{o}{.}\PY{n}{sqrt}\PY{p}{(}\PY{n}{np}\PY{o}{.}\PY{n}{var}\PY{p}{(}\PY{n}{y}\PY{p}{)}\PY{p}{)}\PY{p}{,} \PY{n}{np}\PY{o}{.}\PY{n}{std}\PY{p}{(}\PY{n}{y}\PY{p}{)}
\end{Verbatim}
\end{tcolorbox}

            \begin{tcolorbox}[breakable, size=fbox, boxrule=.5pt, pad at break*=1mm, opacityfill=0]
\prompt{Out}{outcolor}{35}{\boxspacing}
\begin{Verbatim}[commandchars=\\\{\}]
(1.6505576756498128, 1.6505576756498128)
\end{Verbatim}
\end{tcolorbox}
        
    The \texttt{np.mean()}, \texttt{np.var()}, and \texttt{np.std()}
functions can also be applied to the rows and columns of a matrix. To
see this, we construct a \(10 \times 3\) matrix of \(N(0,1)\) random
variables, and consider computing its row sums.

    \begin{tcolorbox}[breakable, size=fbox, boxrule=1pt, pad at break*=1mm,colback=cellbackground, colframe=cellborder]
\prompt{In}{incolor}{36}{\boxspacing}
\begin{Verbatim}[commandchars=\\\{\}]
\PY{n}{X} \PY{o}{=} \PY{n}{rng}\PY{o}{.}\PY{n}{standard\PYZus{}normal}\PY{p}{(}\PY{p}{(}\PY{l+m+mi}{10}\PY{p}{,} \PY{l+m+mi}{3}\PY{p}{)}\PY{p}{)}
\PY{n}{X}
\end{Verbatim}
\end{tcolorbox}

            \begin{tcolorbox}[breakable, size=fbox, boxrule=.5pt, pad at break*=1mm, opacityfill=0]
\prompt{Out}{outcolor}{36}{\boxspacing}
\begin{Verbatim}[commandchars=\\\{\}]
array([[ 0.22578661, -0.35263079, -0.28128742],
       [-0.66804635, -1.05515055, -0.39080098],
       [ 0.48194539, -0.23855361,  0.9577587 ],
       [-0.19980213,  0.02425957,  1.54582085],
       [ 0.54510552, -0.50522874, -0.18283897],
       [ 0.54052513,  1.93508803, -0.26962033],
       [-0.24355868,  1.0023136 , -0.88645994],
       [-0.29172023,  0.88253897,  0.58035002],
       [ 0.0915167 ,  0.67010435, -2.82816231],
       [ 1.02130682, -0.95964476, -1.66861984]])
\end{Verbatim}
\end{tcolorbox}
        
    Since arrays are row-major ordered, the first axis,
i.e.~\texttt{axis=0}, refers to its rows. We pass this argument into the
\texttt{mean()} method for the object \texttt{X}.

    \begin{tcolorbox}[breakable, size=fbox, boxrule=1pt, pad at break*=1mm,colback=cellbackground, colframe=cellborder]
\prompt{In}{incolor}{37}{\boxspacing}
\begin{Verbatim}[commandchars=\\\{\}]
\PY{n}{X}\PY{o}{.}\PY{n}{mean}\PY{p}{(}\PY{n}{axis}\PY{o}{=}\PY{l+m+mi}{0}\PY{p}{)}
\end{Verbatim}
\end{tcolorbox}

            \begin{tcolorbox}[breakable, size=fbox, boxrule=.5pt, pad at break*=1mm, opacityfill=0]
\prompt{Out}{outcolor}{37}{\boxspacing}
\begin{Verbatim}[commandchars=\\\{\}]
array([ 0.15030588,  0.14030961, -0.34238602])
\end{Verbatim}
\end{tcolorbox}
        
    The following yields the same result.

    \begin{tcolorbox}[breakable, size=fbox, boxrule=1pt, pad at break*=1mm,colback=cellbackground, colframe=cellborder]
\prompt{In}{incolor}{38}{\boxspacing}
\begin{Verbatim}[commandchars=\\\{\}]
\PY{n}{X}\PY{o}{.}\PY{n}{mean}\PY{p}{(}\PY{l+m+mi}{0}\PY{p}{)}
\end{Verbatim}
\end{tcolorbox}

            \begin{tcolorbox}[breakable, size=fbox, boxrule=.5pt, pad at break*=1mm, opacityfill=0]
\prompt{Out}{outcolor}{38}{\boxspacing}
\begin{Verbatim}[commandchars=\\\{\}]
array([ 0.15030588,  0.14030961, -0.34238602])
\end{Verbatim}
\end{tcolorbox}
        
    

    \hypertarget{graphics}{%
\subsection{Graphics}\label{graphics}}

In \texttt{Python}, common practice is to use the library
\texttt{matplotlib} for graphics. However, since \texttt{Python} was not
written with data analysis in mind, the notion of plotting is not
intrinsic to the language. We will use the \texttt{subplots()} function
from \texttt{matplotlib.pyplot} to create a figure and the axes onto
which we plot our data. For many more examples of how to make plots in
\texttt{Python}, readers are encouraged to visit
\href{https://matplotlib.org/stable/gallery/index.html}{matplotlib.org/stable/gallery/}.

In \texttt{matplotlib}, a plot consists of a \emph{figure} and one or
more \emph{axes}. You can think of the figure as the blank canvas upon
which one or more plots will be displayed: it is the entire plotting
window. The \emph{axes} contain important information about each plot,
such as its \(x\)- and \(y\)-axis labels, title, and more. (Note that in
\texttt{matplotlib}, the word \emph{axes} is not the plural of
\emph{axis}: a plot's \emph{axes} contains much more information than
just the \(x\)-axis and the \(y\)-axis.)

We begin by importing the \texttt{subplots()} function from
\texttt{matplotlib}. We use this function throughout when creating
figures. The function returns a tuple of length two: a figure object as
well as the relevant axes object. We will typically pass
\texttt{figsize} as a keyword argument. Having created our axes, we
attempt our first plot using its \texttt{plot()} method. To learn more
about it, type \texttt{ax.plot?}.

    \begin{tcolorbox}[breakable, size=fbox, boxrule=1pt, pad at break*=1mm,colback=cellbackground, colframe=cellborder]
\prompt{In}{incolor}{39}{\boxspacing}
\begin{Verbatim}[commandchars=\\\{\}]
\PY{k+kn}{from}\PY{+w}{ }\PY{n+nn}{matplotlib}\PY{n+nn}{.}\PY{n+nn}{pyplot}\PY{+w}{ }\PY{k+kn}{import} \PY{n}{subplots}
\PY{n}{fig}\PY{p}{,} \PY{n}{ax} \PY{o}{=} \PY{n}{subplots}\PY{p}{(}\PY{n}{figsize}\PY{o}{=}\PY{p}{(}\PY{l+m+mi}{8}\PY{p}{,} \PY{l+m+mi}{8}\PY{p}{)}\PY{p}{)}
\PY{n}{x} \PY{o}{=} \PY{n}{rng}\PY{o}{.}\PY{n}{standard\PYZus{}normal}\PY{p}{(}\PY{l+m+mi}{100}\PY{p}{)}
\PY{n}{y} \PY{o}{=} \PY{n}{rng}\PY{o}{.}\PY{n}{standard\PYZus{}normal}\PY{p}{(}\PY{l+m+mi}{100}\PY{p}{)}
\PY{n}{ax}\PY{o}{.}\PY{n}{plot}\PY{p}{(}\PY{n}{x}\PY{p}{,} \PY{n}{y}\PY{p}{)}\PY{p}{;}
\end{Verbatim}
\end{tcolorbox}

    \begin{center}
    \adjustimage{max size={0.9\linewidth}{0.9\paperheight}}{artifacts/latex/Ch02-statlearn-lab_files/artifacts/latex/Ch02-statlearn-lab_96_0.png}
    \end{center}
    { \hspace*{\fill} \\}
    
    We pause here to note that we have \emph{unpacked} the tuple of length
two returned by \texttt{subplots()} into the two distinct variables
\texttt{fig} and \texttt{ax}. Unpacking is typically preferred to the
following equivalent but slightly more verbose code:

    \begin{tcolorbox}[breakable, size=fbox, boxrule=1pt, pad at break*=1mm,colback=cellbackground, colframe=cellborder]
\prompt{In}{incolor}{40}{\boxspacing}
\begin{Verbatim}[commandchars=\\\{\}]
\PY{n}{output} \PY{o}{=} \PY{n}{subplots}\PY{p}{(}\PY{n}{figsize}\PY{o}{=}\PY{p}{(}\PY{l+m+mi}{8}\PY{p}{,} \PY{l+m+mi}{8}\PY{p}{)}\PY{p}{)}
\PY{n}{fig} \PY{o}{=} \PY{n}{output}\PY{p}{[}\PY{l+m+mi}{0}\PY{p}{]}
\PY{n}{ax} \PY{o}{=} \PY{n}{output}\PY{p}{[}\PY{l+m+mi}{1}\PY{p}{]}
\end{Verbatim}
\end{tcolorbox}

    \begin{center}
    \adjustimage{max size={0.9\linewidth}{0.9\paperheight}}{artifacts/latex/Ch02-statlearn-lab_files/artifacts/latex/Ch02-statlearn-lab_98_0.png}
    \end{center}
    { \hspace*{\fill} \\}
    
    We see that our earlier cell produced a line plot, which is the default.
To create a scatterplot, we provide an additional argument to
\texttt{ax.plot()}, indicating that circles should be displayed.

    \begin{tcolorbox}[breakable, size=fbox, boxrule=1pt, pad at break*=1mm,colback=cellbackground, colframe=cellborder]
\prompt{In}{incolor}{41}{\boxspacing}
\begin{Verbatim}[commandchars=\\\{\}]
\PY{n}{fig}\PY{p}{,} \PY{n}{ax} \PY{o}{=} \PY{n}{subplots}\PY{p}{(}\PY{n}{figsize}\PY{o}{=}\PY{p}{(}\PY{l+m+mi}{8}\PY{p}{,} \PY{l+m+mi}{8}\PY{p}{)}\PY{p}{)}
\PY{n}{ax}\PY{o}{.}\PY{n}{plot}\PY{p}{(}\PY{n}{x}\PY{p}{,} \PY{n}{y}\PY{p}{,} \PY{l+s+s1}{\PYZsq{}}\PY{l+s+s1}{o}\PY{l+s+s1}{\PYZsq{}}\PY{p}{)}\PY{p}{;}
\end{Verbatim}
\end{tcolorbox}

    \begin{center}
    \adjustimage{max size={0.9\linewidth}{0.9\paperheight}}{artifacts/latex/Ch02-statlearn-lab_files/artifacts/latex/Ch02-statlearn-lab_100_0.png}
    \end{center}
    { \hspace*{\fill} \\}
    
    Different values of this additional argument can be used to produce
different colored lines as well as different linestyles.

    As an alternative, we could use the \texttt{ax.scatter()} function to
create a scatterplot.

    \begin{tcolorbox}[breakable, size=fbox, boxrule=1pt, pad at break*=1mm,colback=cellbackground, colframe=cellborder]
\prompt{In}{incolor}{42}{\boxspacing}
\begin{Verbatim}[commandchars=\\\{\}]
\PY{n}{fig}\PY{p}{,} \PY{n}{ax} \PY{o}{=} \PY{n}{subplots}\PY{p}{(}\PY{n}{figsize}\PY{o}{=}\PY{p}{(}\PY{l+m+mi}{8}\PY{p}{,} \PY{l+m+mi}{8}\PY{p}{)}\PY{p}{)}
\PY{n}{ax}\PY{o}{.}\PY{n}{scatter}\PY{p}{(}\PY{n}{x}\PY{p}{,} \PY{n}{y}\PY{p}{,} \PY{n}{marker}\PY{o}{=}\PY{l+s+s1}{\PYZsq{}}\PY{l+s+s1}{o}\PY{l+s+s1}{\PYZsq{}}\PY{p}{)}\PY{p}{;}
\end{Verbatim}
\end{tcolorbox}

    \begin{center}
    \adjustimage{max size={0.9\linewidth}{0.9\paperheight}}{artifacts/latex/Ch02-statlearn-lab_files/artifacts/latex/Ch02-statlearn-lab_103_0.png}
    \end{center}
    { \hspace*{\fill} \\}
    
    Notice that in the code blocks above, we have ended the last line with a
semicolon. This prevents \texttt{ax.plot(x,\ y)} from printing text to
the notebook. However, it does not prevent a plot from being produced.
If we omit the trailing semi-colon, then we obtain the following output:

    \begin{tcolorbox}[breakable, size=fbox, boxrule=1pt, pad at break*=1mm,colback=cellbackground, colframe=cellborder]
\prompt{In}{incolor}{43}{\boxspacing}
\begin{Verbatim}[commandchars=\\\{\}]
\PY{n}{fig}\PY{p}{,} \PY{n}{ax} \PY{o}{=} \PY{n}{subplots}\PY{p}{(}\PY{n}{figsize}\PY{o}{=}\PY{p}{(}\PY{l+m+mi}{8}\PY{p}{,} \PY{l+m+mi}{8}\PY{p}{)}\PY{p}{)}
\PY{n}{ax}\PY{o}{.}\PY{n}{scatter}\PY{p}{(}\PY{n}{x}\PY{p}{,} \PY{n}{y}\PY{p}{,} \PY{n}{marker}\PY{o}{=}\PY{l+s+s1}{\PYZsq{}}\PY{l+s+s1}{o}\PY{l+s+s1}{\PYZsq{}}\PY{p}{)}
\end{Verbatim}
\end{tcolorbox}

            \begin{tcolorbox}[breakable, size=fbox, boxrule=.5pt, pad at break*=1mm, opacityfill=0]
\prompt{Out}{outcolor}{43}{\boxspacing}
\begin{Verbatim}[commandchars=\\\{\}]
<matplotlib.collections.PathCollection at 0x11843af60>
\end{Verbatim}
\end{tcolorbox}
        
    \begin{center}
    \adjustimage{max size={0.9\linewidth}{0.9\paperheight}}{artifacts/latex/Ch02-statlearn-lab_files/artifacts/latex/Ch02-statlearn-lab_105_1.png}
    \end{center}
    { \hspace*{\fill} \\}
    
    In what follows, we will use trailing semicolons whenever the text that
would be output is not germane to the discussion at hand.

    To label our plot, we make use of the \texttt{set\_xlabel()},
\texttt{set\_ylabel()}, and \texttt{set\_title()} methods of
\texttt{ax}.

    \begin{tcolorbox}[breakable, size=fbox, boxrule=1pt, pad at break*=1mm,colback=cellbackground, colframe=cellborder]
\prompt{In}{incolor}{44}{\boxspacing}
\begin{Verbatim}[commandchars=\\\{\}]
\PY{n}{fig}\PY{p}{,} \PY{n}{ax} \PY{o}{=} \PY{n}{subplots}\PY{p}{(}\PY{n}{figsize}\PY{o}{=}\PY{p}{(}\PY{l+m+mi}{8}\PY{p}{,} \PY{l+m+mi}{8}\PY{p}{)}\PY{p}{)}
\PY{n}{ax}\PY{o}{.}\PY{n}{scatter}\PY{p}{(}\PY{n}{x}\PY{p}{,} \PY{n}{y}\PY{p}{,} \PY{n}{marker}\PY{o}{=}\PY{l+s+s1}{\PYZsq{}}\PY{l+s+s1}{o}\PY{l+s+s1}{\PYZsq{}}\PY{p}{)}
\PY{n}{ax}\PY{o}{.}\PY{n}{set\PYZus{}xlabel}\PY{p}{(}\PY{l+s+s2}{\PYZdq{}}\PY{l+s+s2}{this is the x\PYZhy{}axis}\PY{l+s+s2}{\PYZdq{}}\PY{p}{)}
\PY{n}{ax}\PY{o}{.}\PY{n}{set\PYZus{}ylabel}\PY{p}{(}\PY{l+s+s2}{\PYZdq{}}\PY{l+s+s2}{this is the y\PYZhy{}axis}\PY{l+s+s2}{\PYZdq{}}\PY{p}{)}
\PY{n}{ax}\PY{o}{.}\PY{n}{set\PYZus{}title}\PY{p}{(}\PY{l+s+s2}{\PYZdq{}}\PY{l+s+s2}{Plot of X vs Y}\PY{l+s+s2}{\PYZdq{}}\PY{p}{)}\PY{p}{;}
\end{Verbatim}
\end{tcolorbox}

    \begin{center}
    \adjustimage{max size={0.9\linewidth}{0.9\paperheight}}{artifacts/latex/Ch02-statlearn-lab_files/artifacts/latex/Ch02-statlearn-lab_108_0.png}
    \end{center}
    { \hspace*{\fill} \\}
    
    Having access to the figure object \texttt{fig} itself means that we can
go in and change some aspects and then redisplay it. Here, we change the
size from \texttt{(8,\ 8)} to \texttt{(12,\ 3)}.

    \begin{tcolorbox}[breakable, size=fbox, boxrule=1pt, pad at break*=1mm,colback=cellbackground, colframe=cellborder]
\prompt{In}{incolor}{45}{\boxspacing}
\begin{Verbatim}[commandchars=\\\{\}]
\PY{n}{fig}\PY{o}{.}\PY{n}{set\PYZus{}size\PYZus{}inches}\PY{p}{(}\PY{l+m+mi}{12}\PY{p}{,}\PY{l+m+mi}{3}\PY{p}{)}
\PY{n}{fig}
\end{Verbatim}
\end{tcolorbox}
 
            
\prompt{Out}{outcolor}{45}{}
    
    \begin{center}
    \adjustimage{max size={0.9\linewidth}{0.9\paperheight}}{artifacts/latex/Ch02-statlearn-lab_files/artifacts/latex/Ch02-statlearn-lab_110_0.png}
    \end{center}
    { \hspace*{\fill} \\}
    

    

    Occasionally we will want to create several plots within a figure. This
can be achieved by passing additional arguments to \texttt{subplots()}.
Below, we create a \(2 \times 3\) grid of plots in a figure of size
determined by the \texttt{figsize} argument. In such situations, there
is often a relationship between the axes in the plots. For example, all
plots may have a common \(x\)-axis. The \texttt{subplots()} function can
automatically handle this situation when passed the keyword argument
\texttt{sharex=True}. The \texttt{axes} object below is an array
pointing to different plots in the figure.

    \begin{tcolorbox}[breakable, size=fbox, boxrule=1pt, pad at break*=1mm,colback=cellbackground, colframe=cellborder]
\prompt{In}{incolor}{46}{\boxspacing}
\begin{Verbatim}[commandchars=\\\{\}]
\PY{n}{fig}\PY{p}{,} \PY{n}{axes} \PY{o}{=} \PY{n}{subplots}\PY{p}{(}\PY{n}{nrows}\PY{o}{=}\PY{l+m+mi}{2}\PY{p}{,}
                     \PY{n}{ncols}\PY{o}{=}\PY{l+m+mi}{3}\PY{p}{,}
                     \PY{n}{figsize}\PY{o}{=}\PY{p}{(}\PY{l+m+mi}{15}\PY{p}{,} \PY{l+m+mi}{5}\PY{p}{)}\PY{p}{)}
\end{Verbatim}
\end{tcolorbox}

    \begin{center}
    \adjustimage{max size={0.9\linewidth}{0.9\paperheight}}{artifacts/latex/Ch02-statlearn-lab_files/artifacts/latex/Ch02-statlearn-lab_113_0.png}
    \end{center}
    { \hspace*{\fill} \\}
    
    We now produce a scatter plot with
\texttt{\textquotesingle{}o\textquotesingle{}} in the second column of
the first row and a scatter plot with
\texttt{\textquotesingle{}+\textquotesingle{}} in the third column of
the second row.

    \begin{tcolorbox}[breakable, size=fbox, boxrule=1pt, pad at break*=1mm,colback=cellbackground, colframe=cellborder]
\prompt{In}{incolor}{47}{\boxspacing}
\begin{Verbatim}[commandchars=\\\{\}]
\PY{n}{axes}\PY{p}{[}\PY{l+m+mi}{0}\PY{p}{,}\PY{l+m+mi}{1}\PY{p}{]}\PY{o}{.}\PY{n}{plot}\PY{p}{(}\PY{n}{x}\PY{p}{,} \PY{n}{y}\PY{p}{,} \PY{l+s+s1}{\PYZsq{}}\PY{l+s+s1}{o}\PY{l+s+s1}{\PYZsq{}}\PY{p}{)}
\PY{n}{axes}\PY{p}{[}\PY{l+m+mi}{1}\PY{p}{,}\PY{l+m+mi}{2}\PY{p}{]}\PY{o}{.}\PY{n}{scatter}\PY{p}{(}\PY{n}{x}\PY{p}{,} \PY{n}{y}\PY{p}{,} \PY{n}{marker}\PY{o}{=}\PY{l+s+s1}{\PYZsq{}}\PY{l+s+s1}{+}\PY{l+s+s1}{\PYZsq{}}\PY{p}{)}
\PY{n}{fig}
\end{Verbatim}
\end{tcolorbox}
 
            
\prompt{Out}{outcolor}{47}{}
    
    \begin{center}
    \adjustimage{max size={0.9\linewidth}{0.9\paperheight}}{artifacts/latex/Ch02-statlearn-lab_files/artifacts/latex/Ch02-statlearn-lab_115_0.png}
    \end{center}
    { \hspace*{\fill} \\}
    

    Type \texttt{subplots?} to learn more about \texttt{subplots()}.

    To save the output of \texttt{fig}, we call its \texttt{savefig()}
method. The argument \texttt{dpi} is the dots per inch, used to
determine how large the figure will be in pixels.

    \begin{tcolorbox}[breakable, size=fbox, boxrule=1pt, pad at break*=1mm,colback=cellbackground, colframe=cellborder]
\prompt{In}{incolor}{48}{\boxspacing}
\begin{Verbatim}[commandchars=\\\{\}]
\PY{n}{fig}\PY{o}{.}\PY{n}{savefig}\PY{p}{(}\PY{l+s+s2}{\PYZdq{}}\PY{l+s+s2}{Figure.png}\PY{l+s+s2}{\PYZdq{}}\PY{p}{,} \PY{n}{dpi}\PY{o}{=}\PY{l+m+mi}{400}\PY{p}{)}
\PY{n}{fig}\PY{o}{.}\PY{n}{savefig}\PY{p}{(}\PY{l+s+s2}{\PYZdq{}}\PY{l+s+s2}{Figure.pdf}\PY{l+s+s2}{\PYZdq{}}\PY{p}{,} \PY{n}{dpi}\PY{o}{=}\PY{l+m+mi}{200}\PY{p}{)}\PY{p}{;}
\end{Verbatim}
\end{tcolorbox}

    We can continue to modify \texttt{fig} using step-by-step updates; for
example, we can modify the range of the \(x\)-axis, re-save the figure,
and even re-display it.

    \begin{tcolorbox}[breakable, size=fbox, boxrule=1pt, pad at break*=1mm,colback=cellbackground, colframe=cellborder]
\prompt{In}{incolor}{49}{\boxspacing}
\begin{Verbatim}[commandchars=\\\{\}]
\PY{n}{axes}\PY{p}{[}\PY{l+m+mi}{0}\PY{p}{,}\PY{l+m+mi}{1}\PY{p}{]}\PY{o}{.}\PY{n}{set\PYZus{}xlim}\PY{p}{(}\PY{p}{[}\PY{o}{\PYZhy{}}\PY{l+m+mi}{1}\PY{p}{,}\PY{l+m+mi}{1}\PY{p}{]}\PY{p}{)}
\PY{n}{fig}\PY{o}{.}\PY{n}{savefig}\PY{p}{(}\PY{l+s+s2}{\PYZdq{}}\PY{l+s+s2}{Figure\PYZus{}updated.jpg}\PY{l+s+s2}{\PYZdq{}}\PY{p}{)}
\PY{n}{fig}
\end{Verbatim}
\end{tcolorbox}
 
            
\prompt{Out}{outcolor}{49}{}
    
    \begin{center}
    \adjustimage{max size={0.9\linewidth}{0.9\paperheight}}{artifacts/latex/Ch02-statlearn-lab_files/artifacts/latex/Ch02-statlearn-lab_120_0.png}
    \end{center}
    { \hspace*{\fill} \\}
    

    We now create some more sophisticated plots. The \texttt{ax.contour()}
method produces a \emph{contour plot} in order to represent
three-dimensional data, similar to a topographical map. It takes three
arguments:

\begin{itemize}
\tightlist
\item
  A vector of \texttt{x} values (the first dimension),
\item
  A vector of \texttt{y} values (the second dimension), and
\item
  A matrix whose elements correspond to the \texttt{z} value (the third
  dimension) for each pair of \texttt{(x,y)} coordinates.
\end{itemize}

To create \texttt{x} and \texttt{y}, we'll use the command
\texttt{np.linspace(a,\ b,\ n)}, which returns a vector of \texttt{n}
numbers starting at \texttt{a} and ending at \texttt{b}.

    \begin{tcolorbox}[breakable, size=fbox, boxrule=1pt, pad at break*=1mm,colback=cellbackground, colframe=cellborder]
\prompt{In}{incolor}{50}{\boxspacing}
\begin{Verbatim}[commandchars=\\\{\}]
\PY{n}{fig}\PY{p}{,} \PY{n}{ax} \PY{o}{=} \PY{n}{subplots}\PY{p}{(}\PY{n}{figsize}\PY{o}{=}\PY{p}{(}\PY{l+m+mi}{8}\PY{p}{,} \PY{l+m+mi}{8}\PY{p}{)}\PY{p}{)}
\PY{n}{x} \PY{o}{=} \PY{n}{np}\PY{o}{.}\PY{n}{linspace}\PY{p}{(}\PY{o}{\PYZhy{}}\PY{n}{np}\PY{o}{.}\PY{n}{pi}\PY{p}{,} \PY{n}{np}\PY{o}{.}\PY{n}{pi}\PY{p}{,} \PY{l+m+mi}{50}\PY{p}{)}
\PY{n}{y} \PY{o}{=} \PY{n}{x}
\PY{n}{f} \PY{o}{=} \PY{n}{np}\PY{o}{.}\PY{n}{multiply}\PY{o}{.}\PY{n}{outer}\PY{p}{(}\PY{n}{np}\PY{o}{.}\PY{n}{cos}\PY{p}{(}\PY{n}{y}\PY{p}{)}\PY{p}{,} \PY{l+m+mi}{1} \PY{o}{/} \PY{p}{(}\PY{l+m+mi}{1} \PY{o}{+} \PY{n}{x}\PY{o}{*}\PY{o}{*}\PY{l+m+mi}{2}\PY{p}{)}\PY{p}{)}
\PY{n}{ax}\PY{o}{.}\PY{n}{contour}\PY{p}{(}\PY{n}{x}\PY{p}{,} \PY{n}{y}\PY{p}{,} \PY{n}{f}\PY{p}{)}\PY{p}{;}
\end{Verbatim}
\end{tcolorbox}

    \begin{center}
    \adjustimage{max size={0.9\linewidth}{0.9\paperheight}}{artifacts/latex/Ch02-statlearn-lab_files/artifacts/latex/Ch02-statlearn-lab_122_0.png}
    \end{center}
    { \hspace*{\fill} \\}
    
    We can increase the resolution by adding more levels to the image.

    \begin{tcolorbox}[breakable, size=fbox, boxrule=1pt, pad at break*=1mm,colback=cellbackground, colframe=cellborder]
\prompt{In}{incolor}{51}{\boxspacing}
\begin{Verbatim}[commandchars=\\\{\}]
\PY{n}{fig}\PY{p}{,} \PY{n}{ax} \PY{o}{=} \PY{n}{subplots}\PY{p}{(}\PY{n}{figsize}\PY{o}{=}\PY{p}{(}\PY{l+m+mi}{8}\PY{p}{,} \PY{l+m+mi}{8}\PY{p}{)}\PY{p}{)}
\PY{n}{ax}\PY{o}{.}\PY{n}{contour}\PY{p}{(}\PY{n}{x}\PY{p}{,} \PY{n}{y}\PY{p}{,} \PY{n}{f}\PY{p}{,} \PY{n}{levels}\PY{o}{=}\PY{l+m+mi}{45}\PY{p}{)}\PY{p}{;}
\end{Verbatim}
\end{tcolorbox}

    \begin{center}
    \adjustimage{max size={0.9\linewidth}{0.9\paperheight}}{artifacts/latex/Ch02-statlearn-lab_files/artifacts/latex/Ch02-statlearn-lab_124_0.png}
    \end{center}
    { \hspace*{\fill} \\}
    
    To fine-tune the output of the \texttt{ax.contour()} function, take a
look at the help file by typing \texttt{ax.contour?}.

The \texttt{ax.imshow()} method is similar to \texttt{ax.contour()},
except that it produces a color-coded plot whose colors depend on the
\texttt{z} value. This is known as a \emph{heatmap}, and is sometimes
used to plot temperature in weather forecasts.

    \begin{tcolorbox}[breakable, size=fbox, boxrule=1pt, pad at break*=1mm,colback=cellbackground, colframe=cellborder]
\prompt{In}{incolor}{52}{\boxspacing}
\begin{Verbatim}[commandchars=\\\{\}]
\PY{n}{fig}\PY{p}{,} \PY{n}{ax} \PY{o}{=} \PY{n}{subplots}\PY{p}{(}\PY{n}{figsize}\PY{o}{=}\PY{p}{(}\PY{l+m+mi}{8}\PY{p}{,} \PY{l+m+mi}{8}\PY{p}{)}\PY{p}{)}
\PY{n}{ax}\PY{o}{.}\PY{n}{imshow}\PY{p}{(}\PY{n}{f}\PY{p}{)}\PY{p}{;}
\end{Verbatim}
\end{tcolorbox}

    \begin{center}
    \adjustimage{max size={0.9\linewidth}{0.9\paperheight}}{artifacts/latex/Ch02-statlearn-lab_files/artifacts/latex/Ch02-statlearn-lab_126_0.png}
    \end{center}
    { \hspace*{\fill} \\}
    
    \hypertarget{sequences-and-slice-notation}{%
\subsection{Sequences and Slice
Notation}\label{sequences-and-slice-notation}}

    As seen above, the function \texttt{np.linspace()} can be used to create
a sequence of numbers.

    \begin{tcolorbox}[breakable, size=fbox, boxrule=1pt, pad at break*=1mm,colback=cellbackground, colframe=cellborder]
\prompt{In}{incolor}{53}{\boxspacing}
\begin{Verbatim}[commandchars=\\\{\}]
\PY{n}{seq1} \PY{o}{=} \PY{n}{np}\PY{o}{.}\PY{n}{linspace}\PY{p}{(}\PY{l+m+mi}{0}\PY{p}{,} \PY{l+m+mi}{10}\PY{p}{,} \PY{l+m+mi}{11}\PY{p}{)}
\PY{n}{seq1}
\end{Verbatim}
\end{tcolorbox}

            \begin{tcolorbox}[breakable, size=fbox, boxrule=.5pt, pad at break*=1mm, opacityfill=0]
\prompt{Out}{outcolor}{53}{\boxspacing}
\begin{Verbatim}[commandchars=\\\{\}]
array([ 0.,  1.,  2.,  3.,  4.,  5.,  6.,  7.,  8.,  9., 10.])
\end{Verbatim}
\end{tcolorbox}
        
    The function \texttt{np.arange()} returns a sequence of numbers spaced
out by \texttt{step}. If \texttt{step} is not specified, then a default
value of \(1\) is used. Let's create a sequence that starts at \(0\) and
ends at \(10\).

    \begin{tcolorbox}[breakable, size=fbox, boxrule=1pt, pad at break*=1mm,colback=cellbackground, colframe=cellborder]
\prompt{In}{incolor}{54}{\boxspacing}
\begin{Verbatim}[commandchars=\\\{\}]
\PY{n}{seq2} \PY{o}{=} \PY{n}{np}\PY{o}{.}\PY{n}{arange}\PY{p}{(}\PY{l+m+mi}{0}\PY{p}{,} \PY{l+m+mi}{10}\PY{p}{)}
\PY{n}{seq2}
\end{Verbatim}
\end{tcolorbox}

            \begin{tcolorbox}[breakable, size=fbox, boxrule=.5pt, pad at break*=1mm, opacityfill=0]
\prompt{Out}{outcolor}{54}{\boxspacing}
\begin{Verbatim}[commandchars=\\\{\}]
array([0, 1, 2, 3, 4, 5, 6, 7, 8, 9])
\end{Verbatim}
\end{tcolorbox}
        
    Why isn't \(10\) output above? This has to do with \emph{slice} notation
in \texttt{Python}. Slice notation\\
is used to index sequences such as lists, tuples and arrays. Suppose we
want to retrieve the fourth through sixth (inclusive) entries of a
string. We obtain a slice of the string using the indexing notation
\texttt{{[}3:6{]}}.

    \begin{tcolorbox}[breakable, size=fbox, boxrule=1pt, pad at break*=1mm,colback=cellbackground, colframe=cellborder]
\prompt{In}{incolor}{55}{\boxspacing}
\begin{Verbatim}[commandchars=\\\{\}]
\PY{l+s+s2}{\PYZdq{}}\PY{l+s+s2}{hello world}\PY{l+s+s2}{\PYZdq{}}\PY{p}{[}\PY{l+m+mi}{3}\PY{p}{:}\PY{l+m+mi}{6}\PY{p}{]}
\end{Verbatim}
\end{tcolorbox}

            \begin{tcolorbox}[breakable, size=fbox, boxrule=.5pt, pad at break*=1mm, opacityfill=0]
\prompt{Out}{outcolor}{55}{\boxspacing}
\begin{Verbatim}[commandchars=\\\{\}]
'lo '
\end{Verbatim}
\end{tcolorbox}
        
    In the code block above, the notation \texttt{3:6} is shorthand for
\texttt{slice(3,6)} when used inside \texttt{{[}{]}}.

    \begin{tcolorbox}[breakable, size=fbox, boxrule=1pt, pad at break*=1mm,colback=cellbackground, colframe=cellborder]
\prompt{In}{incolor}{56}{\boxspacing}
\begin{Verbatim}[commandchars=\\\{\}]
\PY{l+s+s2}{\PYZdq{}}\PY{l+s+s2}{hello world}\PY{l+s+s2}{\PYZdq{}}\PY{p}{[}\PY{n+nb}{slice}\PY{p}{(}\PY{l+m+mi}{3}\PY{p}{,}\PY{l+m+mi}{6}\PY{p}{)}\PY{p}{]}
\end{Verbatim}
\end{tcolorbox}

            \begin{tcolorbox}[breakable, size=fbox, boxrule=.5pt, pad at break*=1mm, opacityfill=0]
\prompt{Out}{outcolor}{56}{\boxspacing}
\begin{Verbatim}[commandchars=\\\{\}]
'lo '
\end{Verbatim}
\end{tcolorbox}
        
    You might have expected \texttt{slice(3,6)} to output the fourth through
seventh characters in the text string (recalling that \texttt{Python}
begins its indexing at zero), but instead it output the fourth through
sixth. This also explains why the earlier \texttt{np.arange(0,\ 10)}
command output only the integers from \(0\) to \(9\). See the
documentation \texttt{slice?} for useful options in creating slices.

    \hypertarget{indexing-data}{%
\subsection{Indexing Data}\label{indexing-data}}

To begin, we create a two-dimensional \texttt{numpy} array.

    \begin{tcolorbox}[breakable, size=fbox, boxrule=1pt, pad at break*=1mm,colback=cellbackground, colframe=cellborder]
\prompt{In}{incolor}{57}{\boxspacing}
\begin{Verbatim}[commandchars=\\\{\}]
\PY{n}{A} \PY{o}{=} \PY{n}{np}\PY{o}{.}\PY{n}{array}\PY{p}{(}\PY{n}{np}\PY{o}{.}\PY{n}{arange}\PY{p}{(}\PY{l+m+mi}{16}\PY{p}{)}\PY{p}{)}\PY{o}{.}\PY{n}{reshape}\PY{p}{(}\PY{p}{(}\PY{l+m+mi}{4}\PY{p}{,} \PY{l+m+mi}{4}\PY{p}{)}\PY{p}{)}
\PY{n}{A}
\end{Verbatim}
\end{tcolorbox}

            \begin{tcolorbox}[breakable, size=fbox, boxrule=.5pt, pad at break*=1mm, opacityfill=0]
\prompt{Out}{outcolor}{57}{\boxspacing}
\begin{Verbatim}[commandchars=\\\{\}]
array([[ 0,  1,  2,  3],
       [ 4,  5,  6,  7],
       [ 8,  9, 10, 11],
       [12, 13, 14, 15]])
\end{Verbatim}
\end{tcolorbox}
        
    Typing \texttt{A{[}1,2{]}} retrieves the element corresponding to the
second row and third column. (As usual, \texttt{Python} indexes from
\(0.\))

    \begin{tcolorbox}[breakable, size=fbox, boxrule=1pt, pad at break*=1mm,colback=cellbackground, colframe=cellborder]
\prompt{In}{incolor}{58}{\boxspacing}
\begin{Verbatim}[commandchars=\\\{\}]
\PY{n}{A}\PY{p}{[}\PY{l+m+mi}{1}\PY{p}{,}\PY{l+m+mi}{2}\PY{p}{]}
\end{Verbatim}
\end{tcolorbox}

            \begin{tcolorbox}[breakable, size=fbox, boxrule=.5pt, pad at break*=1mm, opacityfill=0]
\prompt{Out}{outcolor}{58}{\boxspacing}
\begin{Verbatim}[commandchars=\\\{\}]
6
\end{Verbatim}
\end{tcolorbox}
        
    The first number after the open-bracket symbol \texttt{{[}} refers to
the row, and the second number refers to the column.

\hypertarget{indexing-rows-columns-and-submatrices}{%
\subsubsection{Indexing Rows, Columns, and
Submatrices}\label{indexing-rows-columns-and-submatrices}}

To select multiple rows at a time, we can pass in a list specifying our
selection. For instance, \texttt{{[}1,3{]}} will retrieve the second and
fourth rows:

    \begin{tcolorbox}[breakable, size=fbox, boxrule=1pt, pad at break*=1mm,colback=cellbackground, colframe=cellborder]
\prompt{In}{incolor}{59}{\boxspacing}
\begin{Verbatim}[commandchars=\\\{\}]
\PY{n}{A}\PY{p}{[}\PY{p}{[}\PY{l+m+mi}{1}\PY{p}{,}\PY{l+m+mi}{3}\PY{p}{]}\PY{p}{]}
\end{Verbatim}
\end{tcolorbox}

            \begin{tcolorbox}[breakable, size=fbox, boxrule=.5pt, pad at break*=1mm, opacityfill=0]
\prompt{Out}{outcolor}{59}{\boxspacing}
\begin{Verbatim}[commandchars=\\\{\}]
array([[ 4,  5,  6,  7],
       [12, 13, 14, 15]])
\end{Verbatim}
\end{tcolorbox}
        
    To select the first and third columns, we pass in \texttt{{[}0,2{]}} as
the second argument in the square brackets. In this case we need to
supply the first argument \texttt{:} which selects all rows.

    \begin{tcolorbox}[breakable, size=fbox, boxrule=1pt, pad at break*=1mm,colback=cellbackground, colframe=cellborder]
\prompt{In}{incolor}{60}{\boxspacing}
\begin{Verbatim}[commandchars=\\\{\}]
\PY{n}{A}\PY{p}{[}\PY{p}{:}\PY{p}{,}\PY{p}{[}\PY{l+m+mi}{0}\PY{p}{,}\PY{l+m+mi}{2}\PY{p}{]}\PY{p}{]}
\end{Verbatim}
\end{tcolorbox}

            \begin{tcolorbox}[breakable, size=fbox, boxrule=.5pt, pad at break*=1mm, opacityfill=0]
\prompt{Out}{outcolor}{60}{\boxspacing}
\begin{Verbatim}[commandchars=\\\{\}]
array([[ 0,  2],
       [ 4,  6],
       [ 8, 10],
       [12, 14]])
\end{Verbatim}
\end{tcolorbox}
        
    Now, suppose that we want to select the submatrix made up of the second
and fourth rows as well as the first and third columns. This is where
indexing gets slightly tricky. It is natural to try to use lists to
retrieve the rows and columns:

    \begin{tcolorbox}[breakable, size=fbox, boxrule=1pt, pad at break*=1mm,colback=cellbackground, colframe=cellborder]
\prompt{In}{incolor}{61}{\boxspacing}
\begin{Verbatim}[commandchars=\\\{\}]
\PY{n}{A}\PY{p}{[}\PY{p}{[}\PY{l+m+mi}{1}\PY{p}{,}\PY{l+m+mi}{3}\PY{p}{]}\PY{p}{,}\PY{p}{[}\PY{l+m+mi}{0}\PY{p}{,}\PY{l+m+mi}{2}\PY{p}{]}\PY{p}{]}
\end{Verbatim}
\end{tcolorbox}

            \begin{tcolorbox}[breakable, size=fbox, boxrule=.5pt, pad at break*=1mm, opacityfill=0]
\prompt{Out}{outcolor}{61}{\boxspacing}
\begin{Verbatim}[commandchars=\\\{\}]
array([ 4, 14])
\end{Verbatim}
\end{tcolorbox}
        
    Oops --- what happened? We got a one-dimensional array of length two
identical to

    \begin{tcolorbox}[breakable, size=fbox, boxrule=1pt, pad at break*=1mm,colback=cellbackground, colframe=cellborder]
\prompt{In}{incolor}{62}{\boxspacing}
\begin{Verbatim}[commandchars=\\\{\}]
\PY{n}{np}\PY{o}{.}\PY{n}{array}\PY{p}{(}\PY{p}{[}\PY{n}{A}\PY{p}{[}\PY{l+m+mi}{1}\PY{p}{,}\PY{l+m+mi}{0}\PY{p}{]}\PY{p}{,}\PY{n}{A}\PY{p}{[}\PY{l+m+mi}{3}\PY{p}{,}\PY{l+m+mi}{2}\PY{p}{]}\PY{p}{]}\PY{p}{)}
\end{Verbatim}
\end{tcolorbox}

            \begin{tcolorbox}[breakable, size=fbox, boxrule=.5pt, pad at break*=1mm, opacityfill=0]
\prompt{Out}{outcolor}{62}{\boxspacing}
\begin{Verbatim}[commandchars=\\\{\}]
array([ 4, 14])
\end{Verbatim}
\end{tcolorbox}
        
    Similarly, the following code fails to extract the submatrix comprised
of the second and fourth rows and the first, third, and fourth columns:

    \begin{tcolorbox}[breakable, size=fbox, boxrule=1pt, pad at break*=1mm,colback=cellbackground, colframe=cellborder]
\prompt{In}{incolor}{63}{\boxspacing}
\begin{Verbatim}[commandchars=\\\{\}]
\PY{n}{A}\PY{p}{[}\PY{p}{[}\PY{l+m+mi}{1}\PY{p}{,}\PY{l+m+mi}{3}\PY{p}{]}\PY{p}{,}\PY{p}{[}\PY{l+m+mi}{0}\PY{p}{,}\PY{l+m+mi}{2}\PY{p}{,}\PY{l+m+mi}{3}\PY{p}{]}\PY{p}{]}
\end{Verbatim}
\end{tcolorbox}

    \begin{Verbatim}[commandchars=\\\{\}, frame=single, framerule=2mm, rulecolor=\color{outerrorbackground}]
\textcolor{ansi-red}{---------------------------------------------------------------------------}
\textcolor{ansi-red}{IndexError}                                Traceback (most recent call last)
\textcolor{ansi-cyan}{Cell}\textcolor{ansi-cyan}{ }\textcolor{ansi-green}{In[63]}\textcolor{ansi-green}{, line 1}
\textcolor{ansi-green}{----> }\textcolor{ansi-green}{1} \setlength{\fboxsep}{0pt}\colorbox{ansi-yellow}{A\strut}\setlength{\fboxsep}{0pt}\colorbox{ansi-yellow}{[\strut}\setlength{\fboxsep}{0pt}\colorbox{ansi-yellow}{[\strut}\textcolor{ansi-green}{\setlength{\fboxsep}{0pt}\colorbox{ansi-yellow}{1\strut}}\setlength{\fboxsep}{0pt}\colorbox{ansi-yellow}{,\strut}\textcolor{ansi-green}{\setlength{\fboxsep}{0pt}\colorbox{ansi-yellow}{3\strut}}\setlength{\fboxsep}{0pt}\colorbox{ansi-yellow}{]\strut}\setlength{\fboxsep}{0pt}\colorbox{ansi-yellow}{,\strut}\setlength{\fboxsep}{0pt}\colorbox{ansi-yellow}{[\strut}\textcolor{ansi-green}{\setlength{\fboxsep}{0pt}\colorbox{ansi-yellow}{0\strut}}\setlength{\fboxsep}{0pt}\colorbox{ansi-yellow}{,\strut}\textcolor{ansi-green}{\setlength{\fboxsep}{0pt}\colorbox{ansi-yellow}{2\strut}}\setlength{\fboxsep}{0pt}\colorbox{ansi-yellow}{,\strut}\textcolor{ansi-green}{\setlength{\fboxsep}{0pt}\colorbox{ansi-yellow}{3\strut}}\setlength{\fboxsep}{0pt}\colorbox{ansi-yellow}{]\strut}\setlength{\fboxsep}{0pt}\colorbox{ansi-yellow}{]\strut}

\textcolor{ansi-red}{IndexError}: shape mismatch: indexing arrays could not be broadcast together with shapes (2,) (3,) 
    \end{Verbatim}

    We can see what has gone wrong here. When supplied with two indexing
lists, the \texttt{numpy} interpretation is that these provide pairs of
\(i,j\) indices for a series of entries. That is why the pair of lists
must have the same length. However, that was not our intent, since we
are looking for a submatrix.

One easy way to do this is as follows. We first create a submatrix by
subsetting the rows of \texttt{A}, and then on the fly we make a further
submatrix by subsetting its columns.

    \begin{tcolorbox}[breakable, size=fbox, boxrule=1pt, pad at break*=1mm,colback=cellbackground, colframe=cellborder]
\prompt{In}{incolor}{64}{\boxspacing}
\begin{Verbatim}[commandchars=\\\{\}]
\PY{n}{A}\PY{p}{[}\PY{p}{[}\PY{l+m+mi}{1}\PY{p}{,}\PY{l+m+mi}{3}\PY{p}{]}\PY{p}{]}\PY{p}{[}\PY{p}{:}\PY{p}{,}\PY{p}{[}\PY{l+m+mi}{0}\PY{p}{,}\PY{l+m+mi}{2}\PY{p}{]}\PY{p}{]}
\end{Verbatim}
\end{tcolorbox}

            \begin{tcolorbox}[breakable, size=fbox, boxrule=.5pt, pad at break*=1mm, opacityfill=0]
\prompt{Out}{outcolor}{64}{\boxspacing}
\begin{Verbatim}[commandchars=\\\{\}]
array([[ 4,  6],
       [12, 14]])
\end{Verbatim}
\end{tcolorbox}
        
    

    There are more efficient ways of achieving the same result.

The \emph{convenience function} \texttt{np.ix\_()} allows us to extract
a submatrix using lists, by creating an intermediate \emph{mesh} object.

    \begin{tcolorbox}[breakable, size=fbox, boxrule=1pt, pad at break*=1mm,colback=cellbackground, colframe=cellborder]
\prompt{In}{incolor}{65}{\boxspacing}
\begin{Verbatim}[commandchars=\\\{\}]
\PY{n}{idx} \PY{o}{=} \PY{n}{np}\PY{o}{.}\PY{n}{ix\PYZus{}}\PY{p}{(}\PY{p}{[}\PY{l+m+mi}{1}\PY{p}{,}\PY{l+m+mi}{3}\PY{p}{]}\PY{p}{,}\PY{p}{[}\PY{l+m+mi}{0}\PY{p}{,}\PY{l+m+mi}{2}\PY{p}{,}\PY{l+m+mi}{3}\PY{p}{]}\PY{p}{)}
\PY{n}{A}\PY{p}{[}\PY{n}{idx}\PY{p}{]}
\end{Verbatim}
\end{tcolorbox}

            \begin{tcolorbox}[breakable, size=fbox, boxrule=.5pt, pad at break*=1mm, opacityfill=0]
\prompt{Out}{outcolor}{65}{\boxspacing}
\begin{Verbatim}[commandchars=\\\{\}]
array([[ 4,  6,  7],
       [12, 14, 15]])
\end{Verbatim}
\end{tcolorbox}
        
    Alternatively, we can subset matrices efficiently using slices.

The slice \texttt{1:4:2} captures the second and fourth items of a
sequence, while the slice \texttt{0:3:2} captures the first and third
items (the third element in a slice sequence is the step size).

    \begin{tcolorbox}[breakable, size=fbox, boxrule=1pt, pad at break*=1mm,colback=cellbackground, colframe=cellborder]
\prompt{In}{incolor}{66}{\boxspacing}
\begin{Verbatim}[commandchars=\\\{\}]
\PY{n}{A}\PY{p}{[}\PY{l+m+mi}{1}\PY{p}{:}\PY{l+m+mi}{4}\PY{p}{:}\PY{l+m+mi}{2}\PY{p}{,}\PY{l+m+mi}{0}\PY{p}{:}\PY{l+m+mi}{3}\PY{p}{:}\PY{l+m+mi}{2}\PY{p}{]}
\end{Verbatim}
\end{tcolorbox}

            \begin{tcolorbox}[breakable, size=fbox, boxrule=.5pt, pad at break*=1mm, opacityfill=0]
\prompt{Out}{outcolor}{66}{\boxspacing}
\begin{Verbatim}[commandchars=\\\{\}]
array([[ 4,  6],
       [12, 14]])
\end{Verbatim}
\end{tcolorbox}
        
    

    Why are we able to retrieve a submatrix directly using slices but not
using lists? It's because they are different \texttt{Python} types, and
are treated differently by \texttt{numpy}. Slices can be used to extract
objects from arbitrary sequences, such as strings, lists, and tuples,
while the use of lists for indexing is more limited.

    \hypertarget{boolean-indexing}{%
\subsubsection{Boolean Indexing}\label{boolean-indexing}}

In \texttt{numpy}, a \emph{Boolean} is a type that equals either
\texttt{True} or \texttt{False} (also represented as \(1\) and \(0\),
respectively). The next line creates a vector of \(0\)'s, represented as
Booleans, of length equal to the first dimension of \texttt{A}.

    \begin{tcolorbox}[breakable, size=fbox, boxrule=1pt, pad at break*=1mm,colback=cellbackground, colframe=cellborder]
\prompt{In}{incolor}{67}{\boxspacing}
\begin{Verbatim}[commandchars=\\\{\}]
\PY{n}{keep\PYZus{}rows} \PY{o}{=} \PY{n}{np}\PY{o}{.}\PY{n}{zeros}\PY{p}{(}\PY{n}{A}\PY{o}{.}\PY{n}{shape}\PY{p}{[}\PY{l+m+mi}{0}\PY{p}{]}\PY{p}{,} \PY{n+nb}{bool}\PY{p}{)}
\PY{n}{keep\PYZus{}rows}
\end{Verbatim}
\end{tcolorbox}

            \begin{tcolorbox}[breakable, size=fbox, boxrule=.5pt, pad at break*=1mm, opacityfill=0]
\prompt{Out}{outcolor}{67}{\boxspacing}
\begin{Verbatim}[commandchars=\\\{\}]
array([False, False, False, False])
\end{Verbatim}
\end{tcolorbox}
        
    We now set two of the elements to \texttt{True}.

    \begin{tcolorbox}[breakable, size=fbox, boxrule=1pt, pad at break*=1mm,colback=cellbackground, colframe=cellborder]
\prompt{In}{incolor}{68}{\boxspacing}
\begin{Verbatim}[commandchars=\\\{\}]
\PY{n}{keep\PYZus{}rows}\PY{p}{[}\PY{p}{[}\PY{l+m+mi}{1}\PY{p}{,}\PY{l+m+mi}{3}\PY{p}{]}\PY{p}{]} \PY{o}{=} \PY{k+kc}{True}
\PY{n}{keep\PYZus{}rows}
\end{Verbatim}
\end{tcolorbox}

            \begin{tcolorbox}[breakable, size=fbox, boxrule=.5pt, pad at break*=1mm, opacityfill=0]
\prompt{Out}{outcolor}{68}{\boxspacing}
\begin{Verbatim}[commandchars=\\\{\}]
array([False,  True, False,  True])
\end{Verbatim}
\end{tcolorbox}
        
    Note that the elements of \texttt{keep\_rows}, when viewed as integers,
are the same as the values of \texttt{np.array({[}0,1,0,1{]})}. Below,
we use \texttt{==} to verify their equality. When applied to two arrays,
the \texttt{==} operation is applied elementwise.

    \begin{tcolorbox}[breakable, size=fbox, boxrule=1pt, pad at break*=1mm,colback=cellbackground, colframe=cellborder]
\prompt{In}{incolor}{69}{\boxspacing}
\begin{Verbatim}[commandchars=\\\{\}]
\PY{n}{np}\PY{o}{.}\PY{n}{all}\PY{p}{(}\PY{n}{keep\PYZus{}rows} \PY{o}{==} \PY{n}{np}\PY{o}{.}\PY{n}{array}\PY{p}{(}\PY{p}{[}\PY{l+m+mi}{0}\PY{p}{,}\PY{l+m+mi}{1}\PY{p}{,}\PY{l+m+mi}{0}\PY{p}{,}\PY{l+m+mi}{1}\PY{p}{]}\PY{p}{)}\PY{p}{)}
\end{Verbatim}
\end{tcolorbox}

            \begin{tcolorbox}[breakable, size=fbox, boxrule=.5pt, pad at break*=1mm, opacityfill=0]
\prompt{Out}{outcolor}{69}{\boxspacing}
\begin{Verbatim}[commandchars=\\\{\}]
True
\end{Verbatim}
\end{tcolorbox}
        
    (Here, the function \texttt{np.all()} has checked whether all entries of
an array are \texttt{True}. A similar function, \texttt{np.any()}, can
be used to check whether any entries of an array are \texttt{True}.)

    However, even though \texttt{np.array({[}0,1,0,1{]})} and
\texttt{keep\_rows} are equal according to \texttt{==}, they index
different sets of rows! The former retrieves the first, second, first,
and second rows of \texttt{A}.

    \begin{tcolorbox}[breakable, size=fbox, boxrule=1pt, pad at break*=1mm,colback=cellbackground, colframe=cellborder]
\prompt{In}{incolor}{70}{\boxspacing}
\begin{Verbatim}[commandchars=\\\{\}]
\PY{n}{A}\PY{p}{[}\PY{n}{np}\PY{o}{.}\PY{n}{array}\PY{p}{(}\PY{p}{[}\PY{l+m+mi}{0}\PY{p}{,}\PY{l+m+mi}{1}\PY{p}{,}\PY{l+m+mi}{0}\PY{p}{,}\PY{l+m+mi}{1}\PY{p}{]}\PY{p}{)}\PY{p}{]}
\end{Verbatim}
\end{tcolorbox}

            \begin{tcolorbox}[breakable, size=fbox, boxrule=.5pt, pad at break*=1mm, opacityfill=0]
\prompt{Out}{outcolor}{70}{\boxspacing}
\begin{Verbatim}[commandchars=\\\{\}]
array([[0, 1, 2, 3],
       [4, 5, 6, 7],
       [0, 1, 2, 3],
       [4, 5, 6, 7]])
\end{Verbatim}
\end{tcolorbox}
        
    By contrast, \texttt{keep\_rows} retrieves only the second and fourth
rows of \texttt{A} --- i.e.~the rows for which the Boolean equals
\texttt{True}.

    \begin{tcolorbox}[breakable, size=fbox, boxrule=1pt, pad at break*=1mm,colback=cellbackground, colframe=cellborder]
\prompt{In}{incolor}{71}{\boxspacing}
\begin{Verbatim}[commandchars=\\\{\}]
\PY{n}{A}\PY{p}{[}\PY{n}{keep\PYZus{}rows}\PY{p}{]}
\end{Verbatim}
\end{tcolorbox}

            \begin{tcolorbox}[breakable, size=fbox, boxrule=.5pt, pad at break*=1mm, opacityfill=0]
\prompt{Out}{outcolor}{71}{\boxspacing}
\begin{Verbatim}[commandchars=\\\{\}]
array([[ 4,  5,  6,  7],
       [12, 13, 14, 15]])
\end{Verbatim}
\end{tcolorbox}
        
    This example shows that Booleans and integers are treated differently by
\texttt{numpy}.

    We again make use of the \texttt{np.ix\_()} function to create a mesh
containing the second and fourth rows, and the first, third, and fourth
columns. This time, we apply the function to Booleans, rather than
lists.

    \begin{tcolorbox}[breakable, size=fbox, boxrule=1pt, pad at break*=1mm,colback=cellbackground, colframe=cellborder]
\prompt{In}{incolor}{72}{\boxspacing}
\begin{Verbatim}[commandchars=\\\{\}]
\PY{n}{keep\PYZus{}cols} \PY{o}{=} \PY{n}{np}\PY{o}{.}\PY{n}{zeros}\PY{p}{(}\PY{n}{A}\PY{o}{.}\PY{n}{shape}\PY{p}{[}\PY{l+m+mi}{1}\PY{p}{]}\PY{p}{,} \PY{n+nb}{bool}\PY{p}{)}
\PY{n}{keep\PYZus{}cols}\PY{p}{[}\PY{p}{[}\PY{l+m+mi}{0}\PY{p}{,} \PY{l+m+mi}{2}\PY{p}{,} \PY{l+m+mi}{3}\PY{p}{]}\PY{p}{]} \PY{o}{=} \PY{k+kc}{True}
\PY{n}{idx\PYZus{}bool} \PY{o}{=} \PY{n}{np}\PY{o}{.}\PY{n}{ix\PYZus{}}\PY{p}{(}\PY{n}{keep\PYZus{}rows}\PY{p}{,} \PY{n}{keep\PYZus{}cols}\PY{p}{)}
\PY{n}{A}\PY{p}{[}\PY{n}{idx\PYZus{}bool}\PY{p}{]}
\end{Verbatim}
\end{tcolorbox}

            \begin{tcolorbox}[breakable, size=fbox, boxrule=.5pt, pad at break*=1mm, opacityfill=0]
\prompt{Out}{outcolor}{72}{\boxspacing}
\begin{Verbatim}[commandchars=\\\{\}]
array([[ 4,  6,  7],
       [12, 14, 15]])
\end{Verbatim}
\end{tcolorbox}
        
    We can also mix a list with an array of Booleans in the arguments to
\texttt{np.ix\_()}:

    \begin{tcolorbox}[breakable, size=fbox, boxrule=1pt, pad at break*=1mm,colback=cellbackground, colframe=cellborder]
\prompt{In}{incolor}{73}{\boxspacing}
\begin{Verbatim}[commandchars=\\\{\}]
\PY{n}{idx\PYZus{}mixed} \PY{o}{=} \PY{n}{np}\PY{o}{.}\PY{n}{ix\PYZus{}}\PY{p}{(}\PY{p}{[}\PY{l+m+mi}{1}\PY{p}{,}\PY{l+m+mi}{3}\PY{p}{]}\PY{p}{,} \PY{n}{keep\PYZus{}cols}\PY{p}{)}
\PY{n}{A}\PY{p}{[}\PY{n}{idx\PYZus{}mixed}\PY{p}{]}
\end{Verbatim}
\end{tcolorbox}

            \begin{tcolorbox}[breakable, size=fbox, boxrule=.5pt, pad at break*=1mm, opacityfill=0]
\prompt{Out}{outcolor}{73}{\boxspacing}
\begin{Verbatim}[commandchars=\\\{\}]
array([[ 4,  6,  7],
       [12, 14, 15]])
\end{Verbatim}
\end{tcolorbox}
        
    

    For more details on indexing in \texttt{numpy}, readers are referred to
the \texttt{numpy} tutorial mentioned earlier.

    \hypertarget{loading-data}{%
\subsection{Loading Data}\label{loading-data}}

Data sets often contain different types of data, and may have names
associated with the rows or columns. For these reasons, they typically
are best accommodated using a \emph{data frame}. We can think of a data
frame as a sequence of arrays of identical length; these are the
columns. Entries in the different arrays can be combined to form a row.
The \texttt{pandas} library can be used to create and work with data
frame objects.

    \hypertarget{reading-in-a-data-set}{%
\subsubsection{Reading in a Data Set}\label{reading-in-a-data-set}}

The first step of most analyses involves importing a data set into
\texttt{Python}.\\
Before attempting to load a data set, we must make sure that
\texttt{Python} knows where to find the file containing it. If the file
is in the same location as this notebook file, then we are all set.
Otherwise, the command \texttt{os.chdir()} can be used to \emph{change
directory}. (You will need to call \texttt{import\ os} before calling
\texttt{os.chdir()}.)

    We will begin by reading in \texttt{Auto.csv}, available on the book
website. This is a comma-separated file, and can be read in using
\texttt{pd.read\_csv()}:

    \begin{tcolorbox}[breakable, size=fbox, boxrule=1pt, pad at break*=1mm,colback=cellbackground, colframe=cellborder]
\prompt{In}{incolor}{74}{\boxspacing}
\begin{Verbatim}[commandchars=\\\{\}]
\PY{k+kn}{import}\PY{+w}{ }\PY{n+nn}{pandas}\PY{+w}{ }\PY{k}{as}\PY{+w}{ }\PY{n+nn}{pd}
\PY{n}{Auto} \PY{o}{=} \PY{n}{pd}\PY{o}{.}\PY{n}{read\PYZus{}csv}\PY{p}{(}\PY{l+s+s1}{\PYZsq{}}\PY{l+s+s1}{Auto.csv}\PY{l+s+s1}{\PYZsq{}}\PY{p}{)}
\PY{n}{Auto}
\end{Verbatim}
\end{tcolorbox}

            \begin{tcolorbox}[breakable, size=fbox, boxrule=.5pt, pad at break*=1mm, opacityfill=0]
\prompt{Out}{outcolor}{74}{\boxspacing}
\begin{Verbatim}[commandchars=\\\{\}]
      mpg  cylinders  displacement  horsepower  weight  acceleration  year  \textbackslash{}
0    18.0          8         307.0         130    3504          12.0    70
1    15.0          8         350.0         165    3693          11.5    70
2    18.0          8         318.0         150    3436          11.0    70
3    16.0          8         304.0         150    3433          12.0    70
4    17.0          8         302.0         140    3449          10.5    70
..    {\ldots}        {\ldots}           {\ldots}         {\ldots}     {\ldots}           {\ldots}   {\ldots}
387  27.0          4         140.0          86    2790          15.6    82
388  44.0          4          97.0          52    2130          24.6    82
389  32.0          4         135.0          84    2295          11.6    82
390  28.0          4         120.0          79    2625          18.6    82
391  31.0          4         119.0          82    2720          19.4    82

     origin                       name
0         1  chevrolet chevelle malibu
1         1          buick skylark 320
2         1         plymouth satellite
3         1              amc rebel sst
4         1                ford torino
..      {\ldots}                        {\ldots}
387       1            ford mustang gl
388       2                  vw pickup
389       1              dodge rampage
390       1                ford ranger
391       1                 chevy s-10

[392 rows x 9 columns]
\end{Verbatim}
\end{tcolorbox}
        
    The book website also has a whitespace-delimited version of this data,
called \texttt{Auto.data}. This can be read in as follows:

    \begin{tcolorbox}[breakable, size=fbox, boxrule=1pt, pad at break*=1mm,colback=cellbackground, colframe=cellborder]
\prompt{In}{incolor}{75}{\boxspacing}
\begin{Verbatim}[commandchars=\\\{\}]
\PY{n}{Auto} \PY{o}{=} \PY{n}{pd}\PY{o}{.}\PY{n}{read\PYZus{}csv}\PY{p}{(}\PY{l+s+s1}{\PYZsq{}}\PY{l+s+s1}{Auto.data}\PY{l+s+s1}{\PYZsq{}}\PY{p}{,} \PY{n}{delim\PYZus{}whitespace}\PY{o}{=}\PY{k+kc}{True}\PY{p}{)}
\end{Verbatim}
\end{tcolorbox}

    \begin{Verbatim}[commandchars=\\\{\}]
/var/folders/\_3/k\_9k5d5n5zz57w7qfll9rzs40000gn/T/ipykernel\_64281/1967952672.py:1
: FutureWarning: The 'delim\_whitespace' keyword in pd.read\_csv is deprecated and
will be removed in a future version. Use ``sep='\textbackslash{}s+'`` instead
  Auto = pd.read\_csv('Auto.data', delim\_whitespace=True)
    \end{Verbatim}

    Both \texttt{Auto.csv} and \texttt{Auto.data} are simply text files.
Before loading data into \texttt{Python}, it is a good idea to view it
using a text editor or other software, such as Microsoft Excel.

    We now take a look at the column of \texttt{Auto} corresponding to the
variable \texttt{horsepower}:

    \begin{tcolorbox}[breakable, size=fbox, boxrule=1pt, pad at break*=1mm,colback=cellbackground, colframe=cellborder]
\prompt{In}{incolor}{76}{\boxspacing}
\begin{Verbatim}[commandchars=\\\{\}]
\PY{n}{Auto}\PY{p}{[}\PY{l+s+s1}{\PYZsq{}}\PY{l+s+s1}{horsepower}\PY{l+s+s1}{\PYZsq{}}\PY{p}{]}
\end{Verbatim}
\end{tcolorbox}

            \begin{tcolorbox}[breakable, size=fbox, boxrule=.5pt, pad at break*=1mm, opacityfill=0]
\prompt{Out}{outcolor}{76}{\boxspacing}
\begin{Verbatim}[commandchars=\\\{\}]
0      130.0
1      165.0
2      150.0
3      150.0
4      140.0
       {\ldots}
392    86.00
393    52.00
394    84.00
395    79.00
396    82.00
Name: horsepower, Length: 397, dtype: object
\end{Verbatim}
\end{tcolorbox}
        
    We see that the \texttt{dtype} of this column is \texttt{object}. It
turns out that all values of the \texttt{horsepower} column were
interpreted as strings when reading in the data. We can find out why by
looking at the unique values.

    \begin{tcolorbox}[breakable, size=fbox, boxrule=1pt, pad at break*=1mm,colback=cellbackground, colframe=cellborder]
\prompt{In}{incolor}{77}{\boxspacing}
\begin{Verbatim}[commandchars=\\\{\}]
\PY{n}{np}\PY{o}{.}\PY{n}{unique}\PY{p}{(}\PY{n}{Auto}\PY{p}{[}\PY{l+s+s1}{\PYZsq{}}\PY{l+s+s1}{horsepower}\PY{l+s+s1}{\PYZsq{}}\PY{p}{]}\PY{p}{)}
\end{Verbatim}
\end{tcolorbox}

            \begin{tcolorbox}[breakable, size=fbox, boxrule=.5pt, pad at break*=1mm, opacityfill=0]
\prompt{Out}{outcolor}{77}{\boxspacing}
\begin{Verbatim}[commandchars=\\\{\}]
array(['100.0', '102.0', '103.0', '105.0', '107.0', '108.0', '110.0',
       '112.0', '113.0', '115.0', '116.0', '120.0', '122.0', '125.0',
       '129.0', '130.0', '132.0', '133.0', '135.0', '137.0', '138.0',
       '139.0', '140.0', '142.0', '145.0', '148.0', '149.0', '150.0',
       '152.0', '153.0', '155.0', '158.0', '160.0', '165.0', '167.0',
       '170.0', '175.0', '180.0', '190.0', '193.0', '198.0', '200.0',
       '208.0', '210.0', '215.0', '220.0', '225.0', '230.0', '46.00',
       '48.00', '49.00', '52.00', '53.00', '54.00', '58.00', '60.00',
       '61.00', '62.00', '63.00', '64.00', '65.00', '66.00', '67.00',
       '68.00', '69.00', '70.00', '71.00', '72.00', '74.00', '75.00',
       '76.00', '77.00', '78.00', '79.00', '80.00', '81.00', '82.00',
       '83.00', '84.00', '85.00', '86.00', '87.00', '88.00', '89.00',
       '90.00', '91.00', '92.00', '93.00', '94.00', '95.00', '96.00',
       '97.00', '98.00', '?'], dtype=object)
\end{Verbatim}
\end{tcolorbox}
        
    We see the culprit is the value \texttt{?}, which is being used to
encode missing values.

    To fix the problem, we must provide \texttt{pd.read\_csv()} with an
argument called \texttt{na\_values}. Now, each instance of \texttt{?} in
the file is replaced with the value \texttt{np.nan}, which means
\emph{not a number}:

    \begin{tcolorbox}[breakable, size=fbox, boxrule=1pt, pad at break*=1mm,colback=cellbackground, colframe=cellborder]
\prompt{In}{incolor}{78}{\boxspacing}
\begin{Verbatim}[commandchars=\\\{\}]
\PY{n}{Auto} \PY{o}{=} \PY{n}{pd}\PY{o}{.}\PY{n}{read\PYZus{}csv}\PY{p}{(}\PY{l+s+s1}{\PYZsq{}}\PY{l+s+s1}{Auto.data}\PY{l+s+s1}{\PYZsq{}}\PY{p}{,}
                   \PY{n}{na\PYZus{}values}\PY{o}{=}\PY{p}{[}\PY{l+s+s1}{\PYZsq{}}\PY{l+s+s1}{?}\PY{l+s+s1}{\PYZsq{}}\PY{p}{]}\PY{p}{,}
                   \PY{n}{delim\PYZus{}whitespace}\PY{o}{=}\PY{k+kc}{True}\PY{p}{)}
\PY{n}{Auto}\PY{p}{[}\PY{l+s+s1}{\PYZsq{}}\PY{l+s+s1}{horsepower}\PY{l+s+s1}{\PYZsq{}}\PY{p}{]}\PY{o}{.}\PY{n}{sum}\PY{p}{(}\PY{p}{)}
\end{Verbatim}
\end{tcolorbox}

    \begin{Verbatim}[commandchars=\\\{\}]
/var/folders/\_3/k\_9k5d5n5zz57w7qfll9rzs40000gn/T/ipykernel\_64281/134681464.py:1:
FutureWarning: The 'delim\_whitespace' keyword in pd.read\_csv is deprecated and
will be removed in a future version. Use ``sep='\textbackslash{}s+'`` instead
  Auto = pd.read\_csv('Auto.data',
    \end{Verbatim}

            \begin{tcolorbox}[breakable, size=fbox, boxrule=.5pt, pad at break*=1mm, opacityfill=0]
\prompt{Out}{outcolor}{78}{\boxspacing}
\begin{Verbatim}[commandchars=\\\{\}]
40952.0
\end{Verbatim}
\end{tcolorbox}
        
    The \texttt{Auto.shape} attribute tells us that the data has 397
observations, or rows, and nine variables, or columns.

    \begin{tcolorbox}[breakable, size=fbox, boxrule=1pt, pad at break*=1mm,colback=cellbackground, colframe=cellborder]
\prompt{In}{incolor}{79}{\boxspacing}
\begin{Verbatim}[commandchars=\\\{\}]
\PY{n}{Auto}\PY{o}{.}\PY{n}{shape}
\end{Verbatim}
\end{tcolorbox}

            \begin{tcolorbox}[breakable, size=fbox, boxrule=.5pt, pad at break*=1mm, opacityfill=0]
\prompt{Out}{outcolor}{79}{\boxspacing}
\begin{Verbatim}[commandchars=\\\{\}]
(397, 9)
\end{Verbatim}
\end{tcolorbox}
        
    There are various ways to deal with missing data. In this case, since
only five of the rows contain missing observations, we choose to use the
\texttt{Auto.dropna()} method to simply remove these rows.

    \begin{tcolorbox}[breakable, size=fbox, boxrule=1pt, pad at break*=1mm,colback=cellbackground, colframe=cellborder]
\prompt{In}{incolor}{80}{\boxspacing}
\begin{Verbatim}[commandchars=\\\{\}]
\PY{n}{Auto\PYZus{}new} \PY{o}{=} \PY{n}{Auto}\PY{o}{.}\PY{n}{dropna}\PY{p}{(}\PY{p}{)}
\PY{n}{Auto\PYZus{}new}\PY{o}{.}\PY{n}{shape}
\end{Verbatim}
\end{tcolorbox}

            \begin{tcolorbox}[breakable, size=fbox, boxrule=.5pt, pad at break*=1mm, opacityfill=0]
\prompt{Out}{outcolor}{80}{\boxspacing}
\begin{Verbatim}[commandchars=\\\{\}]
(392, 9)
\end{Verbatim}
\end{tcolorbox}
        
    \hypertarget{basics-of-selecting-rows-and-columns}{%
\subsubsection{Basics of Selecting Rows and
Columns}\label{basics-of-selecting-rows-and-columns}}

We can use \texttt{Auto.columns} to check the variable names.

    \begin{tcolorbox}[breakable, size=fbox, boxrule=1pt, pad at break*=1mm,colback=cellbackground, colframe=cellborder]
\prompt{In}{incolor}{81}{\boxspacing}
\begin{Verbatim}[commandchars=\\\{\}]
\PY{n}{Auto} \PY{o}{=} \PY{n}{Auto\PYZus{}new} \PY{c+c1}{\PYZsh{} overwrite the previous value}
\PY{n}{Auto}\PY{o}{.}\PY{n}{columns}
\end{Verbatim}
\end{tcolorbox}

            \begin{tcolorbox}[breakable, size=fbox, boxrule=.5pt, pad at break*=1mm, opacityfill=0]
\prompt{Out}{outcolor}{81}{\boxspacing}
\begin{Verbatim}[commandchars=\\\{\}]
Index(['mpg', 'cylinders', 'displacement', 'horsepower', 'weight',
       'acceleration', 'year', 'origin', 'name'],
      dtype='object')
\end{Verbatim}
\end{tcolorbox}
        
    Accessing the rows and columns of a data frame is similar, but not
identical, to accessing the rows and columns of an array. Recall that
the first argument to the \texttt{{[}{]}} method is always applied to
the rows of the array.\\
Similarly, passing in a slice to the \texttt{{[}{]}} method creates a
data frame whose \emph{rows} are determined by the slice:

    \begin{tcolorbox}[breakable, size=fbox, boxrule=1pt, pad at break*=1mm,colback=cellbackground, colframe=cellborder]
\prompt{In}{incolor}{82}{\boxspacing}
\begin{Verbatim}[commandchars=\\\{\}]
\PY{n}{Auto}\PY{p}{[}\PY{p}{:}\PY{l+m+mi}{3}\PY{p}{]}
\end{Verbatim}
\end{tcolorbox}

            \begin{tcolorbox}[breakable, size=fbox, boxrule=.5pt, pad at break*=1mm, opacityfill=0]
\prompt{Out}{outcolor}{82}{\boxspacing}
\begin{Verbatim}[commandchars=\\\{\}]
    mpg  cylinders  displacement  horsepower  weight  acceleration  year  \textbackslash{}
0  18.0          8         307.0       130.0  3504.0          12.0    70
1  15.0          8         350.0       165.0  3693.0          11.5    70
2  18.0          8         318.0       150.0  3436.0          11.0    70

   origin                       name
0       1  chevrolet chevelle malibu
1       1          buick skylark 320
2       1         plymouth satellite
\end{Verbatim}
\end{tcolorbox}
        
    Similarly, an array of Booleans can be used to subset the rows:

    \begin{tcolorbox}[breakable, size=fbox, boxrule=1pt, pad at break*=1mm,colback=cellbackground, colframe=cellborder]
\prompt{In}{incolor}{83}{\boxspacing}
\begin{Verbatim}[commandchars=\\\{\}]
\PY{n}{idx\PYZus{}80} \PY{o}{=} \PY{n}{Auto}\PY{p}{[}\PY{l+s+s1}{\PYZsq{}}\PY{l+s+s1}{year}\PY{l+s+s1}{\PYZsq{}}\PY{p}{]} \PY{o}{\PYZgt{}} \PY{l+m+mi}{80}
\PY{n}{Auto}\PY{p}{[}\PY{n}{idx\PYZus{}80}\PY{p}{]}
\end{Verbatim}
\end{tcolorbox}

            \begin{tcolorbox}[breakable, size=fbox, boxrule=.5pt, pad at break*=1mm, opacityfill=0]
\prompt{Out}{outcolor}{83}{\boxspacing}
\begin{Verbatim}[commandchars=\\\{\}]
      mpg  cylinders  displacement  horsepower  weight  acceleration  year  \textbackslash{}
338  27.2          4         135.0        84.0  2490.0          15.7    81
339  26.6          4         151.0        84.0  2635.0          16.4    81
340  25.8          4         156.0        92.0  2620.0          14.4    81
341  23.5          6         173.0       110.0  2725.0          12.6    81
342  30.0          4         135.0        84.0  2385.0          12.9    81
343  39.1          4          79.0        58.0  1755.0          16.9    81
344  39.0          4          86.0        64.0  1875.0          16.4    81
345  35.1          4          81.0        60.0  1760.0          16.1    81
346  32.3          4          97.0        67.0  2065.0          17.8    81
347  37.0          4          85.0        65.0  1975.0          19.4    81
348  37.7          4          89.0        62.0  2050.0          17.3    81
349  34.1          4          91.0        68.0  1985.0          16.0    81
350  34.7          4         105.0        63.0  2215.0          14.9    81
351  34.4          4          98.0        65.0  2045.0          16.2    81
352  29.9          4          98.0        65.0  2380.0          20.7    81
353  33.0          4         105.0        74.0  2190.0          14.2    81
355  33.7          4         107.0        75.0  2210.0          14.4    81
356  32.4          4         108.0        75.0  2350.0          16.8    81
357  32.9          4         119.0       100.0  2615.0          14.8    81
358  31.6          4         120.0        74.0  2635.0          18.3    81
359  28.1          4         141.0        80.0  3230.0          20.4    81
360  30.7          6         145.0        76.0  3160.0          19.6    81
361  25.4          6         168.0       116.0  2900.0          12.6    81
362  24.2          6         146.0       120.0  2930.0          13.8    81
363  22.4          6         231.0       110.0  3415.0          15.8    81
364  26.6          8         350.0       105.0  3725.0          19.0    81
365  20.2          6         200.0        88.0  3060.0          17.1    81
366  17.6          6         225.0        85.0  3465.0          16.6    81
367  28.0          4         112.0        88.0  2605.0          19.6    82
368  27.0          4         112.0        88.0  2640.0          18.6    82
369  34.0          4         112.0        88.0  2395.0          18.0    82
370  31.0          4         112.0        85.0  2575.0          16.2    82
371  29.0          4         135.0        84.0  2525.0          16.0    82
372  27.0          4         151.0        90.0  2735.0          18.0    82
373  24.0          4         140.0        92.0  2865.0          16.4    82
374  36.0          4         105.0        74.0  1980.0          15.3    82
375  37.0          4          91.0        68.0  2025.0          18.2    82
376  31.0          4          91.0        68.0  1970.0          17.6    82
377  38.0          4         105.0        63.0  2125.0          14.7    82
378  36.0          4          98.0        70.0  2125.0          17.3    82
379  36.0          4         120.0        88.0  2160.0          14.5    82
380  36.0          4         107.0        75.0  2205.0          14.5    82
381  34.0          4         108.0        70.0  2245.0          16.9    82
382  38.0          4          91.0        67.0  1965.0          15.0    82
383  32.0          4          91.0        67.0  1965.0          15.7    82
384  38.0          4          91.0        67.0  1995.0          16.2    82
385  25.0          6         181.0       110.0  2945.0          16.4    82
386  38.0          6         262.0        85.0  3015.0          17.0    82
387  26.0          4         156.0        92.0  2585.0          14.5    82
388  22.0          6         232.0       112.0  2835.0          14.7    82
389  32.0          4         144.0        96.0  2665.0          13.9    82
390  36.0          4         135.0        84.0  2370.0          13.0    82
391  27.0          4         151.0        90.0  2950.0          17.3    82
392  27.0          4         140.0        86.0  2790.0          15.6    82
393  44.0          4          97.0        52.0  2130.0          24.6    82
394  32.0          4         135.0        84.0  2295.0          11.6    82
395  28.0          4         120.0        79.0  2625.0          18.6    82
396  31.0          4         119.0        82.0  2720.0          19.4    82

     origin                               name
338       1                   plymouth reliant
339       1                      buick skylark
340       1             dodge aries wagon (sw)
341       1                 chevrolet citation
342       1                   plymouth reliant
343       3                     toyota starlet
344       1                     plymouth champ
345       3                   honda civic 1300
346       3                             subaru
347       3                     datsun 210 mpg
348       3                      toyota tercel
349       3                        mazda glc 4
350       1                 plymouth horizon 4
351       1                     ford escort 4w
352       1                     ford escort 2h
353       2                   volkswagen jetta
355       3                      honda prelude
356       3                     toyota corolla
357       3                       datsun 200sx
358       3                          mazda 626
359       2          peugeot 505s turbo diesel
360       2                       volvo diesel
361       3                    toyota cressida
362       3                  datsun 810 maxima
363       1                      buick century
364       1              oldsmobile cutlass ls
365       1                    ford granada gl
366       1             chrysler lebaron salon
367       1                 chevrolet cavalier
368       1           chevrolet cavalier wagon
369       1          chevrolet cavalier 2-door
370       1         pontiac j2000 se hatchback
371       1                     dodge aries se
372       1                    pontiac phoenix
373       1               ford fairmont futura
374       2                volkswagen rabbit l
375       3                 mazda glc custom l
376       3                   mazda glc custom
377       1             plymouth horizon miser
378       1                     mercury lynx l
379       3                   nissan stanza xe
380       3                       honda accord
381       3                     toyota corolla
382       3                        honda civic
383       3                 honda civic (auto)
384       3                      datsun 310 gx
385       1              buick century limited
386       1  oldsmobile cutlass ciera (diesel)
387       1         chrysler lebaron medallion
388       1                     ford granada l
389       3                   toyota celica gt
390       1                  dodge charger 2.2
391       1                   chevrolet camaro
392       1                    ford mustang gl
393       2                          vw pickup
394       1                      dodge rampage
395       1                        ford ranger
396       1                         chevy s-10
\end{Verbatim}
\end{tcolorbox}
        
    However, if we pass in a list of strings to the \texttt{{[}{]}} method,
then we obtain a data frame containing the corresponding set of
\emph{columns}.

    \begin{tcolorbox}[breakable, size=fbox, boxrule=1pt, pad at break*=1mm,colback=cellbackground, colframe=cellborder]
\prompt{In}{incolor}{84}{\boxspacing}
\begin{Verbatim}[commandchars=\\\{\}]
\PY{n}{Auto}\PY{p}{[}\PY{p}{[}\PY{l+s+s1}{\PYZsq{}}\PY{l+s+s1}{mpg}\PY{l+s+s1}{\PYZsq{}}\PY{p}{,} \PY{l+s+s1}{\PYZsq{}}\PY{l+s+s1}{horsepower}\PY{l+s+s1}{\PYZsq{}}\PY{p}{]}\PY{p}{]}
\end{Verbatim}
\end{tcolorbox}

            \begin{tcolorbox}[breakable, size=fbox, boxrule=.5pt, pad at break*=1mm, opacityfill=0]
\prompt{Out}{outcolor}{84}{\boxspacing}
\begin{Verbatim}[commandchars=\\\{\}]
      mpg  horsepower
0    18.0       130.0
1    15.0       165.0
2    18.0       150.0
3    16.0       150.0
4    17.0       140.0
..    {\ldots}         {\ldots}
392  27.0        86.0
393  44.0        52.0
394  32.0        84.0
395  28.0        79.0
396  31.0        82.0

[392 rows x 2 columns]
\end{Verbatim}
\end{tcolorbox}
        
    Since we did not specify an \emph{index} column when we loaded our data
frame, the rows are labeled using integers 0 to 396.

    \begin{tcolorbox}[breakable, size=fbox, boxrule=1pt, pad at break*=1mm,colback=cellbackground, colframe=cellborder]
\prompt{In}{incolor}{85}{\boxspacing}
\begin{Verbatim}[commandchars=\\\{\}]
\PY{n}{Auto}\PY{o}{.}\PY{n}{index}
\end{Verbatim}
\end{tcolorbox}

            \begin{tcolorbox}[breakable, size=fbox, boxrule=.5pt, pad at break*=1mm, opacityfill=0]
\prompt{Out}{outcolor}{85}{\boxspacing}
\begin{Verbatim}[commandchars=\\\{\}]
Index([  0,   1,   2,   3,   4,   5,   6,   7,   8,   9,
       {\ldots}
       387, 388, 389, 390, 391, 392, 393, 394, 395, 396],
      dtype='int64', length=392)
\end{Verbatim}
\end{tcolorbox}
        
    We can use the \texttt{set\_index()} method to re-name the rows using
the contents of
\texttt{Auto{[}\textquotesingle{}name\textquotesingle{}{]}}.

    \begin{tcolorbox}[breakable, size=fbox, boxrule=1pt, pad at break*=1mm,colback=cellbackground, colframe=cellborder]
\prompt{In}{incolor}{86}{\boxspacing}
\begin{Verbatim}[commandchars=\\\{\}]
\PY{n}{Auto\PYZus{}re} \PY{o}{=} \PY{n}{Auto}\PY{o}{.}\PY{n}{set\PYZus{}index}\PY{p}{(}\PY{l+s+s1}{\PYZsq{}}\PY{l+s+s1}{name}\PY{l+s+s1}{\PYZsq{}}\PY{p}{)}
\PY{n}{Auto\PYZus{}re}
\end{Verbatim}
\end{tcolorbox}

            \begin{tcolorbox}[breakable, size=fbox, boxrule=.5pt, pad at break*=1mm, opacityfill=0]
\prompt{Out}{outcolor}{86}{\boxspacing}
\begin{Verbatim}[commandchars=\\\{\}]
                            mpg  cylinders  displacement  horsepower  weight  \textbackslash{}
name
chevrolet chevelle malibu  18.0          8         307.0       130.0  3504.0
buick skylark 320          15.0          8         350.0       165.0  3693.0
plymouth satellite         18.0          8         318.0       150.0  3436.0
amc rebel sst              16.0          8         304.0       150.0  3433.0
ford torino                17.0          8         302.0       140.0  3449.0
{\ldots}                         {\ldots}        {\ldots}           {\ldots}         {\ldots}     {\ldots}
ford mustang gl            27.0          4         140.0        86.0  2790.0
vw pickup                  44.0          4          97.0        52.0  2130.0
dodge rampage              32.0          4         135.0        84.0  2295.0
ford ranger                28.0          4         120.0        79.0  2625.0
chevy s-10                 31.0          4         119.0        82.0  2720.0

                           acceleration  year  origin
name
chevrolet chevelle malibu          12.0    70       1
buick skylark 320                  11.5    70       1
plymouth satellite                 11.0    70       1
amc rebel sst                      12.0    70       1
ford torino                        10.5    70       1
{\ldots}                                 {\ldots}   {\ldots}     {\ldots}
ford mustang gl                    15.6    82       1
vw pickup                          24.6    82       2
dodge rampage                      11.6    82       1
ford ranger                        18.6    82       1
chevy s-10                         19.4    82       1

[392 rows x 8 columns]
\end{Verbatim}
\end{tcolorbox}
        
    \begin{tcolorbox}[breakable, size=fbox, boxrule=1pt, pad at break*=1mm,colback=cellbackground, colframe=cellborder]
\prompt{In}{incolor}{87}{\boxspacing}
\begin{Verbatim}[commandchars=\\\{\}]
\PY{n}{Auto\PYZus{}re}\PY{o}{.}\PY{n}{columns}
\end{Verbatim}
\end{tcolorbox}

            \begin{tcolorbox}[breakable, size=fbox, boxrule=.5pt, pad at break*=1mm, opacityfill=0]
\prompt{Out}{outcolor}{87}{\boxspacing}
\begin{Verbatim}[commandchars=\\\{\}]
Index(['mpg', 'cylinders', 'displacement', 'horsepower', 'weight',
       'acceleration', 'year', 'origin'],
      dtype='object')
\end{Verbatim}
\end{tcolorbox}
        
    We see that the column \texttt{\textquotesingle{}name\textquotesingle{}}
is no longer there.

Now that the index has been set to \texttt{name}, we can access rows of
the data frame by \texttt{name} using the \texttt{\{loc{[}{]}}\} method
of \texttt{Auto}:

    \begin{tcolorbox}[breakable, size=fbox, boxrule=1pt, pad at break*=1mm,colback=cellbackground, colframe=cellborder]
\prompt{In}{incolor}{88}{\boxspacing}
\begin{Verbatim}[commandchars=\\\{\}]
\PY{n}{rows} \PY{o}{=} \PY{p}{[}\PY{l+s+s1}{\PYZsq{}}\PY{l+s+s1}{amc rebel sst}\PY{l+s+s1}{\PYZsq{}}\PY{p}{,} \PY{l+s+s1}{\PYZsq{}}\PY{l+s+s1}{ford torino}\PY{l+s+s1}{\PYZsq{}}\PY{p}{]}
\PY{n}{Auto\PYZus{}re}\PY{o}{.}\PY{n}{loc}\PY{p}{[}\PY{n}{rows}\PY{p}{]}
\end{Verbatim}
\end{tcolorbox}

            \begin{tcolorbox}[breakable, size=fbox, boxrule=.5pt, pad at break*=1mm, opacityfill=0]
\prompt{Out}{outcolor}{88}{\boxspacing}
\begin{Verbatim}[commandchars=\\\{\}]
                mpg  cylinders  displacement  horsepower  weight  \textbackslash{}
name
amc rebel sst  16.0          8         304.0       150.0  3433.0
ford torino    17.0          8         302.0       140.0  3449.0

               acceleration  year  origin
name
amc rebel sst          12.0    70       1
ford torino            10.5    70       1
\end{Verbatim}
\end{tcolorbox}
        
    As an alternative to using the index name, we could retrieve the 4th and
5th rows of \texttt{Auto} using the \texttt{\{iloc{[}{]}}\} method:

    \begin{tcolorbox}[breakable, size=fbox, boxrule=1pt, pad at break*=1mm,colback=cellbackground, colframe=cellborder]
\prompt{In}{incolor}{89}{\boxspacing}
\begin{Verbatim}[commandchars=\\\{\}]
\PY{n}{Auto\PYZus{}re}\PY{o}{.}\PY{n}{iloc}\PY{p}{[}\PY{p}{[}\PY{l+m+mi}{3}\PY{p}{,}\PY{l+m+mi}{4}\PY{p}{]}\PY{p}{]}
\end{Verbatim}
\end{tcolorbox}

            \begin{tcolorbox}[breakable, size=fbox, boxrule=.5pt, pad at break*=1mm, opacityfill=0]
\prompt{Out}{outcolor}{89}{\boxspacing}
\begin{Verbatim}[commandchars=\\\{\}]
                mpg  cylinders  displacement  horsepower  weight  \textbackslash{}
name
amc rebel sst  16.0          8         304.0       150.0  3433.0
ford torino    17.0          8         302.0       140.0  3449.0

               acceleration  year  origin
name
amc rebel sst          12.0    70       1
ford torino            10.5    70       1
\end{Verbatim}
\end{tcolorbox}
        
    We can also use it to retrieve the 1st, 3rd and and 4th columns of
\texttt{Auto\_re}:

    \begin{tcolorbox}[breakable, size=fbox, boxrule=1pt, pad at break*=1mm,colback=cellbackground, colframe=cellborder]
\prompt{In}{incolor}{90}{\boxspacing}
\begin{Verbatim}[commandchars=\\\{\}]
\PY{n}{Auto\PYZus{}re}\PY{o}{.}\PY{n}{iloc}\PY{p}{[}\PY{p}{:}\PY{p}{,}\PY{p}{[}\PY{l+m+mi}{0}\PY{p}{,}\PY{l+m+mi}{2}\PY{p}{,}\PY{l+m+mi}{3}\PY{p}{]}\PY{p}{]}
\end{Verbatim}
\end{tcolorbox}

            \begin{tcolorbox}[breakable, size=fbox, boxrule=.5pt, pad at break*=1mm, opacityfill=0]
\prompt{Out}{outcolor}{90}{\boxspacing}
\begin{Verbatim}[commandchars=\\\{\}]
                            mpg  displacement  horsepower
name
chevrolet chevelle malibu  18.0         307.0       130.0
buick skylark 320          15.0         350.0       165.0
plymouth satellite         18.0         318.0       150.0
amc rebel sst              16.0         304.0       150.0
ford torino                17.0         302.0       140.0
{\ldots}                         {\ldots}           {\ldots}         {\ldots}
ford mustang gl            27.0         140.0        86.0
vw pickup                  44.0          97.0        52.0
dodge rampage              32.0         135.0        84.0
ford ranger                28.0         120.0        79.0
chevy s-10                 31.0         119.0        82.0

[392 rows x 3 columns]
\end{Verbatim}
\end{tcolorbox}
        
    We can extract the 4th and 5th rows, as well as the 1st, 3rd and 4th
columns, using a single call to \texttt{iloc{[}{]}}:

    \begin{tcolorbox}[breakable, size=fbox, boxrule=1pt, pad at break*=1mm,colback=cellbackground, colframe=cellborder]
\prompt{In}{incolor}{91}{\boxspacing}
\begin{Verbatim}[commandchars=\\\{\}]
\PY{n}{Auto\PYZus{}re}\PY{o}{.}\PY{n}{iloc}\PY{p}{[}\PY{p}{[}\PY{l+m+mi}{3}\PY{p}{,}\PY{l+m+mi}{4}\PY{p}{]}\PY{p}{,}\PY{p}{[}\PY{l+m+mi}{0}\PY{p}{,}\PY{l+m+mi}{2}\PY{p}{,}\PY{l+m+mi}{3}\PY{p}{]}\PY{p}{]}
\end{Verbatim}
\end{tcolorbox}

            \begin{tcolorbox}[breakable, size=fbox, boxrule=.5pt, pad at break*=1mm, opacityfill=0]
\prompt{Out}{outcolor}{91}{\boxspacing}
\begin{Verbatim}[commandchars=\\\{\}]
                mpg  displacement  horsepower
name
amc rebel sst  16.0         304.0       150.0
ford torino    17.0         302.0       140.0
\end{Verbatim}
\end{tcolorbox}
        
    Index entries need not be unique: there are several cars in the data
frame named \texttt{ford\ galaxie\ 500}.

    \begin{tcolorbox}[breakable, size=fbox, boxrule=1pt, pad at break*=1mm,colback=cellbackground, colframe=cellborder]
\prompt{In}{incolor}{92}{\boxspacing}
\begin{Verbatim}[commandchars=\\\{\}]
\PY{n}{Auto\PYZus{}re}\PY{o}{.}\PY{n}{loc}\PY{p}{[}\PY{l+s+s1}{\PYZsq{}}\PY{l+s+s1}{ford galaxie 500}\PY{l+s+s1}{\PYZsq{}}\PY{p}{,} \PY{p}{[}\PY{l+s+s1}{\PYZsq{}}\PY{l+s+s1}{mpg}\PY{l+s+s1}{\PYZsq{}}\PY{p}{,} \PY{l+s+s1}{\PYZsq{}}\PY{l+s+s1}{origin}\PY{l+s+s1}{\PYZsq{}}\PY{p}{]}\PY{p}{]}
\end{Verbatim}
\end{tcolorbox}

            \begin{tcolorbox}[breakable, size=fbox, boxrule=.5pt, pad at break*=1mm, opacityfill=0]
\prompt{Out}{outcolor}{92}{\boxspacing}
\begin{Verbatim}[commandchars=\\\{\}]
                   mpg  origin
name
ford galaxie 500  15.0       1
ford galaxie 500  14.0       1
ford galaxie 500  14.0       1
\end{Verbatim}
\end{tcolorbox}
        
    \hypertarget{more-on-selecting-rows-and-columns}{%
\subsubsection{More on Selecting Rows and
Columns}\label{more-on-selecting-rows-and-columns}}

Suppose now that we want to create a data frame consisting of the
\texttt{weight} and \texttt{origin} of the subset of cars with
\texttt{year} greater than 80 --- i.e.~those built after 1980. To do
this, we first create a Boolean array that indexes the rows. The
\texttt{loc{[}{]}} method allows for Boolean entries as well as strings:

    \begin{tcolorbox}[breakable, size=fbox, boxrule=1pt, pad at break*=1mm,colback=cellbackground, colframe=cellborder]
\prompt{In}{incolor}{93}{\boxspacing}
\begin{Verbatim}[commandchars=\\\{\}]
\PY{n}{idx\PYZus{}80} \PY{o}{=} \PY{n}{Auto\PYZus{}re}\PY{p}{[}\PY{l+s+s1}{\PYZsq{}}\PY{l+s+s1}{year}\PY{l+s+s1}{\PYZsq{}}\PY{p}{]} \PY{o}{\PYZgt{}} \PY{l+m+mi}{80}
\PY{n}{Auto\PYZus{}re}\PY{o}{.}\PY{n}{loc}\PY{p}{[}\PY{n}{idx\PYZus{}80}\PY{p}{,} \PY{p}{[}\PY{l+s+s1}{\PYZsq{}}\PY{l+s+s1}{weight}\PY{l+s+s1}{\PYZsq{}}\PY{p}{,} \PY{l+s+s1}{\PYZsq{}}\PY{l+s+s1}{origin}\PY{l+s+s1}{\PYZsq{}}\PY{p}{]}\PY{p}{]}
\end{Verbatim}
\end{tcolorbox}

            \begin{tcolorbox}[breakable, size=fbox, boxrule=.5pt, pad at break*=1mm, opacityfill=0]
\prompt{Out}{outcolor}{93}{\boxspacing}
\begin{Verbatim}[commandchars=\\\{\}]
                                   weight  origin
name
plymouth reliant                   2490.0       1
buick skylark                      2635.0       1
dodge aries wagon (sw)             2620.0       1
chevrolet citation                 2725.0       1
plymouth reliant                   2385.0       1
toyota starlet                     1755.0       3
plymouth champ                     1875.0       1
honda civic 1300                   1760.0       3
subaru                             2065.0       3
datsun 210 mpg                     1975.0       3
toyota tercel                      2050.0       3
mazda glc 4                        1985.0       3
plymouth horizon 4                 2215.0       1
ford escort 4w                     2045.0       1
ford escort 2h                     2380.0       1
volkswagen jetta                   2190.0       2
honda prelude                      2210.0       3
toyota corolla                     2350.0       3
datsun 200sx                       2615.0       3
mazda 626                          2635.0       3
peugeot 505s turbo diesel          3230.0       2
volvo diesel                       3160.0       2
toyota cressida                    2900.0       3
datsun 810 maxima                  2930.0       3
buick century                      3415.0       1
oldsmobile cutlass ls              3725.0       1
ford granada gl                    3060.0       1
chrysler lebaron salon             3465.0       1
chevrolet cavalier                 2605.0       1
chevrolet cavalier wagon           2640.0       1
chevrolet cavalier 2-door          2395.0       1
pontiac j2000 se hatchback         2575.0       1
dodge aries se                     2525.0       1
pontiac phoenix                    2735.0       1
ford fairmont futura               2865.0       1
volkswagen rabbit l                1980.0       2
mazda glc custom l                 2025.0       3
mazda glc custom                   1970.0       3
plymouth horizon miser             2125.0       1
mercury lynx l                     2125.0       1
nissan stanza xe                   2160.0       3
honda accord                       2205.0       3
toyota corolla                     2245.0       3
honda civic                        1965.0       3
honda civic (auto)                 1965.0       3
datsun 310 gx                      1995.0       3
buick century limited              2945.0       1
oldsmobile cutlass ciera (diesel)  3015.0       1
chrysler lebaron medallion         2585.0       1
ford granada l                     2835.0       1
toyota celica gt                   2665.0       3
dodge charger 2.2                  2370.0       1
chevrolet camaro                   2950.0       1
ford mustang gl                    2790.0       1
vw pickup                          2130.0       2
dodge rampage                      2295.0       1
ford ranger                        2625.0       1
chevy s-10                         2720.0       1
\end{Verbatim}
\end{tcolorbox}
        
    To do this more concisely, we can use an anonymous function called a
\texttt{lambda}:

    \begin{tcolorbox}[breakable, size=fbox, boxrule=1pt, pad at break*=1mm,colback=cellbackground, colframe=cellborder]
\prompt{In}{incolor}{94}{\boxspacing}
\begin{Verbatim}[commandchars=\\\{\}]
\PY{n}{Auto\PYZus{}re}\PY{o}{.}\PY{n}{loc}\PY{p}{[}\PY{k}{lambda} \PY{n}{df}\PY{p}{:} \PY{n}{df}\PY{p}{[}\PY{l+s+s1}{\PYZsq{}}\PY{l+s+s1}{year}\PY{l+s+s1}{\PYZsq{}}\PY{p}{]} \PY{o}{\PYZgt{}} \PY{l+m+mi}{80}\PY{p}{,} \PY{p}{[}\PY{l+s+s1}{\PYZsq{}}\PY{l+s+s1}{weight}\PY{l+s+s1}{\PYZsq{}}\PY{p}{,} \PY{l+s+s1}{\PYZsq{}}\PY{l+s+s1}{origin}\PY{l+s+s1}{\PYZsq{}}\PY{p}{]}\PY{p}{]}
\end{Verbatim}
\end{tcolorbox}

            \begin{tcolorbox}[breakable, size=fbox, boxrule=.5pt, pad at break*=1mm, opacityfill=0]
\prompt{Out}{outcolor}{94}{\boxspacing}
\begin{Verbatim}[commandchars=\\\{\}]
                                   weight  origin
name
plymouth reliant                   2490.0       1
buick skylark                      2635.0       1
dodge aries wagon (sw)             2620.0       1
chevrolet citation                 2725.0       1
plymouth reliant                   2385.0       1
toyota starlet                     1755.0       3
plymouth champ                     1875.0       1
honda civic 1300                   1760.0       3
subaru                             2065.0       3
datsun 210 mpg                     1975.0       3
toyota tercel                      2050.0       3
mazda glc 4                        1985.0       3
plymouth horizon 4                 2215.0       1
ford escort 4w                     2045.0       1
ford escort 2h                     2380.0       1
volkswagen jetta                   2190.0       2
honda prelude                      2210.0       3
toyota corolla                     2350.0       3
datsun 200sx                       2615.0       3
mazda 626                          2635.0       3
peugeot 505s turbo diesel          3230.0       2
volvo diesel                       3160.0       2
toyota cressida                    2900.0       3
datsun 810 maxima                  2930.0       3
buick century                      3415.0       1
oldsmobile cutlass ls              3725.0       1
ford granada gl                    3060.0       1
chrysler lebaron salon             3465.0       1
chevrolet cavalier                 2605.0       1
chevrolet cavalier wagon           2640.0       1
chevrolet cavalier 2-door          2395.0       1
pontiac j2000 se hatchback         2575.0       1
dodge aries se                     2525.0       1
pontiac phoenix                    2735.0       1
ford fairmont futura               2865.0       1
volkswagen rabbit l                1980.0       2
mazda glc custom l                 2025.0       3
mazda glc custom                   1970.0       3
plymouth horizon miser             2125.0       1
mercury lynx l                     2125.0       1
nissan stanza xe                   2160.0       3
honda accord                       2205.0       3
toyota corolla                     2245.0       3
honda civic                        1965.0       3
honda civic (auto)                 1965.0       3
datsun 310 gx                      1995.0       3
buick century limited              2945.0       1
oldsmobile cutlass ciera (diesel)  3015.0       1
chrysler lebaron medallion         2585.0       1
ford granada l                     2835.0       1
toyota celica gt                   2665.0       3
dodge charger 2.2                  2370.0       1
chevrolet camaro                   2950.0       1
ford mustang gl                    2790.0       1
vw pickup                          2130.0       2
dodge rampage                      2295.0       1
ford ranger                        2625.0       1
chevy s-10                         2720.0       1
\end{Verbatim}
\end{tcolorbox}
        
    The \texttt{lambda} call creates a function that takes a single
argument, here \texttt{df}, and returns
\texttt{df{[}\textquotesingle{}year\textquotesingle{}{]}\textgreater{}80}.
Since it is created inside the \texttt{loc{[}{]}} method for the
dataframe \texttt{Auto\_re}, that dataframe will be the argument
supplied. As another example of using a \texttt{lambda}, suppose that we
want all cars built after 1980 that achieve greater than 30 miles per
gallon:

    \begin{tcolorbox}[breakable, size=fbox, boxrule=1pt, pad at break*=1mm,colback=cellbackground, colframe=cellborder]
\prompt{In}{incolor}{95}{\boxspacing}
\begin{Verbatim}[commandchars=\\\{\}]
\PY{n}{Auto\PYZus{}re}\PY{o}{.}\PY{n}{loc}\PY{p}{[}\PY{k}{lambda} \PY{n}{df}\PY{p}{:} \PY{p}{(}\PY{n}{df}\PY{p}{[}\PY{l+s+s1}{\PYZsq{}}\PY{l+s+s1}{year}\PY{l+s+s1}{\PYZsq{}}\PY{p}{]} \PY{o}{\PYZgt{}} \PY{l+m+mi}{80}\PY{p}{)} \PY{o}{\PYZam{}} \PY{p}{(}\PY{n}{df}\PY{p}{[}\PY{l+s+s1}{\PYZsq{}}\PY{l+s+s1}{mpg}\PY{l+s+s1}{\PYZsq{}}\PY{p}{]} \PY{o}{\PYZgt{}} \PY{l+m+mi}{30}\PY{p}{)}\PY{p}{,}
            \PY{p}{[}\PY{l+s+s1}{\PYZsq{}}\PY{l+s+s1}{weight}\PY{l+s+s1}{\PYZsq{}}\PY{p}{,} \PY{l+s+s1}{\PYZsq{}}\PY{l+s+s1}{origin}\PY{l+s+s1}{\PYZsq{}}\PY{p}{]}
           \PY{p}{]}
\end{Verbatim}
\end{tcolorbox}

            \begin{tcolorbox}[breakable, size=fbox, boxrule=.5pt, pad at break*=1mm, opacityfill=0]
\prompt{Out}{outcolor}{95}{\boxspacing}
\begin{Verbatim}[commandchars=\\\{\}]
                                   weight  origin
name
toyota starlet                     1755.0       3
plymouth champ                     1875.0       1
honda civic 1300                   1760.0       3
subaru                             2065.0       3
datsun 210 mpg                     1975.0       3
toyota tercel                      2050.0       3
mazda glc 4                        1985.0       3
plymouth horizon 4                 2215.0       1
ford escort 4w                     2045.0       1
volkswagen jetta                   2190.0       2
honda prelude                      2210.0       3
toyota corolla                     2350.0       3
datsun 200sx                       2615.0       3
mazda 626                          2635.0       3
volvo diesel                       3160.0       2
chevrolet cavalier 2-door          2395.0       1
pontiac j2000 se hatchback         2575.0       1
volkswagen rabbit l                1980.0       2
mazda glc custom l                 2025.0       3
mazda glc custom                   1970.0       3
plymouth horizon miser             2125.0       1
mercury lynx l                     2125.0       1
nissan stanza xe                   2160.0       3
honda accord                       2205.0       3
toyota corolla                     2245.0       3
honda civic                        1965.0       3
honda civic (auto)                 1965.0       3
datsun 310 gx                      1995.0       3
oldsmobile cutlass ciera (diesel)  3015.0       1
toyota celica gt                   2665.0       3
dodge charger 2.2                  2370.0       1
vw pickup                          2130.0       2
dodge rampage                      2295.0       1
chevy s-10                         2720.0       1
\end{Verbatim}
\end{tcolorbox}
        
    The symbol \texttt{\&} computes an element-wise \emph{and} operation. As
another example, suppose that we want to retrieve all \texttt{Ford} and
\texttt{Datsun} cars with \texttt{displacement} less than 300. We check
whether each \texttt{name} entry contains either the string
\texttt{ford} or \texttt{datsun} using the \texttt{str.contains()}
method of the \texttt{index} attribute of the dataframe:

    \begin{tcolorbox}[breakable, size=fbox, boxrule=1pt, pad at break*=1mm,colback=cellbackground, colframe=cellborder]
\prompt{In}{incolor}{96}{\boxspacing}
\begin{Verbatim}[commandchars=\\\{\}]
\PY{n}{Auto\PYZus{}re}\PY{o}{.}\PY{n}{loc}\PY{p}{[}\PY{k}{lambda} \PY{n}{df}\PY{p}{:} \PY{p}{(}\PY{n}{df}\PY{p}{[}\PY{l+s+s1}{\PYZsq{}}\PY{l+s+s1}{displacement}\PY{l+s+s1}{\PYZsq{}}\PY{p}{]} \PY{o}{\PYZlt{}} \PY{l+m+mi}{300}\PY{p}{)}
                       \PY{o}{\PYZam{}} \PY{p}{(}\PY{n}{df}\PY{o}{.}\PY{n}{index}\PY{o}{.}\PY{n}{str}\PY{o}{.}\PY{n}{contains}\PY{p}{(}\PY{l+s+s1}{\PYZsq{}}\PY{l+s+s1}{ford}\PY{l+s+s1}{\PYZsq{}}\PY{p}{)}
                       \PY{o}{|} \PY{n}{df}\PY{o}{.}\PY{n}{index}\PY{o}{.}\PY{n}{str}\PY{o}{.}\PY{n}{contains}\PY{p}{(}\PY{l+s+s1}{\PYZsq{}}\PY{l+s+s1}{datsun}\PY{l+s+s1}{\PYZsq{}}\PY{p}{)}\PY{p}{)}\PY{p}{,}
            \PY{p}{[}\PY{l+s+s1}{\PYZsq{}}\PY{l+s+s1}{weight}\PY{l+s+s1}{\PYZsq{}}\PY{p}{,} \PY{l+s+s1}{\PYZsq{}}\PY{l+s+s1}{origin}\PY{l+s+s1}{\PYZsq{}}\PY{p}{]}
           \PY{p}{]}
\end{Verbatim}
\end{tcolorbox}

            \begin{tcolorbox}[breakable, size=fbox, boxrule=.5pt, pad at break*=1mm, opacityfill=0]
\prompt{Out}{outcolor}{96}{\boxspacing}
\begin{Verbatim}[commandchars=\\\{\}]
                       weight  origin
name
ford maverick          2587.0       1
datsun pl510           2130.0       3
datsun pl510           2130.0       3
ford torino 500        3302.0       1
ford mustang           3139.0       1
datsun 1200            1613.0       3
ford pinto runabout    2226.0       1
ford pinto (sw)        2395.0       1
datsun 510 (sw)        2288.0       3
ford maverick          3021.0       1
datsun 610             2379.0       3
ford pinto             2310.0       1
datsun b210            1950.0       3
ford pinto             2451.0       1
datsun 710             2003.0       3
ford maverick          3158.0       1
ford pinto             2639.0       1
datsun 710             2545.0       3
ford pinto             2984.0       1
ford maverick          3012.0       1
ford granada ghia      3574.0       1
datsun b-210           1990.0       3
ford pinto             2565.0       1
datsun f-10 hatchback  1945.0       3
ford granada           3525.0       1
ford mustang ii 2+2    2755.0       1
datsun 810             2815.0       3
ford fiesta            1800.0       1
datsun b210 gx         2070.0       3
ford fairmont (auto)   2965.0       1
ford fairmont (man)    2720.0       1
datsun 510             2300.0       3
datsun 200-sx          2405.0       3
ford fairmont 4        2890.0       1
datsun 210             2020.0       3
datsun 310             2019.0       3
ford fairmont          2870.0       1
datsun 510 hatchback   2434.0       3
datsun 210             2110.0       3
datsun 280-zx          2910.0       3
datsun 210 mpg         1975.0       3
ford escort 4w         2045.0       1
ford escort 2h         2380.0       1
datsun 200sx           2615.0       3
datsun 810 maxima      2930.0       3
ford granada gl        3060.0       1
ford fairmont futura   2865.0       1
datsun 310 gx          1995.0       3
ford granada l         2835.0       1
ford mustang gl        2790.0       1
ford ranger            2625.0       1
\end{Verbatim}
\end{tcolorbox}
        
    Here, the symbol \texttt{\textbar{}} computes an element-wise \emph{or}
operation.

In summary, a powerful set of operations is available to index the rows
and columns of data frames. For integer based queries, use the
\texttt{iloc{[}{]}} method. For string and Boolean selections, use the
\texttt{loc{[}{]}} method. For functional queries that filter rows, use
the \texttt{loc{[}{]}} method with a function (typically a
\texttt{lambda}) in the rows argument.

\hypertarget{for-loops}{%
\subsection{For Loops}\label{for-loops}}

A \texttt{for} loop is a standard tool in many languages that repeatedly
evaluates some chunk of code while varying different values inside the
code. For example, suppose we loop over elements of a list and compute
their sum.

    \begin{tcolorbox}[breakable, size=fbox, boxrule=1pt, pad at break*=1mm,colback=cellbackground, colframe=cellborder]
\prompt{In}{incolor}{97}{\boxspacing}
\begin{Verbatim}[commandchars=\\\{\}]
\PY{n}{total} \PY{o}{=} \PY{l+m+mi}{0}
\PY{k}{for} \PY{n}{value} \PY{o+ow}{in} \PY{p}{[}\PY{l+m+mi}{3}\PY{p}{,}\PY{l+m+mi}{2}\PY{p}{,}\PY{l+m+mi}{19}\PY{p}{]}\PY{p}{:}
    \PY{n}{total} \PY{o}{+}\PY{o}{=} \PY{n}{value}
\PY{n+nb}{print}\PY{p}{(}\PY{l+s+s1}{\PYZsq{}}\PY{l+s+s1}{Total is: }\PY{l+s+si}{\PYZob{}0\PYZcb{}}\PY{l+s+s1}{\PYZsq{}}\PY{o}{.}\PY{n}{format}\PY{p}{(}\PY{n}{total}\PY{p}{)}\PY{p}{)}
\end{Verbatim}
\end{tcolorbox}

    \begin{Verbatim}[commandchars=\\\{\}]
Total is: 24
    \end{Verbatim}

    The indented code beneath the line with the \texttt{for} statement is
run for each value in the sequence specified in the \texttt{for}
statement. The loop ends either when the cell ends or when code is
indented at the same level as the original \texttt{for} statement. We
see that the final line above which prints the total is executed only
once after the for loop has terminated. Loops can be nested by
additional indentation.

    \begin{tcolorbox}[breakable, size=fbox, boxrule=1pt, pad at break*=1mm,colback=cellbackground, colframe=cellborder]
\prompt{In}{incolor}{98}{\boxspacing}
\begin{Verbatim}[commandchars=\\\{\}]
\PY{n}{total} \PY{o}{=} \PY{l+m+mi}{0}
\PY{k}{for} \PY{n}{value} \PY{o+ow}{in} \PY{p}{[}\PY{l+m+mi}{2}\PY{p}{,}\PY{l+m+mi}{3}\PY{p}{,}\PY{l+m+mi}{19}\PY{p}{]}\PY{p}{:}
    \PY{k}{for} \PY{n}{weight} \PY{o+ow}{in} \PY{p}{[}\PY{l+m+mi}{3}\PY{p}{,} \PY{l+m+mi}{2}\PY{p}{,} \PY{l+m+mi}{1}\PY{p}{]}\PY{p}{:}
        \PY{n}{total} \PY{o}{+}\PY{o}{=} \PY{n}{value} \PY{o}{*} \PY{n}{weight}
\PY{n+nb}{print}\PY{p}{(}\PY{l+s+s1}{\PYZsq{}}\PY{l+s+s1}{Total is: }\PY{l+s+si}{\PYZob{}0\PYZcb{}}\PY{l+s+s1}{\PYZsq{}}\PY{o}{.}\PY{n}{format}\PY{p}{(}\PY{n}{total}\PY{p}{)}\PY{p}{)}
\end{Verbatim}
\end{tcolorbox}

    \begin{Verbatim}[commandchars=\\\{\}]
Total is: 144
    \end{Verbatim}

    Above, we summed over each combination of \texttt{value} and
\texttt{weight}. We also took advantage of the \emph{increment} notation
in \texttt{Python}: the expression \texttt{a\ +=\ b} is equivalent to
\texttt{a\ =\ a\ +\ b}. Besides being a convenient notation, this can
save time in computationally heavy tasks in which the intermediate value
of \texttt{a+b} need not be explicitly created.

Perhaps a more common task would be to sum over
\texttt{(value,\ weight)} pairs. For instance, to compute the average
value of a random variable that takes on possible values 2, 3 or 19 with
probability 0.2, 0.3, 0.5 respectively we would compute the weighted
sum. Tasks such as this can often be accomplished using the
\texttt{zip()} function that loops over a sequence of tuples.

    \begin{tcolorbox}[breakable, size=fbox, boxrule=1pt, pad at break*=1mm,colback=cellbackground, colframe=cellborder]
\prompt{In}{incolor}{99}{\boxspacing}
\begin{Verbatim}[commandchars=\\\{\}]
\PY{n}{total} \PY{o}{=} \PY{l+m+mi}{0}
\PY{k}{for} \PY{n}{value}\PY{p}{,} \PY{n}{weight} \PY{o+ow}{in} \PY{n+nb}{zip}\PY{p}{(}\PY{p}{[}\PY{l+m+mi}{2}\PY{p}{,}\PY{l+m+mi}{3}\PY{p}{,}\PY{l+m+mi}{19}\PY{p}{]}\PY{p}{,}
                         \PY{p}{[}\PY{l+m+mf}{0.2}\PY{p}{,}\PY{l+m+mf}{0.3}\PY{p}{,}\PY{l+m+mf}{0.5}\PY{p}{]}\PY{p}{)}\PY{p}{:}
    \PY{n}{total} \PY{o}{+}\PY{o}{=} \PY{n}{weight} \PY{o}{*} \PY{n}{value}
\PY{n+nb}{print}\PY{p}{(}\PY{l+s+s1}{\PYZsq{}}\PY{l+s+s1}{Weighted average is: }\PY{l+s+si}{\PYZob{}0\PYZcb{}}\PY{l+s+s1}{\PYZsq{}}\PY{o}{.}\PY{n}{format}\PY{p}{(}\PY{n}{total}\PY{p}{)}\PY{p}{)}
\end{Verbatim}
\end{tcolorbox}

    \begin{Verbatim}[commandchars=\\\{\}]
Weighted average is: 10.8
    \end{Verbatim}

    \hypertarget{string-formatting}{%
\subsubsection{String Formatting}\label{string-formatting}}

In the code chunk above we also printed a string displaying the total.
However, the object \texttt{total} is an integer and not a string.
Inserting the value of something into a string is a common task, made
simple using some of the powerful string formatting tools in
\texttt{Python}. Many data cleaning tasks involve manipulating and
programmatically producing strings.

For example we may want to loop over the columns of a data frame and
print the percent missing in each column. Let's create a data frame
\texttt{D} with columns in which 20\% of the entries are missing
i.e.~set to \texttt{np.nan}. We'll create the values in \texttt{D} from
a normal distribution with mean 0 and variance 1 using
\texttt{rng.standard\_normal()} and then overwrite some random entries
using \texttt{rng.choice()}.

    \begin{tcolorbox}[breakable, size=fbox, boxrule=1pt, pad at break*=1mm,colback=cellbackground, colframe=cellborder]
\prompt{In}{incolor}{100}{\boxspacing}
\begin{Verbatim}[commandchars=\\\{\}]
\PY{n}{rng} \PY{o}{=} \PY{n}{np}\PY{o}{.}\PY{n}{random}\PY{o}{.}\PY{n}{default\PYZus{}rng}\PY{p}{(}\PY{l+m+mi}{1}\PY{p}{)}
\PY{n}{A} \PY{o}{=} \PY{n}{rng}\PY{o}{.}\PY{n}{standard\PYZus{}normal}\PY{p}{(}\PY{p}{(}\PY{l+m+mi}{127}\PY{p}{,} \PY{l+m+mi}{5}\PY{p}{)}\PY{p}{)}
\PY{n}{M} \PY{o}{=} \PY{n}{rng}\PY{o}{.}\PY{n}{choice}\PY{p}{(}\PY{p}{[}\PY{l+m+mi}{0}\PY{p}{,} \PY{n}{np}\PY{o}{.}\PY{n}{nan}\PY{p}{]}\PY{p}{,} \PY{n}{p}\PY{o}{=}\PY{p}{[}\PY{l+m+mf}{0.8}\PY{p}{,}\PY{l+m+mf}{0.2}\PY{p}{]}\PY{p}{,} \PY{n}{size}\PY{o}{=}\PY{n}{A}\PY{o}{.}\PY{n}{shape}\PY{p}{)}
\PY{n}{A} \PY{o}{+}\PY{o}{=} \PY{n}{M}
\PY{n}{D} \PY{o}{=} \PY{n}{pd}\PY{o}{.}\PY{n}{DataFrame}\PY{p}{(}\PY{n}{A}\PY{p}{,} \PY{n}{columns}\PY{o}{=}\PY{p}{[}\PY{l+s+s1}{\PYZsq{}}\PY{l+s+s1}{food}\PY{l+s+s1}{\PYZsq{}}\PY{p}{,}
                             \PY{l+s+s1}{\PYZsq{}}\PY{l+s+s1}{bar}\PY{l+s+s1}{\PYZsq{}}\PY{p}{,}
                             \PY{l+s+s1}{\PYZsq{}}\PY{l+s+s1}{pickle}\PY{l+s+s1}{\PYZsq{}}\PY{p}{,}
                             \PY{l+s+s1}{\PYZsq{}}\PY{l+s+s1}{snack}\PY{l+s+s1}{\PYZsq{}}\PY{p}{,}
                             \PY{l+s+s1}{\PYZsq{}}\PY{l+s+s1}{popcorn}\PY{l+s+s1}{\PYZsq{}}\PY{p}{]}\PY{p}{)}
\PY{n}{D}\PY{p}{[}\PY{p}{:}\PY{l+m+mi}{3}\PY{p}{]}
\end{Verbatim}
\end{tcolorbox}

            \begin{tcolorbox}[breakable, size=fbox, boxrule=.5pt, pad at break*=1mm, opacityfill=0]
\prompt{Out}{outcolor}{100}{\boxspacing}
\begin{Verbatim}[commandchars=\\\{\}]
       food       bar    pickle     snack   popcorn
0  0.345584  0.821618  0.330437 -1.303157       NaN
1       NaN -0.536953  0.581118  0.364572  0.294132
2       NaN  0.546713       NaN -0.162910 -0.482119
\end{Verbatim}
\end{tcolorbox}
        
    \begin{tcolorbox}[breakable, size=fbox, boxrule=1pt, pad at break*=1mm,colback=cellbackground, colframe=cellborder]
\prompt{In}{incolor}{101}{\boxspacing}
\begin{Verbatim}[commandchars=\\\{\}]
\PY{k}{for} \PY{n}{col} \PY{o+ow}{in} \PY{n}{D}\PY{o}{.}\PY{n}{columns}\PY{p}{:}
    \PY{n}{template} \PY{o}{=} \PY{l+s+s1}{\PYZsq{}}\PY{l+s+s1}{Column }\PY{l+s+s1}{\PYZdq{}}\PY{l+s+si}{\PYZob{}0\PYZcb{}}\PY{l+s+s1}{\PYZdq{}}\PY{l+s+s1}{ has }\PY{l+s+si}{\PYZob{}1:.2\PYZpc{}\PYZcb{}}\PY{l+s+s1}{ missing values}\PY{l+s+s1}{\PYZsq{}}
    \PY{n+nb}{print}\PY{p}{(}\PY{n}{template}\PY{o}{.}\PY{n}{format}\PY{p}{(}\PY{n}{col}\PY{p}{,}
          \PY{n}{np}\PY{o}{.}\PY{n}{isnan}\PY{p}{(}\PY{n}{D}\PY{p}{[}\PY{n}{col}\PY{p}{]}\PY{p}{)}\PY{o}{.}\PY{n}{mean}\PY{p}{(}\PY{p}{)}\PY{p}{)}\PY{p}{)}
\end{Verbatim}
\end{tcolorbox}

    \begin{Verbatim}[commandchars=\\\{\}]
Column "food" has 16.54\% missing values
Column "bar" has 25.98\% missing values
Column "pickle" has 29.13\% missing values
Column "snack" has 21.26\% missing values
Column "popcorn" has 22.83\% missing values
    \end{Verbatim}

    We see that the \texttt{template.format()} method expects two arguments
\texttt{\{0\}} and \texttt{\{1:.2\%\}}, and the latter includes some
formatting information. In particular, it specifies that the second
argument should be expressed as a percent with two decimal digits.

The reference
\href{https://docs.python.org/3/library/string.html}{docs.python.org/3/library/string.html}
includes many helpful and more complex examples.

    \hypertarget{additional-graphical-and-numerical-summaries}{%
\subsection{Additional Graphical and Numerical
Summaries}\label{additional-graphical-and-numerical-summaries}}

We can use the \texttt{ax.plot()} or \texttt{ax.scatter()} functions to
display the quantitative variables. However, simply typing the variable
names will produce an error message, because \texttt{Python} does not
know to look in the \texttt{Auto} data set for those variables.

    \begin{tcolorbox}[breakable, size=fbox, boxrule=1pt, pad at break*=1mm,colback=cellbackground, colframe=cellborder]
\prompt{In}{incolor}{102}{\boxspacing}
\begin{Verbatim}[commandchars=\\\{\}]
\PY{n}{fig}\PY{p}{,} \PY{n}{ax} \PY{o}{=} \PY{n}{subplots}\PY{p}{(}\PY{n}{figsize}\PY{o}{=}\PY{p}{(}\PY{l+m+mi}{8}\PY{p}{,} \PY{l+m+mi}{8}\PY{p}{)}\PY{p}{)}
\PY{n}{ax}\PY{o}{.}\PY{n}{plot}\PY{p}{(}\PY{n}{horsepower}\PY{p}{,} \PY{n}{mpg}\PY{p}{,} \PY{l+s+s1}{\PYZsq{}}\PY{l+s+s1}{o}\PY{l+s+s1}{\PYZsq{}}\PY{p}{)}\PY{p}{;}
\end{Verbatim}
\end{tcolorbox}

    \begin{Verbatim}[commandchars=\\\{\}, frame=single, framerule=2mm, rulecolor=\color{outerrorbackground}]
\textcolor{ansi-red}{---------------------------------------------------------------------------}
\textcolor{ansi-red}{NameError}                                 Traceback (most recent call last)
\textcolor{ansi-cyan}{Cell}\textcolor{ansi-cyan}{ }\textcolor{ansi-green}{In[102]}\textcolor{ansi-green}{, line 2}
\textcolor{ansi-green}{      1} fig, ax = subplots(figsize=(\textcolor{ansi-green}{8}, \textcolor{ansi-green}{8}))
\textcolor{ansi-green}{----> }\textcolor{ansi-green}{2} ax.plot(\setlength{\fboxsep}{0pt}\colorbox{ansi-yellow}{horsepower\strut}, mpg, \textcolor{ansi-yellow}{'}\textcolor{ansi-yellow}{o}\textcolor{ansi-yellow}{'});

\textcolor{ansi-red}{NameError}: name 'horsepower' is not defined
    \end{Verbatim}

    \begin{center}
    \adjustimage{max size={0.9\linewidth}{0.9\paperheight}}{artifacts/latex/Ch02-statlearn-lab_files/artifacts/latex/Ch02-statlearn-lab_238_1.png}
    \end{center}
    { \hspace*{\fill} \\}
    
    We can address this by accessing the columns directly:

    \begin{tcolorbox}[breakable, size=fbox, boxrule=1pt, pad at break*=1mm,colback=cellbackground, colframe=cellborder]
\prompt{In}{incolor}{103}{\boxspacing}
\begin{Verbatim}[commandchars=\\\{\}]
\PY{n}{fig}\PY{p}{,} \PY{n}{ax} \PY{o}{=} \PY{n}{subplots}\PY{p}{(}\PY{n}{figsize}\PY{o}{=}\PY{p}{(}\PY{l+m+mi}{8}\PY{p}{,} \PY{l+m+mi}{8}\PY{p}{)}\PY{p}{)}
\PY{n}{ax}\PY{o}{.}\PY{n}{plot}\PY{p}{(}\PY{n}{Auto}\PY{p}{[}\PY{l+s+s1}{\PYZsq{}}\PY{l+s+s1}{horsepower}\PY{l+s+s1}{\PYZsq{}}\PY{p}{]}\PY{p}{,} \PY{n}{Auto}\PY{p}{[}\PY{l+s+s1}{\PYZsq{}}\PY{l+s+s1}{mpg}\PY{l+s+s1}{\PYZsq{}}\PY{p}{]}\PY{p}{,} \PY{l+s+s1}{\PYZsq{}}\PY{l+s+s1}{o}\PY{l+s+s1}{\PYZsq{}}\PY{p}{)}\PY{p}{;}
\end{Verbatim}
\end{tcolorbox}

    \begin{center}
    \adjustimage{max size={0.9\linewidth}{0.9\paperheight}}{artifacts/latex/Ch02-statlearn-lab_files/artifacts/latex/Ch02-statlearn-lab_240_0.png}
    \end{center}
    { \hspace*{\fill} \\}
    
    Alternatively, we can use the \texttt{plot()} method with the call
\texttt{Auto.plot()}. Using this method, the variables can be accessed
by name. The plot methods of a data frame return a familiar object: an
axes. We can use it to update the plot as we did previously:

    \begin{tcolorbox}[breakable, size=fbox, boxrule=1pt, pad at break*=1mm,colback=cellbackground, colframe=cellborder]
\prompt{In}{incolor}{104}{\boxspacing}
\begin{Verbatim}[commandchars=\\\{\}]
\PY{n}{ax} \PY{o}{=} \PY{n}{Auto}\PY{o}{.}\PY{n}{plot}\PY{o}{.}\PY{n}{scatter}\PY{p}{(}\PY{l+s+s1}{\PYZsq{}}\PY{l+s+s1}{horsepower}\PY{l+s+s1}{\PYZsq{}}\PY{p}{,} \PY{l+s+s1}{\PYZsq{}}\PY{l+s+s1}{mpg}\PY{l+s+s1}{\PYZsq{}}\PY{p}{)}
\PY{n}{ax}\PY{o}{.}\PY{n}{set\PYZus{}title}\PY{p}{(}\PY{l+s+s1}{\PYZsq{}}\PY{l+s+s1}{Horsepower vs. MPG}\PY{l+s+s1}{\PYZsq{}}\PY{p}{)}\PY{p}{;}
\end{Verbatim}
\end{tcolorbox}

    \begin{center}
    \adjustimage{max size={0.9\linewidth}{0.9\paperheight}}{artifacts/latex/Ch02-statlearn-lab_files/artifacts/latex/Ch02-statlearn-lab_242_0.png}
    \end{center}
    { \hspace*{\fill} \\}
    
    If we want to save the figure that contains a given axes, we can find
the relevant figure by accessing the \texttt{figure} attribute:

    \begin{tcolorbox}[breakable, size=fbox, boxrule=1pt, pad at break*=1mm,colback=cellbackground, colframe=cellborder]
\prompt{In}{incolor}{105}{\boxspacing}
\begin{Verbatim}[commandchars=\\\{\}]
\PY{n}{fig} \PY{o}{=} \PY{n}{ax}\PY{o}{.}\PY{n}{figure}
\PY{n}{fig}\PY{o}{.}\PY{n}{savefig}\PY{p}{(}\PY{l+s+s1}{\PYZsq{}}\PY{l+s+s1}{horsepower\PYZus{}mpg.png}\PY{l+s+s1}{\PYZsq{}}\PY{p}{)}\PY{p}{;}
\end{Verbatim}
\end{tcolorbox}

    We can further instruct the data frame to plot to a particular axes
object. In this case the corresponding \texttt{plot()} method will
return the modified axes we passed in as an argument. Note that when we
request a one-dimensional grid of plots, the object \texttt{axes} is
similarly one-dimensional. We place our scatter plot in the middle plot
of a row of three plots within a figure.

    \begin{tcolorbox}[breakable, size=fbox, boxrule=1pt, pad at break*=1mm,colback=cellbackground, colframe=cellborder]
\prompt{In}{incolor}{106}{\boxspacing}
\begin{Verbatim}[commandchars=\\\{\}]
\PY{n}{fig}\PY{p}{,} \PY{n}{axes} \PY{o}{=} \PY{n}{subplots}\PY{p}{(}\PY{n}{ncols}\PY{o}{=}\PY{l+m+mi}{3}\PY{p}{,} \PY{n}{figsize}\PY{o}{=}\PY{p}{(}\PY{l+m+mi}{15}\PY{p}{,} \PY{l+m+mi}{5}\PY{p}{)}\PY{p}{)}
\PY{n}{Auto}\PY{o}{.}\PY{n}{plot}\PY{o}{.}\PY{n}{scatter}\PY{p}{(}\PY{l+s+s1}{\PYZsq{}}\PY{l+s+s1}{horsepower}\PY{l+s+s1}{\PYZsq{}}\PY{p}{,} \PY{l+s+s1}{\PYZsq{}}\PY{l+s+s1}{mpg}\PY{l+s+s1}{\PYZsq{}}\PY{p}{,} \PY{n}{ax}\PY{o}{=}\PY{n}{axes}\PY{p}{[}\PY{l+m+mi}{1}\PY{p}{]}\PY{p}{)}\PY{p}{;}
\end{Verbatim}
\end{tcolorbox}

    \begin{center}
    \adjustimage{max size={0.9\linewidth}{0.9\paperheight}}{artifacts/latex/Ch02-statlearn-lab_files/artifacts/latex/Ch02-statlearn-lab_246_0.png}
    \end{center}
    { \hspace*{\fill} \\}
    
    Note also that the columns of a data frame can be accessed as
attributes: try typing in \texttt{Auto.horsepower}.

    We now consider the \texttt{cylinders} variable. Typing in
\texttt{Auto.cylinders.dtype} reveals that it is being treated as a
quantitative variable. However, since there is only a small number of
possible values for this variable, we may wish to treat it as
qualitative. Below, we replace the \texttt{cylinders} column with a
categorical version of \texttt{Auto.cylinders}. The function
\texttt{pd.Series()} owes its name to the fact that \texttt{pandas} is
often used in time series applications.

    \begin{tcolorbox}[breakable, size=fbox, boxrule=1pt, pad at break*=1mm,colback=cellbackground, colframe=cellborder]
\prompt{In}{incolor}{107}{\boxspacing}
\begin{Verbatim}[commandchars=\\\{\}]
\PY{n}{Auto}\PY{o}{.}\PY{n}{cylinders} \PY{o}{=} \PY{n}{pd}\PY{o}{.}\PY{n}{Series}\PY{p}{(}\PY{n}{Auto}\PY{o}{.}\PY{n}{cylinders}\PY{p}{,} \PY{n}{dtype}\PY{o}{=}\PY{l+s+s1}{\PYZsq{}}\PY{l+s+s1}{category}\PY{l+s+s1}{\PYZsq{}}\PY{p}{)}
\PY{n}{Auto}\PY{o}{.}\PY{n}{cylinders}\PY{o}{.}\PY{n}{dtype}
\end{Verbatim}
\end{tcolorbox}

            \begin{tcolorbox}[breakable, size=fbox, boxrule=.5pt, pad at break*=1mm, opacityfill=0]
\prompt{Out}{outcolor}{107}{\boxspacing}
\begin{Verbatim}[commandchars=\\\{\}]
CategoricalDtype(categories=[3, 4, 5, 6, 8], ordered=False,
categories\_dtype=int64)
\end{Verbatim}
\end{tcolorbox}
        
    Now that \texttt{cylinders} is qualitative, we can display it using the
\texttt{boxplot()} method.

    \begin{tcolorbox}[breakable, size=fbox, boxrule=1pt, pad at break*=1mm,colback=cellbackground, colframe=cellborder]
\prompt{In}{incolor}{108}{\boxspacing}
\begin{Verbatim}[commandchars=\\\{\}]
\PY{n}{fig}\PY{p}{,} \PY{n}{ax} \PY{o}{=} \PY{n}{subplots}\PY{p}{(}\PY{n}{figsize}\PY{o}{=}\PY{p}{(}\PY{l+m+mi}{8}\PY{p}{,} \PY{l+m+mi}{8}\PY{p}{)}\PY{p}{)}
\PY{n}{Auto}\PY{o}{.}\PY{n}{boxplot}\PY{p}{(}\PY{l+s+s1}{\PYZsq{}}\PY{l+s+s1}{mpg}\PY{l+s+s1}{\PYZsq{}}\PY{p}{,} \PY{n}{by}\PY{o}{=}\PY{l+s+s1}{\PYZsq{}}\PY{l+s+s1}{cylinders}\PY{l+s+s1}{\PYZsq{}}\PY{p}{,} \PY{n}{ax}\PY{o}{=}\PY{n}{ax}\PY{p}{)}\PY{p}{;}
\end{Verbatim}
\end{tcolorbox}

    \begin{center}
    \adjustimage{max size={0.9\linewidth}{0.9\paperheight}}{artifacts/latex/Ch02-statlearn-lab_files/artifacts/latex/Ch02-statlearn-lab_251_0.png}
    \end{center}
    { \hspace*{\fill} \\}
    
    The \texttt{hist()} method can be used to plot a \emph{histogram}.

    \begin{tcolorbox}[breakable, size=fbox, boxrule=1pt, pad at break*=1mm,colback=cellbackground, colframe=cellborder]
\prompt{In}{incolor}{109}{\boxspacing}
\begin{Verbatim}[commandchars=\\\{\}]
\PY{n}{fig}\PY{p}{,} \PY{n}{ax} \PY{o}{=} \PY{n}{subplots}\PY{p}{(}\PY{n}{figsize}\PY{o}{=}\PY{p}{(}\PY{l+m+mi}{8}\PY{p}{,} \PY{l+m+mi}{8}\PY{p}{)}\PY{p}{)}
\PY{n}{Auto}\PY{o}{.}\PY{n}{hist}\PY{p}{(}\PY{l+s+s1}{\PYZsq{}}\PY{l+s+s1}{mpg}\PY{l+s+s1}{\PYZsq{}}\PY{p}{,} \PY{n}{ax}\PY{o}{=}\PY{n}{ax}\PY{p}{)}\PY{p}{;}
\end{Verbatim}
\end{tcolorbox}

    \begin{center}
    \adjustimage{max size={0.9\linewidth}{0.9\paperheight}}{artifacts/latex/Ch02-statlearn-lab_files/artifacts/latex/Ch02-statlearn-lab_253_0.png}
    \end{center}
    { \hspace*{\fill} \\}
    
    The color of the bars and the number of bins can be changed:

    \begin{tcolorbox}[breakable, size=fbox, boxrule=1pt, pad at break*=1mm,colback=cellbackground, colframe=cellborder]
\prompt{In}{incolor}{110}{\boxspacing}
\begin{Verbatim}[commandchars=\\\{\}]
\PY{n}{fig}\PY{p}{,} \PY{n}{ax} \PY{o}{=} \PY{n}{subplots}\PY{p}{(}\PY{n}{figsize}\PY{o}{=}\PY{p}{(}\PY{l+m+mi}{8}\PY{p}{,} \PY{l+m+mi}{8}\PY{p}{)}\PY{p}{)}
\PY{n}{Auto}\PY{o}{.}\PY{n}{hist}\PY{p}{(}\PY{l+s+s1}{\PYZsq{}}\PY{l+s+s1}{mpg}\PY{l+s+s1}{\PYZsq{}}\PY{p}{,} \PY{n}{color}\PY{o}{=}\PY{l+s+s1}{\PYZsq{}}\PY{l+s+s1}{red}\PY{l+s+s1}{\PYZsq{}}\PY{p}{,} \PY{n}{bins}\PY{o}{=}\PY{l+m+mi}{12}\PY{p}{,} \PY{n}{ax}\PY{o}{=}\PY{n}{ax}\PY{p}{)}\PY{p}{;}
\end{Verbatim}
\end{tcolorbox}

    \begin{center}
    \adjustimage{max size={0.9\linewidth}{0.9\paperheight}}{artifacts/latex/Ch02-statlearn-lab_files/artifacts/latex/Ch02-statlearn-lab_255_0.png}
    \end{center}
    { \hspace*{\fill} \\}
    
    See \texttt{Auto.hist?} for more plotting options.

We can use the \texttt{pd.plotting.scatter\_matrix()} function to create
a \emph{scatterplot matrix} to visualize all of the pairwise
relationships between the columns in a data frame.

    \begin{tcolorbox}[breakable, size=fbox, boxrule=1pt, pad at break*=1mm,colback=cellbackground, colframe=cellborder]
\prompt{In}{incolor}{111}{\boxspacing}
\begin{Verbatim}[commandchars=\\\{\}]
\PY{n}{pd}\PY{o}{.}\PY{n}{plotting}\PY{o}{.}\PY{n}{scatter\PYZus{}matrix}\PY{p}{(}\PY{n}{Auto}\PY{p}{)}\PY{p}{;}
\end{Verbatim}
\end{tcolorbox}

    \begin{center}
    \adjustimage{max size={0.9\linewidth}{0.9\paperheight}}{artifacts/latex/Ch02-statlearn-lab_files/artifacts/latex/Ch02-statlearn-lab_257_0.png}
    \end{center}
    { \hspace*{\fill} \\}
    
    We can also produce scatterplots for a subset of the variables.

    \begin{tcolorbox}[breakable, size=fbox, boxrule=1pt, pad at break*=1mm,colback=cellbackground, colframe=cellborder]
\prompt{In}{incolor}{112}{\boxspacing}
\begin{Verbatim}[commandchars=\\\{\}]
\PY{n}{pd}\PY{o}{.}\PY{n}{plotting}\PY{o}{.}\PY{n}{scatter\PYZus{}matrix}\PY{p}{(}\PY{n}{Auto}\PY{p}{[}\PY{p}{[}\PY{l+s+s1}{\PYZsq{}}\PY{l+s+s1}{mpg}\PY{l+s+s1}{\PYZsq{}}\PY{p}{,}
                                 \PY{l+s+s1}{\PYZsq{}}\PY{l+s+s1}{displacement}\PY{l+s+s1}{\PYZsq{}}\PY{p}{,}
                                 \PY{l+s+s1}{\PYZsq{}}\PY{l+s+s1}{weight}\PY{l+s+s1}{\PYZsq{}}\PY{p}{]}\PY{p}{]}\PY{p}{)}\PY{p}{;}
\end{Verbatim}
\end{tcolorbox}

    \begin{center}
    \adjustimage{max size={0.9\linewidth}{0.9\paperheight}}{artifacts/latex/Ch02-statlearn-lab_files/artifacts/latex/Ch02-statlearn-lab_259_0.png}
    \end{center}
    { \hspace*{\fill} \\}
    
    The \texttt{describe()} method produces a numerical summary of each
column in a data frame.

    \begin{tcolorbox}[breakable, size=fbox, boxrule=1pt, pad at break*=1mm,colback=cellbackground, colframe=cellborder]
\prompt{In}{incolor}{113}{\boxspacing}
\begin{Verbatim}[commandchars=\\\{\}]
\PY{n}{Auto}\PY{p}{[}\PY{p}{[}\PY{l+s+s1}{\PYZsq{}}\PY{l+s+s1}{mpg}\PY{l+s+s1}{\PYZsq{}}\PY{p}{,} \PY{l+s+s1}{\PYZsq{}}\PY{l+s+s1}{weight}\PY{l+s+s1}{\PYZsq{}}\PY{p}{]}\PY{p}{]}\PY{o}{.}\PY{n}{describe}\PY{p}{(}\PY{p}{)}
\end{Verbatim}
\end{tcolorbox}

            \begin{tcolorbox}[breakable, size=fbox, boxrule=.5pt, pad at break*=1mm, opacityfill=0]
\prompt{Out}{outcolor}{113}{\boxspacing}
\begin{Verbatim}[commandchars=\\\{\}]
              mpg       weight
count  392.000000   392.000000
mean    23.445918  2977.584184
std      7.805007   849.402560
min      9.000000  1613.000000
25\%     17.000000  2225.250000
50\%     22.750000  2803.500000
75\%     29.000000  3614.750000
max     46.600000  5140.000000
\end{Verbatim}
\end{tcolorbox}
        
    We can also produce a summary of just a single column.

    \begin{tcolorbox}[breakable, size=fbox, boxrule=1pt, pad at break*=1mm,colback=cellbackground, colframe=cellborder]
\prompt{In}{incolor}{114}{\boxspacing}
\begin{Verbatim}[commandchars=\\\{\}]
\PY{n}{Auto}\PY{p}{[}\PY{l+s+s1}{\PYZsq{}}\PY{l+s+s1}{cylinders}\PY{l+s+s1}{\PYZsq{}}\PY{p}{]}\PY{o}{.}\PY{n}{describe}\PY{p}{(}\PY{p}{)}
\PY{n}{Auto}\PY{p}{[}\PY{l+s+s1}{\PYZsq{}}\PY{l+s+s1}{mpg}\PY{l+s+s1}{\PYZsq{}}\PY{p}{]}\PY{o}{.}\PY{n}{describe}\PY{p}{(}\PY{p}{)}
\end{Verbatim}
\end{tcolorbox}

            \begin{tcolorbox}[breakable, size=fbox, boxrule=.5pt, pad at break*=1mm, opacityfill=0]
\prompt{Out}{outcolor}{114}{\boxspacing}
\begin{Verbatim}[commandchars=\\\{\}]
count    392.000000
mean      23.445918
std        7.805007
min        9.000000
25\%       17.000000
50\%       22.750000
75\%       29.000000
max       46.600000
Name: mpg, dtype: float64
\end{Verbatim}
\end{tcolorbox}
        
    To exit \texttt{Jupyter}, select \texttt{File\ /\ Shut\ Down}.


    % Add a bibliography block to the postdoc
    
    
    
\end{document}
